\documentclass[11pt,a4paper]{article}
\usepackage[margin=1in]{geometry}
\usepackage{amsmath,amssymb,amsthm,mathtools}
\usepackage{enumitem}
\usepackage{booktabs}
\usepackage[colorlinks=true,linkcolor=black,citecolor=black,urlcolor=black]{hyperref}
\usepackage[most]{tcolorbox}
\tcbuselibrary{skins,breakable}
\usepackage{fancyhdr}
\usepackage{caption}

% ── Theorem environments ──
\newtheorem{theorem}{Theorem}[section]
\newtheorem{lemma}[theorem]{Lemma}
\newtheorem{proposition}[theorem]{Proposition}
\newtheorem{corollary}[theorem]{Corollary}
\theoremstyle{definition}
\newtheorem{definition}[theorem]{Definition}
\newtheorem{remark}[theorem]{Remark}

% ── Box style (matches Paper 1) ──
\tcbset{
    apfbox/.style={
        enhanced,
        colback=black!3,
        frame hidden,
        borderline={0.5pt}{0pt}{black},
        arc=0pt,
        left=8pt, right=8pt, top=8pt, bottom=8pt,
        fonttitle=\bfseries\color{black},
        coltitle=black,
        detach title,
        before upper={\tcbtitle\par\smallskip},
        before skip=14pt,
        after skip=14pt,
        grow to left by=-8pt,
        grow to right by=-8pt,
        sharp corners,
        breakable,
        pad after break=4pt,
        pad before break=4pt,
        lines before break=3,
    }
}

% ── Shortcuts ──
\newcommand{\eps}{\varepsilon}
\newcommand{\Dgamma}{\mathcal{D}(\Gamma)}
\newcommand{\End}{\mathrm{End}}
\newcommand{\im}{\mathrm{im}}
\newcommand{\tr}{\mathrm{tr}}
\newcommand{\Aut}{\mathrm{Aut}}
\newcommand{\Gr}{\mathrm{Gr}}
\newcommand{\sa}{\mathrm{sa}}
\newcommand{\op}{\mathrm{op}}
\DeclareMathOperator{\Ann}{Ann}
\DeclareMathOperator{\ran}{ran}
\DeclareMathOperator{\spn}{span}
\DeclareMathOperator{\conv}{conv}
\DeclareMathOperator{\diag}{diag}
\DeclareMathOperator{\rad}{rad}

\pagestyle{fancy}
\fancyhf{}
\fancyhead[L]{\small APF Formal Supplement}
\fancyhead[R]{\small E.\,S.\ Brooke}
\fancyfoot[C]{\thepage}

% ═══════════════════════════════════════════════════════════════
\begin{document}

\begin{center}
{\LARGE\bfseries Formal Supplement: End-to-End Derivation\\[4pt]
from Finite Enforcement Capacity to Quantum Structure}

\bigskip
{\large E.\,S.\ Brooke}\\[4pt]
{\small Companion to ``The Enforceability of Distinction''}

\bigskip
{\small\today}
\end{center}

\bigskip

\begin{abstract}
This supplement presents a self-contained reconstruction of a finite-dimensional noncommutative operator skeleton from the Admissibility Physics Framework's axiom set (A1: finite enforcement capacity; MD: uniform per-test cost floor; BW: budget-window richness) together with constitutive admissibility semantics (SP, K3, FD1--FD6).  Every auxiliary structure (vector space, inner product, adjoint, algebra, Hilbert space) is either derived or explicitly constructed from these inputs, with its logical status flagged at the point of use.

The derivation has three tiers.  \textit{Core chain} (A1 + MD + BW + constitutive semantics): finite-dimensional substrate, sector decomposition, sector-orthogonal bilinear form, self-adjoint enforcement projections, superadditivity, noncommutative $*$-algebra, Wedderburn--Artin decomposition, GNS Hilbert space, faithful positive state.  \textit{Field selection} (additional closure assumptions): real-block exclusion from state-separation completeness (single-interface, T2c.1); quaternionic-block exclusion from composition closure (multi-interface, T2c.2); complex field as conclusion.  \textit{Standard quantum corollaries} (invoking external theorems): Born rule via Busch's generalized Gleason theorem; Tsirelson bound via the Cirel'son identity; no-broadcasting and monogamy from the derived algebra.

The document is standalone; no result from the main paper is assumed without re-statement.  Regularity conditions are flagged at the point of entry.  Conditional and deferred claims are given precise status throughout.
\end{abstract}

\medskip
\begin{tcolorbox}[apfbox, title={What Is Assumed, What Is Proved, What Is Deferred}]
\textbf{Inputs (three, and only three).}
\begin{enumerate}[nosep,label=\arabic*.,leftmargin=2em]
\item \textbf{A1} (axiom): enforcement capacity is finite and positive.
\item \textbf{MD} (regularity condition): uniform positive cost per unit of enforcement complexity ($\eps(d) \ge \mu^* \cdot n(d)$ with $\mu^* > 0$).  Not derivable from A1; explicit countermodel given.  The strong reading (cost isotropy, $L_{\mathrm{iso}}$) is a strengthening used only by $L_{\mathrm{cost}}$; the core chain requires only the floor.
\item \textbf{BW} (richness premise): the cost spectrum witnesses a budget-window triple.  Generic at any non-trivial interface ($n_{\max} \ge 3$); enters only at T1.
\end{enumerate}

\smallskip\noindent\textbf{Constitutive principles (definitional, not empirical).}
\begin{enumerate}[nosep,label=\arabic*.,leftmargin=2em]
\item \textbf{SP} (Substrate Faithfulness): the physical substrate is the quotient by enforcement-indistinguishable configurations.  This defines what ``physical'' means within A1's scope.
\item \textbf{K3} (Forced Additivity): the unique scalar admissibility semantics on disjoint supports compatible with exact budget accounting, support locality, and parallel realizability.  Not a theorem of bare A1; derived from A1 plus the requirement that the budget semantics be sound, complete, and support-local.
\item \textbf{FD1--FD6}: formal definitions of interfaces, distinctions, anchors, perturbation cost, enforcement operations, and joint enforcement.  FD1's predicate set is unrestricted (every constructible binary predicate with positive enforcement cost is included); this is a definitional choice, not a physical axiom.
\end{enumerate}

\smallskip\noindent\textbf{Structural outputs proved in this supplement.}
\begin{enumerate}[nosep,label=\arabic*.,leftmargin=2em]
\item Finite-dimensional real vector space $V$ and canonical embedding (OR1, $T_{\mathrm{form}}$, $T_{\mathrm{embed}}$).
\item Sector decomposition into disjoint anchor sets plus pool ($T_{\mathrm{sep}}$).
\item Uniform cost floor $\eps \ge \eps^* > 0$ ($L_{\eps^*}$).  \textit{Conditional on $L_{\mathrm{iso}}$:} exact integrality $\eps(d) = n(d) \cdot \eps^*$ ($L_{\mathrm{cost}}$).
\item Cost bilinear form with sector orthogonality ($L_\omega$).
\item Self-adjoint enforcement projections ($T_{\mathrm{adj}}$).  \textit{Structurally invariant across the admissible $\kappa$-class}: any cost function satisfying K1--K3 produces the same projection structure (Proposition~\ref{prop:kappa_class}).
\item Superadditivity of co-located enforcement costs ($L_\Delta$).
\item Direct-sum failure and codespace rotation ($L_{\mathrm{blk}}$, $L_\Pi$).
\item Finite-dimensional associative noncommutative real $*$-algebra ($T_{\mathrm{alg}}$).
\item Wedderburn--Artin decomposition and GNS representation (T2a, $T_{\mathrm{GNS}}$).
\item Complex field selection per block: real exclusion via state-separation completeness (T2c.1), quaternionic exclusion via composition closure (T2c.2).
\end{enumerate}

\smallskip\noindent\textbf{Short corollaries in the derived algebra.}
\begin{enumerate}[nosep,label=\arabic*.,leftmargin=2em]
\item Born rule ($T_{\mathrm{Born}}$): from T2c + $L_{\mathrm{cert}}$ + $L_{\mathrm{prob}}$ + Busch's theorem.
\item Tsirelson bound ($T_{\mathrm{Tsirelson}}$): from T2a + Cirel'son identity.  Saturation conditional on $\mathbb{F} = \mathbb{C}$.
\item Interface monogamy ($T_M$): from A1 + $L_{\mathrm{loc}}$ + $L_\Delta$.
\item No-broadcasting ($T_{\mathrm{NB}}$): Route~1 unconditional; Route~2 conditional on $T_{\mathrm{tensor}}$.
\end{enumerate}

\smallskip\noindent\textbf{Not proved from A1 alone (honest flags).}
\begin{enumerate}[nosep,label=\arabic*.,leftmargin=2em]
\item MD and BW are not derived; their epistemic status is stated at point of entry.
\item Cost isotropy ($L_{\mathrm{iso}}$: all irreducible distinctions have equal cost) follows from the intended reading of MD's ``regardless'' clause but could alternatively be viewed as a mild strengthening.  $L_{\mathrm{cost}}$'s integrality ($\eps(d) = n(d) \cdot \eps^*$) is conditional on $L_{\mathrm{iso}}$.  The core chain ($L_\Delta$ through $T_{\mathrm{Born}}$) does not depend on either.
\item The tomographic confirmation route ($P_{\mathrm{tom}}$) is not fully derived as a standalone theorem; it is retained as an independent corroborative dimension-count argument (T2c.1, corroboration paragraph).
\item Tensor-product identification ($T_{\mathrm{tensor}}$) is conditional on T2c; the algebraic route to no-broadcasting (Route~2) inherits this condition.
\item CPTP dynamics ($T_{\mathrm{CPTP}}$) is not proved here; it requires a dynamical framework beyond the scope of this document.
\end{enumerate}
\end{tcolorbox}

\smallskip\noindent\textit{Non-A1 ingredients: explicit inventory.}  Beyond A1, this supplement uses three kinds of non-A1 ingredients, none of which is claimed to be derivable from bare A1:  (i)~\emph{regularity conditions} (MD: uniform per-test cost floor; BW: budget-window richness) --- these exclude pathological infinite-dimensional or degenerate cases;  (ii)~\emph{constitutive admissibility semantics} (SP: substrate faithfulness; K3: forced additivity on disjoint supports; FD1--FD6: formal definitions) --- these define what ``admissible enforcement'' means within A1's scope, analogous to the axioms that define a metric space given the concept of distance;  (iii)~\emph{additional closure assumptions for field selection and composition} (state-separation completeness: Definition~\ref{def:SSC}; composition closure: T2c.2) --- these enter only at the field-selection tier and are flagged at point of use.

\medskip
\begin{center}
\small
\begin{tabular}{@{}llll@{}}
\toprule
\textbf{Stage} & \textbf{Key result} & \textbf{Uses} & \textbf{What it establishes} \\
\midrule
Arena & $T_{\mathrm{form}}$ & A1 + MD + SP & Finite-dimensional substrate \\
Arena & $T_{\mathrm{sep}}$ & A1 + FD3 & Sector decomposition \\
Arena & $L_{\mathrm{cost}}$ & A1 + MD + $L_{\mathrm{iso}}$ & Unique cost functional (cond.) \\
Bridge & $L_\Delta$ & A1 + K2 + K3 + $T_{\mathrm{sep}}$ & Superadditive enforcement costs \\
Bridge & T1 & $L_\Delta$ + BW & Order-dependent enforcement \\
Bridge & $T_{\mathrm{alg}}$ & $L_{\mathrm{blk}}$ + $T_{\mathrm{adj}}$ & Noncommutative $*$-algebra \\
Skeleton & T2a & $T_{\mathrm{alg}}$ + O3 + O4\textsuperscript{c} & Wedderburn--Artin decomposition \\
Skeleton & $T_{\mathrm{GNS}}$ & T2a + O4\textsuperscript{c} & Hilbert-space representation \\
Skeleton & T2c.1 & T2a + state-sep.\ (Lem.~\ref{lem:state_sep}) & $\mathbb{R}$-exclusion (proved) \\
Skeleton & T2c.2 & T2a + comp.\ closure & $\mathbb{H}$-exclusion (conditional) \\
\bottomrule
\end{tabular}
\end{center}
\smallskip\noindent\textit{Regularity entry points:}  MD enters at $L_{\eps^*}$ (cost floor); BW enters only at T1 (order witness).  Constitutive principles (SP, K3, FD1--FD6) enter at the definitions and are used throughout.  \textsuperscript{c}O4 is \emph{constructed} from $\tr$ on $\End(V)$ and the enforcement costs $\eps(d)$ (Proposition~\ref{prop:O4}); it is not an additional input.  Complete dependency table: Section~\ref{sec:dependency}.

\smallskip\noindent\textit{Scope of this document.}  Every theorem in this supplement is proved from the three stated inputs (A1, MD, BW), the constitutive principles (SP, K3), and the formal definitions (FD1--FD6), using only the constructions given herein.  No result depends on any companion paper.  References to ``Paper~2,'' ``Paper~3,'' ``Paper~6,'' etc.\ appear in this document solely as forward pointers: they indicate where the derivation chain continues beyond this supplement's scope, not results that have been peer-validated or that this supplement's conclusions depend on.  Claims about gauge content, field content, cosmological parameters, and spectral triples attributed to companion papers should be understood as the intended program of the series, not as established results.  The structural claims that \emph{are} established here --- Hilbert space, $C^*$-algebra, Born rule, Tsirelson bound, no-broadcasting --- stand or fall on the proofs in this document alone.

\smallskip\noindent\textit{Structure.}  Part~I (Sections 1--6) contains the complete derivation chain: inputs, definitions, lemmas, theorems, worked examples, and dependency map.  Part~II (Sections 7--9) collects extended commentary, comparison with other reconstruction programs, and scope analysis.  A verifying reader need only read Part~I.

\medskip\noindent\textbf{Symbol legend.}

\smallskip
\begin{center}\small
\renewcommand{\arraystretch}{1.15}
\begin{tabular}{@{}lllp{4.5cm}@{}}
\toprule
\textbf{APF symbol} & \textbf{Standard equiv.} & \textbf{Defined} & \textbf{Meaning} \\
\midrule
$\Gamma$ & --- & FD1 & Enforcement interface (physical boundary) \\
$S_\Gamma$ & State space & FD1 & Substrate: set of configurations at $\Gamma$ \\
$\mathcal{D}(\Gamma)$ & Effect set & FD1 & Enforceable distinctions (binary predicates, $\eps > 0$) \\
$\eps(d)$ & Weight / cost & FD4 & Enforcement cost: $\inf_{\mathcal{P}(d)} \kappa$ \\
$C(\Gamma)$ & Capacity & A1 & Total enforcement budget at $\Gamma$ \\
$\eps^*$ & Min.\ weight & $L_{\eps^*}$ & Uniform cost floor: $\eps(d) \ge \eps^*$ for all $d$ \\
$\mu^*$ & Cost density & MD & Cost per unit complexity: $\eps(d) \ge \mu^* n(d)$ \\
$M_d$ & Face / range & FD3 & Anchor set: minimal perturbation-support subspace \\
$\Pi$ & Pool & $T_{\mathrm{sep}}$ & Shared substrate: $S_\Gamma \ominus \bigoplus_d M_d$ \\
$\kappa(p)$ & Cost fn.\ & FD4 & Perturbation cost function on $\End(S_\Gamma)$ \\
$E_d$ & Compression / proj. & FD5 & Enforcement operation: $B$-orth.\ proj.\ onto $M_d$ \\
$B(\cdot,\cdot)$ & Inner product & $L_\omega$ & Sector-orthogonal bilinear form on $V$ \\
$\mathcal{A}$ & Obs.\ algebra & O3 & $\mathrm{Alg}_\mathbb{R}(\{E_d, E_{\{d_i,d_j\}}\}) \subseteq \End(V)$ \\
$\omega$ & Functional (constructed) & O4 & Faithful positive state on $\mathcal{A}$; built from $\tr$ + $\eps(d)$ + $E_d$ \\
$\mathcal{H}_\omega$ & Hilbert space & $T_{\mathrm{GNS}}$ & GNS representation space from $(\mathcal{A}, \omega)$ \\
$\Delta(d_1,d_2)$ & Superadd.\ surplus & $L_\Delta$ & $\eps(\{d_1,d_2\}) - \eps(d_1) - \eps(d_2) > 0$ \\
\bottomrule
\end{tabular}
\end{center}

\medskip\noindent\textbf{APF $\leftrightarrow$ standard terminology crosswalk.}

\smallskip
\begin{center}\small
\renewcommand{\arraystretch}{1.15}
\begin{tabular}{@{}ll@{}}
\toprule
\textbf{APF term} & \textbf{Standard / GPT equivalent} \\
\midrule
Distinction $d$ & Effect / binary test / observable outcome \\
Enforcement projection $E_d$ & Sharp effect / compression / filter \\
Anchor set $M_d$ & Support / range sector of $E_d$ \\
Pool $\Pi$ & Shared latent support / coupling sector \\
Enforcement cost $\eps(d)$ & Weight / resource cost of maintaining $d$ \\
Budget $C(\Gamma)$ & Total capacity / resource bound at interface \\
Admissible state & Budget-compatible operational state \\
Superadditivity surplus $\Delta$ & Nonclassicality witness / coupling measure \\
Sector decomposition ($T_{\mathrm{sep}}$) & Face decomposition of state space \\
\bottomrule
\end{tabular}
\end{center}

\tableofcontents
\newpage

% ╔═══════════════════════════════════════════════════════════════╗
% ║                    PART I: CORE DERIVATION                    ║
% ║  Definitions → Lemmas → Theorems → QED.                      ║
% ╚═══════════════════════════════════════════════════════════════╝

\part{Core Derivation}
\label{part:core}

% ═══════════════════════════════════════════════════════════════
\section{Inputs: One Axiom, One Regularity Condition, One Richness Premise}
\label{sec:inputs}

The derivation uses exactly three inputs.  Every subsequent structure is either defined (carrying no empirical content) or derived as a theorem.

\begin{tcolorbox}[apfbox, title={A1: Enforcement Capacity is Finite and Positive (Axiom)}]
For every physical region $R$, there exists a finite positive real number $C(R) \in \mathbb{R}_{>0}$ such that the total enforcement cost of all simultaneously maintained distinctions in $R$ does not exceed $C(R)$:
\[
\sum_{d \in S_\sigma} \eps(d) \;\le\; C(R) \;<\; \infty,
\]
where $S_\sigma \subseteq \mathcal{D}(\sigma, R)$ is the set of actively maintained distinctions in admissible state $\sigma$, and $\eps(d) > 0$ for every $d \in \mathcal{D}$.  The presupposition $|\mathcal{D}| \ge 2$ (at least two non-identical distinctions exist) is folded into A1 itself.
\end{tcolorbox}

\begin{tcolorbox}[apfbox, title={MD: Uniform Per-Test Cost Floor (Regularity Condition)}]
There exists $\mu^* > 0$ such that for every $d \in \mathcal{D}$:
\[
\eps(d) \;\ge\; \mu^* \cdot n(d),
\]
where $n(d)$ is the enforcement complexity of $d$ (the minimum number of independent binary tests required to determine the enforcement status of $d$; after the vector-space construction, $n(d) = \dim(M_d)$).  MD asserts that each unit of enforcement complexity costs at least $\mu^*$, regardless of which distinction is being enforced.  \textit{Note:} MD is stated in operational terms ($n(d)$ as binary-test count).  Prior to $T_{\mathrm{embed}}$, $n(d)$ carries this operational meaning; after $T_{\mathrm{embed}}$ establishes the vector-space structure, the operational $n(d)$ coincides with $\dim(M_d)$, and all subsequent uses of MD invoke this mathematical identification.  \textit{Why the identification is exact:}  $T_{\mathrm{embed}}$'s injection $\varphi$ maps each distinction to a cost-coordinate direction in $V$; two operationally independent tests (distinct perturbation classes $\mathcal{P}(d_1) \ne \mathcal{P}(d_2)$) necessarily map to linearly independent anchor directions, because if they mapped to the same direction, the $\mathcal{D}$-quotient would identify them as the same distinction --- contradicting operational independence.  Therefore operational independence of tests is isomorphic to linear independence in $V$, and $n(d) = \dim(M_d)$ exactly.

\smallskip
\noindent\textit{Status.}  Regularity condition, not derivable from A1 alone.  Explicit countermodel: $\eps(d_n) = 1/2^n$ with $n(d_n) = 1$ for all $n$ satisfies A1 but has $\eps^* = 0$.  Analogous to Hardy's continuous reversibility or CDP's purification.  Independently motivated by the Bekenstein--Hawking entropy bound, which suggests $\mu^* \ge 4\ell_P^2/A$ per binary test; that dimensional identification is developed in a companion paper and is not required here.

\smallskip
\noindent\textit{Operational meaning of $n(d)$.}  The enforcement complexity $n(d)$ is the minimum number of independent binary yes/no tests required to determine $d$'s enforcement status.  For a spin-$\tfrac{1}{2}$ system, the distinction ``spin is up along $\hat{z}$'' requires one binary test ($n = 1$, $\eps = \eps^*$).  For a two-qubit system, the distinction ``both spins aligned'' requires two independent tests ($n = 2$, $\eps = 2\eps^*$).  MD asserts that this complexity--cost scaling is uniform: no distinction receives a per-test discount.  The Bekenstein--Hawking bound suggests a physical value of $\mu^*$: the maximum number of independent binary tests at an interface of area $A$ is $n_{\max} \le A / 4\ell_P^2$, giving $\mu^* \ge 4\ell_P^2 / A$ per test.  That dimensional connection is developed in a companion paper; MD enters the present work as a regularity condition whose physical content is independently motivated but not derived here.

\smallskip
\noindent\textit{Necessity: what breaks without MD.}  MD is not a convenience; it is load-bearing at four specific points in the derivation chain.  Below, each failure is exhibited via the countermodel $\eps(d_n) = 1/2^n$ with $n(d_n) = 1$ for all $n$, which satisfies A1 (with $C = 2$, since $\sum_{n=1}^\infty 1/2^n = 1 < 2$) but has $\eps^* = \inf_n \eps(d_n) = 0$.

\noindent(1)~\textit{$L_{\eps^*}$ fails $\Rightarrow$ no uniform cost floor.}  Without $\eps^* > 0$, distinctions can be arbitrarily cheap.  The cost spectrum has no gap above zero, so arbitrarily many distinctions can coexist at finite total cost.

\noindent(2)~\textit{$T_{\mathrm{form}}$ fails $\Rightarrow$ infinite-dimensional substrate.}  The bound $\dim(S_\Gamma) \le C/\mu^*$ diverges as $\mu^* \to 0$.  In the countermodel, infinitely many independent distinctions (each with $n(d_n) = 1$) fit within $C = 2$, so $\dim(S_\Gamma) = \infty$.  This is analogous to a quantum harmonic oscillator: finite energy, infinite-dimensional Hilbert space.  MD excludes this by imposing $\mu^* > 0$, ensuring finite capacity supports only finitely many independent degrees of freedom.

\noindent(3)~\textit{$L_{\mathrm{cost}}$ fails $\Rightarrow$ no unique cost functional.}  Without MD, the step H1 (ledger completeness) in $L_{\mathrm{cost}}$'s proof does not collapse cost to a function of channel count alone.  The costs $\eps(d_n) = 1/2^n$ depend on the specific distinction, not just on $n(d) = 1$, so $f(1)$ is not well-defined.

\noindent(4)~\textit{Wedderburn--Artin fails $\Rightarrow$ no matrix-algebra decomposition.}  T2a requires $\dim(\mathcal{A}) < \infty$, which follows from OR3 ($\dim(V) < \infty$, from $T_{\mathrm{form}}$).  Without $T_{\mathrm{form}}$, the algebra $\mathcal{A} \subseteq \End(V)$ is infinite-dimensional, and the Wedderburn--Artin theorem (which applies to finite-dimensional semisimple algebras) is inapplicable.

\noindent\textit{Why MD is not too restrictive.}  MD excludes \emph{only} models with $\eps^* = 0$: those where enforcement costs can be made arbitrarily small.  It does not require $\eps(d) = \mu^* \cdot n(d)$ (equality); it requires only the floor $\ge \mu^* \cdot n(d)$.  Distinctions can cost more than the floor.  MD does not exclude fine structure across changing contexts, nonuniform costs above the floor, or renormalization across energy scales; it constrains only the infimum cost within a single fixed-capacity interface.  MD is strictly weaker than Hardy's ``continuous reversibility'' (which imposes a global symmetry group) and CDP's ``purification'' (which imposes a structural principle about extensions).  Scale-free or fractal-like cost spectra with $\eps(d) \to 0$ along some sequence are excluded, but such spectra require infinite hierarchies of increasingly cheap distinctions --- a feature incompatible with the finite-capacity philosophy of A1.  In this sense, MD is the minimal regularity condition that makes A1's finiteness constraint operationally effective.

\smallskip\noindent\textit{Weaker alternative.}  A logically weaker condition $\mathrm{MD}_{\mathrm{weak}}$ would assert only that no infinite family of independent non-null distinctions has finite total cost: $\sum_{i=1}^\infty \eps(d_i) = \infty$ for every infinite independent sequence.  $\mathrm{MD}_{\mathrm{weak}}$ suffices for finite dimensionality ($\dim(S_\Gamma) < \infty$) without providing an explicit numerical bound.  The strong version (MD as stated, with uniform floor $\mu^* > 0$) gives the quantitative bound $\dim(S_\Gamma) \le C/\mu^*$ and is used throughout this supplement for concreteness.

\smallskip\noindent\textit{Precise survival under $\mathrm{MD}_{\mathrm{weak}}$.}  The following results survive with $\mathrm{MD}_{\mathrm{weak}}$ replacing MD (qualitative finite-dimensionality without a numerical bound):  $T_{\mathrm{sep}}$ (sector decomposition: uses K3, not MD directly); $L_\omega$ (sector-orthogonal bilinear form: uses K3); $T_{\mathrm{adj}}$ (self-adjointness: uses $L_\omega$); $L_\Delta$ (superadditivity: uses K3 + SP); $L_{\mathrm{blk}}$, $L_\Pi$ (codespace rotation: uses compactness of the Grassmannian, which requires $\dim < \infty$ but not a specific bound); $T_{\mathrm{alg}}$ (noncommutativity); O3, O4 (algebra and faithful state); T2a (Wedderburn--Artin: requires $\dim(\mathcal{A}) < \infty$); $T_{\mathrm{GNS}}$ (Hilbert space); T2c.1--T2c.2, $T_{\mathrm{Born}}$, $T_{\mathrm{Tsirelson}}$, $T_M$, $T_{\mathrm{NB}}$.  \textit{What is lost:} the explicit bounds $\dim(S_\Gamma) \le C/\mu^*$, $n_{\max} = \lfloor C/\eps^*\rfloor$, and $n_{\mathrm{bc}} \le \lfloor C/\Delta_0\rfloor$ (the finite broadcasting number).  Also lost: $L_{\mathrm{cost}}$'s integrality $\eps(d) = n(d) \cdot \eps^*$ (which additionally requires $L_{\mathrm{iso}}$) and the NT corollary's explicit two-value witness.  \textit{Summary:} the qualitative algebraic skeleton (Hilbert space, $C^*$-algebra, Born rule, Tsirelson bound) is fully robust to the weakening; the quantitative capacity bounds are not.

\smallskip\noindent\textit{MD and the Markov quiet regime --- a sharper characterization.}
A companion paper (\emph{Markov Breakdown and the Hard Problems}) provides a
stronger characterization of MD's role by connecting it to the enforcement Markov
theory.  Three routes to deriving MD from A1 were systematically assessed:
\begin{enumerate}[nosep,leftmargin=2em]
  \item \textit{Budget-exhaustion via QEC.}  If $\eps^* = 0$, QEC (forced by FD6)
    must protect against arbitrarily cheap erasures.  But protection costs are
    sequential, not additive; they do not sum to a contradiction with $C < \infty$.
    Route fails.
  \item \textit{$L_{\mathrm{NZ}}$ (no-Zeno).}  $L_{\mathrm{NZ}}$ governs
    temporal enforcement \emph{histories}, not static budget allocations.  The
    countermodel $\eps(d_n) = 1/2^n$ involves simultaneous maintenance costs, not
    sequential acts.  Route fails.
  \item \textit{QEC code dimensionality via FD6.}  FD6 is proved after OR3,
    which requires MD.  Using FD6 to derive MD is circular.  Route fails.
\end{enumerate}
The countermodel therefore stands and the four failure points above remain the
correct account of MD's load-bearing role.  What the Markov theory establishes
instead is an equivalence that reframes MD's physical content without changing its
logical status.  Under A1 and $L_{\mathrm{loc}}$:
\[
  \mathrm{MD}
  \;\Longleftrightarrow\;
  \exists\,\sigma \text{ on } ABC :
  \Delta_{\mathrm{SSA}}(\sigma) = 0
  \quad\bigl[\text{``Phase~I is non-empty''}\bigr],
\]
where $\Delta_{\mathrm{SSA}} = \varepsilon_{AC} + \varepsilon_{ABC}$ is the
enforcement Markov gap and $\Delta_{\mathrm{SSA}} = 0$ is the condition that
interface $B$ screens $A$ from $C$ exactly with zero cross-gap budget.  Without MD
($\eps^* = 0$), every state has $\Delta_{\mathrm{SSA}} > 0$: cross-gap
enforcement at zero marginal cost appears universally, the Markov condition never
holds, and physics has no quiet regime.  The equivalence gives MD a new physical
reading alongside the substrate formulation already given here: MD is the minimal
condition for the enforcement framework to admit states in which the Markov
condition holds exactly.  Both readings are equivalent under A1 $+$
$L_{\mathrm{loc}}$; the structural formulation is the one the Markov theory makes
natural.  \emph{MD is the axiom that says: the world has simple cases.}
\end{tcolorbox}

\begin{tcolorbox}[apfbox, title={BW: Budget-Window Richness (Richness Premise)}]
The cost spectrum contains a triple $(d_1, d_2, d_3)$ at a single interface $\Gamma$ such that $\eps(d_1) < \eps(d_2)$ and $C - \eps(d_1) \ge \eps(d_3) > C - \eps(d_2)$.

\smallskip
\noindent\textit{Status.}  Richness premise.  Holds for any cost spectrum with $\ge 3$ distinct values at a single interface (see Genericity below).  Enters only at T1 (order-dependence witness).

\smallskip
\noindent\textit{Genericity.}  BW holds generically under mild non-degeneracy.  The witness construction uses only the MD floor ($\eps(d) \ge \mu^* \cdot n(d)$ with $\mu^* > 0$), \emph{not} the $L_{\mathrm{cost}}$ integrality $\eps(d) = n(d) \cdot \eps^*$ (which requires $L_{\mathrm{iso}}$).

At any interface with $\lfloor C/\eps^*\rfloor \ge 3$ (i.e., $C \ge 3\eps^*$, where $\eps^* := \inf_d \eps(d) > 0$ by $L_{\eps^*}$): let $d_1$ be an irreducible distinction ($n(d_1) = 1$) achieving the infimum $\eps(d_1) = \eps^*$, and let $d_2$ be a distinction with $n(d_2) = 2$.  By MD's floor: $\eps(d_2) \ge 2\mu^* \ge 2\eps^*$, so $\eps(d_1) < \eps(d_2)$.  The budget window $W := (C - \eps(d_2),\, C - \eps(d_1)]$ has width $\eps(d_2) - \eps(d_1) \ge \eps^* > 0$.  If the cost spectrum at $\Gamma$ contains at least one value in $W$, BW is witnessed.  At non-degenerate interfaces --- those whose cost spectrum is not confined to a single value --- this is a mild condition.  The non-degeneracy of the cost spectrum is a premise, not derived from A1+MD alone; it is physically generic but not automatic.

BW fails only at interfaces with $\lfloor C/\eps^*\rfloor \le 2$ (at most two independent distinctions) or with a completely degenerate cost spectrum.  Both are non-generic.

\smallskip\noindent\textit{Alternative (stronger) genericity under $L_{\mathrm{iso}}$.}  If $L_{\mathrm{iso}}$ holds (strong reading of MD), then $\eps(d) = n(d) \cdot \eps^*$ and the cost spectrum is $\{\eps^*, 2\eps^*, 3\eps^*, \ldots\}$.  In this case, the witness is explicit: $\eps(d_1) = \eps^*$, $\eps(d_2) = 2\eps^*$, $\eps(d_3) = C - \eps^* - \eta$ for small $\eta > 0$.  The core chain does not depend on this stronger form.
\end{tcolorbox}

% ═══════════════════════════════════════════════════════════════
\section{Definitions (FD1--FD6)}
\label{sec:defs}

These carry no empirical content; they are the vocabulary of the framework.

\begin{definition}[FD1: Enforcement Interface]\label{def:FD1}
An \emph{enforcement interface} $\Gamma$ is a triple $(S_\Gamma,\, \mathcal{D}(\Gamma),\, C(\Gamma))$ where $S_\Gamma$ is the \emph{enforcement substrate} (a set of possible substrate configurations), $\mathcal{D}(\Gamma)$ is the set of enforceable distinctions at $\Gamma$, and $C(\Gamma) \in \mathbb{R}_{>0}$ is the enforcement capacity exclusively assigned to $\Gamma$.

\smallskip\noindent\textit{Formal predicate class.}  A \emph{distinction} at $\Gamma$ is a function $d: \Omega(\Gamma) \to \{0, 1\}$ that partitions admissible configurations into two non-empty classes ($d^{-1}(0) \ne \emptyset$, $d^{-1}(1) \ne \emptyset$) with positive enforcement cost ($\eps(d) > 0$).  The set $\mathcal{D}(\Gamma)$ consists of \emph{all} such functions --- every binary partition of the admissible state space that costs positive capacity to maintain is included.

\smallskip\noindent\textit{What ``constructible'' means (no measurability or definability constraint needed).}  The predicate class $\mathcal{D}(\Gamma)$ is defined set-theoretically, not measure-theoretically: $d$ is any function $\Omega(\Gamma) \to \{0,1\}$ with $\eps(d) > 0$.  No $\sigma$-algebra, Borel structure, topology, or definability hierarchy is imposed on $\Omega(\Gamma)$ or on the predicates.  This is possible because the framework is fundamentally finite: $T_{\mathrm{form}}$ (proved later, from A1 + MD) bounds $|\mathcal{B}| \le C/\mu^* < \infty$, so only finitely many pairwise-independent distinctions coexist.  The state space $\Omega(\Gamma)$ is then a finite union of cells indexed by subsets of $\mathcal{B}$ with residual budget $\ge 0$, and every binary partition is trivially measurable.  After $T_{\mathrm{embed}}$ constructs the finite-dimensional vector space $V$, each $d \in \mathcal{D}(\Gamma)$ is represented by its pullback along $\varphi: \Omega(\Gamma) \hookrightarrow V$, and Borel measurability is automatic (every subset of a finite-dimensional space is Borel).

\smallskip\noindent\textit{No linear structure before $T_{\mathrm{embed}}$.}  The definition of $\mathcal{D}(\Gamma)$ uses only set-theoretic operations (binary partition, non-emptiness of classes) and the cost functional $\eps$ (a real-valued function on partitions).  No vector space, inner product, coordinate system, or ``$v$-component'' language enters the definition.  Directional language (``anchor subspace $M_d = \spn\{v\}$,'' ``support in $M_d$'') appears only after $T_{\mathrm{embed}}$ and $T_{\mathrm{sep}}$ have constructed the linear and sector structure.  This separation is verified in Remark~\ref{rem:bootstrap} (bootstrapping order).

\smallskip\noindent\textit{Definitional completeness.}  The unrestricted definition of $\mathcal{D}(\Gamma)$ (every positive-cost binary partition) is a \emph{definitional choice}, not a physical axiom.  A framework with a restricted predicate set (only certain tests ``allowed'') would have a smaller $\mathcal{D}(\Gamma)$ and would need $L_{\mathrm{pred}}$ as an additional axiom.  In the APF, $L_{\mathrm{pred}}$ follows from the unrestricted definition.

\smallskip\noindent\textit{Two-layer predicate structure.}  The full predicate universe $\mathcal{D}(\Gamma)$ contains every positive-cost binary partition, but the vector-space construction uses only a structurally controlled subclass:
\begin{itemize}[nosep,leftmargin=2em,itemsep=3pt]
\item $\mathcal{D}(\Gamma)$: the \emph{full predicate universe}.  All positive-cost binary partitions of $\Omega(\Gamma)$.  This may contain composite, redundant, or pathological predicates --- and that is fine: no result depends on excluding them.
\item $\mathcal{D}_{\mathrm{irr}}(\Gamma) \subseteq \mathcal{D}(\Gamma)$: the \emph{irreducible distinctions}.  A distinction $d$ is \emph{irreducible} if its enforcement complexity is $n(d) = 1$ (it requires exactly one binary yes/no test, with cost $\eps(d) = \eps^*$).  Equivalently, after $T_{\mathrm{embed}}$, the anchor set $M_d$ is one-dimensional ($\dim(M_d) = 1$).  The irreducible distinctions are the generators of the linearized structure.
\item $\mathcal{B} \subseteq \mathcal{D}_{\mathrm{irr}}(\Gamma)$: the \emph{maximal antichain} of pairwise independently enforceable irreducible distinctions.  $T_{\mathrm{form}}$ bounds $|\mathcal{B}| \le C/\mu^*$, and $T_{\mathrm{embed}}$ builds $V$ using $\mathcal{B}$ as its coordinate basis (Lemma~\ref{lem:affine_indep}: operational independence $\Rightarrow$ linear independence of the embedding coordinates).
\end{itemize}
The construction thus proceeds: $\mathcal{D}(\Gamma) \supseteq \mathcal{D}_{\mathrm{irr}}(\Gamma) \supseteq \mathcal{B}$.  The role of the full universe $\mathcal{D}(\Gamma)$ is to guarantee completeness: unrestrictedness ensures that no structurally necessary distinction is missing from the predicate set.  The role of $\mathcal{D}_{\mathrm{irr}}(\Gamma)$ is to identify the generators: irreducible distinctions are the atomic units of enforcement complexity, and the maximal antichain $\mathcal{B}$ among them provides the embedding coordinates.  Composite distinctions ($n(d) > 1$) are expressible as conjunctions of irreducibles and contribute no additional dimensions to $V$.  Pathological partitions (those with no clean anchor interpretation) are harmless: they exist in $\mathcal{D}(\Gamma)$ but are not selected into $\mathcal{B}$ and play no role in the construction.
\end{definition}

\begin{definition}[FD2: Admissible State]
An \emph{admissible state} at $\Gamma$ is a pair $\sigma = (S_\sigma,\, \delta_\sigma)$ where $S_\sigma \subseteq \mathcal{D}(\Gamma)$ is the set of actively maintained distinctions and $\delta_\sigma = C(\Gamma) - \sum_{d \in S_\sigma} \eps(d) \ge 0$ is the residual budget.  The set of all admissible states is $\Omega(\Gamma)$.
\end{definition}

\begin{definition}[FD3: Anchor Set]
For each $d \in \mathcal{D}(\Gamma)$, the \emph{anchor set} $M_d \subseteq S_\Gamma$ is the minimal subspace containing the support of every perturbation that can threaten $d$.  After the vector-space construction (Theorem~\ref{thm:Tembed}), $M_d$ is a linear subspace of $S_\Gamma$ with $\dim(M_d) \ge 1$.
\end{definition}

\begin{definition}[FD4: Perturbation Cost Function]\label{def:FD4}

\textit{Pre-linear formulation (operational, before $T_{\mathrm{embed}}$).}  A \emph{perturbation} of distinction $d$ is a state-transformation $p$ that can change the enforcement status of $d$.  The \emph{affected region} of $p$ is the set of substrate configurations that $p$ can modify: $\mathrm{Aff}(p) := \{\sigma \in S_\Gamma : p(\sigma) \ne \sigma\}$.  Two perturbations $p_1, p_2$ are \emph{independent} if their affected regions do not overlap: $\mathrm{Aff}(p_1) \cap \mathrm{Aff}(p_2) = \emptyset$.  The \emph{perturbation cost function} $\kappa$ assigns a non-negative real cost to each perturbation and satisfies:

\begin{itemize}[nosep,leftmargin=2em,itemsep=2pt]
\item \textbf{K1} (Non-negativity): $\kappa(p) \ge 0$ for all $p$.
\item \textbf{K2} (Positivity): If $p$ is non-trivial ($\mathrm{Aff}(p) \ne \emptyset$) and $p \in \mathcal{P}(d)$ for some $d \in \mathcal{D}$, then $\kappa(p) > 0$.
\item \textbf{K3} (Independence additivity): For independent perturbations, $\kappa(p_1 \circ p_2) = \kappa(p_1) + \kappa(p_2)$.
\end{itemize}

These three properties are stated without reference to linear structure, vector spaces, or endomorphisms.  They use only: a set $S_\Gamma$ of configurations, a cost function $\kappa: \{\text{perturbations}\} \to \mathbb{R}_{\ge 0}$, and the notion of non-overlapping affected regions.

\smallskip\noindent
\textit{Post-linear formulation (after $T_{\mathrm{embed}}$).}  Once $T_{\mathrm{embed}}$ constructs the vector space $V$ and identifies $S_\Gamma$ with $V$, perturbations become linear maps $p \in \End(V)$, and the operational notions upgrade as follows:

\smallskip
\begin{center}\small
\begin{tabular}{@{}ll@{}}
\toprule
\textbf{Pre-linear (operational)} & \textbf{Post-linear (algebraic)} \\
\midrule
Perturbation: state-transformation & Linear map $p \in \End(V)$ \\
Affected region: $\mathrm{Aff}(p) = \{\sigma : p(\sigma) \ne \sigma\}$ & Support: $\mathrm{supp}(p) := \im(p - \mathrm{id})$ \\
Independent: $\mathrm{Aff}(p_1) \cap \mathrm{Aff}(p_2) = \emptyset$ & Disjoint supports: $\mathrm{supp}(p_1) \cap \mathrm{supp}(p_2) = \{0\}$ \\
Combined: $p_1 \circ p_2$ & Direct sum: $p_1 \oplus p_2$ \\
\bottomrule
\end{tabular}
\end{center}

\smallskip\noindent
\textit{Equivalence.}  The $T_{\mathrm{embed}}$ injection $\varphi: \Omega \hookrightarrow \mathbb{R}^{N+1}$ maps each configuration to its cost-coordinate vector.  A perturbation $p$ that changes configurations $\sigma$ to $p(\sigma)$ is linearized to $\hat{p} \in \End(V)$ by $\hat{p}(\varphi(\sigma)) := \varphi(p(\sigma))$.  The operational affected region $\mathrm{Aff}(p)$ maps to the linear support: $\varphi(\mathrm{Aff}(p)) = \mathrm{supp}(\hat{p})$, because a configuration is changed iff its cost-coordinate vector is changed (the $\mathcal{D}$-quotient identifies configurations that agree on all enforcement costs).  Independence ($\mathrm{Aff}(p_1) \cap \mathrm{Aff}(p_2) = \emptyset$) maps to disjoint supports ($\mathrm{supp}(\hat{p}_1) \cap \mathrm{supp}(\hat{p}_2) = \{0\}$) because disjoint configuration sets map to linearly independent direction sets in $V$ (the same D-quotient argument as in MD's operational/algebraic bridge, Section~\ref{sec:inputs}).  The composition $p_1 \circ p_2$ maps to the direct sum $\hat{p}_1 \oplus \hat{p}_2$ because independent perturbations act on non-overlapping regions, so their joint action decomposes.

The remainder of this definition uses the post-linear formulation for mathematical precision; the operational content is the pre-linear version above, and the equivalence ensures no circularity.

\smallskip\noindent
The \emph{support} of a linear map $p \in \End(S_\Gamma)$ is the minimal linear subspace $\mathrm{supp}(p) \subseteq S_\Gamma$ such that $p$ restricts to the identity on every algebraic complement of $\mathrm{supp}(p)$; equivalently, $\mathrm{supp}(p) := \im(p - \mathrm{id})$.  For two maps $p_1, p_2$ with $\mathrm{supp}(p_1) \cap \mathrm{supp}(p_2) = \{0\}$, the direct sum $p_1 \oplus p_2$ is the unique map that agrees with $p_1$ on $\mathrm{supp}(p_1)$, with $p_2$ on $\mathrm{supp}(p_2)$, and with the identity on the complement of $\mathrm{supp}(p_1) \oplus \mathrm{supp}(p_2)$.

A \emph{perturbation} of distinction $d$ is a linear map $p \in \End(S_\Gamma)$ with $\mathrm{supp}(p) \subseteq M_d$ that can destroy $d$.  The set of all such perturbations is $\mathcal{P}(d)$.  In the post-linear formulation, K1--K3 above become:

\smallskip\noindent
\textbf{K1} (Non-negativity): $\kappa(p) \ge 0$ for all $p$.  \textit{Derivation:} consuming negative capacity would create capacity from nothing, violating A1's finite bound.

\smallskip\noindent
\textbf{K2} (Positivity for non-trivial perturbations): If $p \ne 0$ and $p \in \mathcal{P}(d)$ for some $d \in \mathcal{D}$, then $\kappa(p) > 0$.  \textit{Derivation:} if $\kappa(p) = 0$, then $p$ destroys $d$ at zero cost, but $d \in \mathcal{D}$ means $\eps(d) > 0$, contradicting $d \in \mathcal{D}$.

\smallskip\noindent
\textbf{K3} (Additivity on disjoint supports): If $\mathrm{supp}(p_1) \cap \mathrm{supp}(p_2) = \{0\}$, then $\kappa(p_1 \oplus p_2) = \kappa(p_1) + \kappa(p_2)$.

\noindent\textit{Status: forced by admissibility semantics.}  K3 is not an additional dynamical axiom.  It is the unique form that any scalar admissibility semantics can take on disjoint supports if it is to be sound, complete, and support-local (Theorem~\ref{thm:K3} below).  A1 asserts that the ``total'' enforcement cost does not exceed $C$; K3 makes ``total'' mean ``sum.''  Any other aggregation rule (e.g., ``total'' = ``maximum'') would break either soundness or completeness of the budget.  See Theorem~\ref{thm:K3} for the precise statement and proof.

\smallskip\noindent
\textbf{K4} (Continuity): $\kappa$ is continuous on $\End(S_\Gamma)$ (automatic once $\dim(S_\Gamma) < \infty$, from $T_{\mathrm{form}}$).

\smallskip\noindent
The \emph{enforcement cost} of distinction $d$ is $\eps(d) := \inf_{p \in \mathcal{P}(d)} \kappa(p) > 0$.
\end{definition}

\begin{remark}[Bootstrapping order: prelinear and postlinear $\kappa$]\label{rem:bootstrap}
FD4 and K1--K4 are stated above in their post-$T_{\mathrm{embed}}$ form (using $\End(S_\Gamma)$, linear maps, supports as subspaces) for mathematical precision.  The operational content is well-defined before the vector-space construction:
\begin{itemize}[nosep,leftmargin=2em,itemsep=2pt]
\item \textit{Prelinear perturbation:} any enforcement act that can change the status of a distinction at $\Gamma$.
\item \textit{Prelinear support:} the set of substrate degrees of freedom on which the act operates.
\item \textit{Prelinear disjointness:} two acts operate on non-overlapping DOF.
\item \textit{Prelinear K3:} acts on non-overlapping DOF have additive costs.
\end{itemize}
No result prior to $T_{\mathrm{embed}}$ (OR1, $L_{\eps^*}$, $T_{\mathrm{form}}$) invokes $\End(V)$, direct sums, or operator norms.  $T_{\mathrm{form}}$'s ``independence'' is independence of enforcement (disjoint anchors under FD3), not linear independence in a vector space.  After $T_{\mathrm{embed}}$ constructs $V$, the prelinear notions become the linear ones by the embedding's injectivity: operationally independent enforcement acts map to maps with disjoint linear supports, and prelinear K3 becomes the linear K3 on $\End(V)$.  Theorem~K3 below is stated for the general (prelinear) setting: $\mathsf{E}_\Gamma$ is the set of enforcement acts, not restricted to $\End(V)$.
\end{remark}

\begin{theorem}[K3: Forced Additivity on Disjoint Supports]\label{thm:K3}
Let $\kappa: \mathsf{E}_\Gamma \to \mathbb{R}_{\ge 0}$ be any scalar admissibility semantics for enforcement acts at interface $\Gamma$, where $\mathsf{E}_\Gamma$ denotes the set of enforcement acts and $\mathrm{supp}: \mathsf{E}_\Gamma \to \mathcal{P}(S_\Gamma)$ assigns to each act the subset of substrate degrees of freedom on which it operates.  For acts $x, y \in \mathsf{E}_\Gamma$ with $\mathrm{supp}(x) \cap \mathrm{supp}(y) = \{0\}$, write $x \oplus y$ for their disjoint parallel composition.  Suppose $\kappa$ satisfies:

\begin{enumerate}[nosep,leftmargin=2em,label=(S\arabic*)]
\item \textit{Budget semantics:} $x$ is admissible at budget $C$ $\iff$ $\kappa(x) \le C$.  (That is, $\kappa(x)$ is the scalar quantity that the admissibility rule compares to $C$.)
\item \textit{Support locality:} if $\mathrm{supp}(x) \cap \mathrm{supp}(y) = \{0\}$, then composing $x$ and $y$ introduces no new enforcement burden beyond the two supports already present --- there is no interaction term carried by a third, shared support channel.  (This is the operational content of FD2 in the disjoint-support regime.)
\item \textit{Parallel realizability:} if $\mathrm{supp}(x) \cap \mathrm{supp}(y) = \{0\}$, then $x$ and $y$ can be implemented in parallel at the same interface, with actual committed capacity equal to the sum of the two separate commitments.
\item \textit{Exactness:} $\kappa$ is both sound (never underestimates actual committed capacity) and complete (never overestimates it) for admissibility.
\end{enumerate}

Then for all $x, y$ with disjoint supports:
\[
\kappa(x \oplus y) = \kappa(x) + \kappa(y).
\]

\smallskip\noindent\textit{Proof.}
Take $x, y \in \mathsf{E}_\Gamma$ with disjoint supports.  By (S3), the disjoint parallel act $x \oplus y$ is physically realizable by running the two enforcement acts side by side.  By (S2), there is no interaction channel on disjoint supports, so the actual committed capacity of $x \oplus y$ is exactly the sum of the actual committed capacities of $x$ and $y$.

\textit{Excluding subadditivity.}  Suppose $\kappa(x \oplus y) < \kappa(x) + \kappa(y)$.  Choose $C$ with $\kappa(x \oplus y) \le C < \kappa(x) + \kappa(y)$.  By (S1), $x \oplus y$ is declared admissible at budget $C$.  But by (S3), its actual committed capacity equals $\kappa(x) + \kappa(y) > C$: the act is not physically realizable within $C$.  This contradicts soundness in (S4).

\textit{Excluding superadditivity.}  Suppose $\kappa(x \oplus y) > \kappa(x) + \kappa(y)$.  Set $C = \kappa(x) + \kappa(y)$.  By (S3), the disjoint parallel implementation exists and commits exactly $C$, so it is physically realizable.  But by (S1), $\kappa(x \oplus y) > C$ declares it inadmissible.  This contradicts completeness in (S4).

Both strict inequalities are impossible.  Therefore $\kappa(x \oplus y) = \kappa(x) + \kappa(y)$.  \qed

\smallskip\noindent\textit{Interpretation.}  K3 is not an independent axiom imposed alongside A1.  It is the \emph{unique} form that a scalar admissibility semantics can take on disjoint supports while remaining sound, complete, and support-local.  The theorem does not derive K3 from bare A1 alone; it derives K3 from the requirement that the admissibility semantics actually do the job the framework needs --- track real capacity usage faithfully.  A reader who prefers to view K3 as an explicit input may do so; the downstream derivation chain is unaffected.

\noindent\textit{Conditions used:} A1 (budget bound), FD2 (admissibility definition), K1--K2 (from A1), support locality (from FD2+FD3).
\end{theorem}

\begin{tcolorbox}[apfbox, title={Status of S3 and Models That Violate It}]
\textbf{What S3 is.}  S3 (parallel realizability) is a \emph{physical} claim: when two enforcement acts operate on non-overlapping substrate DOF, their actual capacity commitments do not interact, so the total committed capacity is exactly the sum.  K3 is the \emph{semantic} consequence: the scalar function $\kappa$ must assign the sum.  The K3 theorem bridges from the physical fact (S3) to the semantic requirement (K3).

\smallskip\noindent\textbf{What S3 is not.}  S3 is not derived from bare A1.  It is the physical content of support locality (FD2 + FD3).  A framework that does not assume FD3's anchor-set locality could have models satisfying A1 where S3 fails.

\smallskip\noindent\textbf{Concrete A1-compliant model violating S3.}  Consider two enforcement acts $x, y$ on disjoint substrate regions that share a common noise bath (a fluctuating environment coupling to both regions).  The noise bath creates a hidden interaction channel: running $x$ and $y$ in parallel partially correlates their capacity draws through the shared bath, so the actual committed capacity satisfies $\kappa(x \oplus y) = \sqrt{\kappa(x)^2 + \kappa(y)^2}$ (root-sum-of-squares, not sum).  This model satisfies A1 (finite total capacity: $\kappa(x \oplus y) \le \kappa(x) + \kappa(y) \le C$) but violates S3: the actual committed capacity is strictly less than the sum.  The budget semantics under this $\kappa$ would be subadditive and would violate completeness (S4): it would declare some actually enforceable states inadmissible.  FD3's anchor-set locality excludes such models by requiring that each distinction's enforcement operates entirely within its own substrate region, with no shared resource channels crossing the support boundary.

\smallskip\noindent\textbf{What the framework selects.}  Any $\kappa$ satisfying K1--K3 produces a faithful budget semantics.  Two countermodels showing what fails without K3:

\smallskip\noindent(a)~$\kappa_{\max}(p) := \|p - \mathrm{id}\|_{\op}$ (operator norm).  Satisfies K1--K2.  Gives $\kappa_{\max}(p_1 \oplus p_2) = \max(\kappa_{\max}(p_1), \kappa_{\max}(p_2))$: subadditive.  Under $\kappa_{\max}$, the semantics rejects states that are actually enforceable --- violating completeness (S4).

\smallskip\noindent(b)~$\kappa_{\mathrm{sq}}(p) := \|p - \mathrm{id}\|_F^2$ (squared Frobenius norm).  Satisfies K1--K3: $\|A \oplus B\|_F^2 = \|A\|_F^2 + \|B\|_F^2$.  The semantics is faithful.  This is the $\kappa$ used in the worked models.
\end{tcolorbox}

\begin{proposition}[Robustness of the Derivation Chain Under Approximate K3]\label{prop:robustness}
Suppose $\kappa$ satisfies K1, K2, K4, and an \emph{approximate} version of K3: for all disjoint-support acts $x, y$,
\[
|\kappa(x \oplus y) - \kappa(x) - \kappa(y)| \;\le\; \delta
\]
for some $\delta \ge 0$.  (Exact K3 is the case $\delta = 0$.)  Then the downstream structure degrades as follows:

\begin{enumerate}[nosep,label=(\roman*),leftmargin=2em]
\item \textit{Bilinear form:} the sector-orthogonal bilinear form $B$ acquires cross-sector terms of magnitude $O(\delta)$: $|B(u_d, v_{d'})| \le C_1 \delta \|u_d\| \|v_{d'}\|$ for $d \ne d'$, where $C_1$ depends only on $\dim(V)$.  Sectors are $O(\delta)$-approximately orthogonal.  \textit{Proof sketch:} $L_\omega$.5 derives $B(u_d, v_{d'}) = 0$ from $\|u_d + v_{d'}\|_B^2 = \|u_d\|_B^2 + \|v_{d'}\|_B^2$ (exact K3).  Under approximate K3: $|\|u_d + v_{d'}\|_B^2 - \|u_d\|_B^2 - \|v_{d'}\|_B^2| \le \delta$, so $|2B(u_d, v_{d'})| \le \delta$.
\item \textit{Enforcement projections:} the sector maps $E_d$ satisfy $\|E_d - E_d^\dagger\| = O(\delta)$ and $\|E_d^2 - E_d\| = O(\delta)$: they are approximate self-adjoint projections.  \textit{Proof sketch:} self-adjointness follows from range $\perp$ kernel ($T_{\mathrm{sector}}$(iv)).  With $O(\delta)$ cross-sector terms in $B$, the sectors are only approximately orthogonal, so range and kernel are only approximately $B$-orthogonal.
\item \textit{Algebra structure:} for $\delta < \delta_{\mathrm{crit}}$ (a threshold depending only on $\dim(V)$ and the spectral gaps of the unperturbed algebra), the Wedderburn--Artin block structure is \textbf{exactly preserved}: the number and sizes of simple blocks are unchanged.  This follows from the \emph{rigidity of semisimple algebras}: $H^2(\mathcal{A}, \mathcal{A}) = 0$ for semisimple $\mathcal{A}$ (Gerstenhaber 1964), so small deformations of the multiplication are equivalent to the undeformed algebra by an inner automorphism.  \textit{Proof sketch:} the deformed multiplication $m_\delta(A,B) = AB + \delta\,\mu(A,B)$ defines a cocycle in $H^2(\mathcal{A}, \mathcal{A})$; vanishing of $H^2$ means $\mu$ is a coboundary, hence the deformed algebra is isomorphic to the undeformed one.  The threshold $\delta_{\mathrm{crit}}$ is the radius of the ball in multiplication space within which this isomorphism exists; it is positive and depends continuously on the spectral gaps of $\mathcal{A}$.
\item \textit{Field selection:} the choice $\mathbb{R}/\mathbb{C}/\mathbb{H}$ per block is a discrete invariant and cannot change under continuous deformation.  Therefore field selection is \textbf{exactly stable} for $\delta < \delta_{\mathrm{crit}}$.
\item \textit{Born rule:} the probability assignment $p(P) = \tr(\rho\,P)$ is perturbed by $O(\delta)$: $|p_\delta(P) - p_0(P)| \le C_2\,\delta$ for all projections $P$.  \textit{Proof sketch:} the state $\omega_\delta$ is a continuous function of the algebra structure; projections $P_\delta$ are $O(\delta)$-close to exact projections $P_0$; therefore $|\omega_\delta(P_\delta) - \omega_0(P_0)| \le C_2\,\delta$ by Lipschitz continuity of the trace functional on bounded operators.
\end{enumerate}

\smallskip\noindent\textit{Physical interpretation.}  Boundary interactions, finite-size effects, or gravitational corrections may introduce cross-support couplings of size $\delta \sim o(\eps^*)$.  As long as $\delta \ll \eps^*$, the quantum skeleton (Hilbert space, complex field, Born rule) is preserved exactly in its discrete structure and $O(\delta)$-approximately in its continuous quantities.  The framework degrades gracefully, not catastrophically.

\smallskip\noindent\textit{Where exact K3 is used vs.\ where approximate K3 suffices.}  The derivation chain uses exact K3 at two points: (a)~$T_{\mathrm{sep}}$'s ``disjoint anchors $\Leftrightarrow$ additive costs'' biconditional, and (b)~$L_\Delta$'s cost decomposition $\eps(\{d_1,d_2\}) = \eps(d_1) + \eps(d_2) + \kappa_\Pi$.  Under approximate K3, (a) becomes an approximate biconditional (additive to within $\delta$) and (b) gives $\Delta = \kappa_\Pi + O(\delta)$, which is still positive for $\kappa_\Pi > \delta$.  The noncommutativity conclusion ($T_{\mathrm{alg}}$) survives whenever the pool defense cost exceeds the cross-talk defect: $\kappa_\Pi > \delta$.  This is the \emph{non-classical regime} in the approximate theory.
\end{proposition}

\begin{definition}[FD5: Enforcement Operation (two stages)]
\textbf{FD5a (post-$T_{\mathrm{embed}}$, pre-inner-product).}  The enforcement operation $E_d: V \to V$ is any idempotent linear map with $\im(E_d) = M_d$ and $E_d|_{M_d} = \mathrm{id}$.  That is: $E_d^2 = E_d$, and $E_d$ restricts to the identity on the anchor set of $d$.  At this stage, $E_d$ is constrained only by these two properties; its kernel (equivalently, its action on $S_\Gamma \setminus M_d$) is \emph{not} fixed by definition.  In particular, FD5a does not assert that $\ker(E_d)$ equals the $T_{\mathrm{sep}}$ complement $N_d := \bigoplus_{d' \ne d} M_{d'} \oplus \Pi$; that identification is a \emph{theorem} ($T_{\mathrm{adj}}$, Section~\ref{sec:bridge}), derived from cost minimization.

\smallskip\noindent
\textbf{FD5b (post-$T_{\mathrm{adj}}$).}  $T_{\mathrm{adj}}$ proves that the unique cost-minimizing idempotent with $\im(E_d) = M_d$ has $\ker(E_d) = N_d$ (zero spillover from $\Pi$ into $M_d$).  Once $L_\omega$ constructs the cost bilinear form $B$, this identification coincides with $B$-orthogonal complementation: $N_d = M_d^{\perp_B}$, so $E_d = \pi_{M_d}^{(B)}$, the $B$-orthogonal projection.  Self-adjointness $E_d = E_d^\dagger$ follows.
\end{definition}

\begin{definition}[FD6: Joint Enforcement Operation (post-$L_\Pi$)]
For co-located distinctions $d_1, d_2$ with $\Delta(d_1, d_2) > 0$, the joint enforcement operation is $E_{\{d_1,d_2\}} := \pi_W$, where $W$ is the minimum-cost joint defense codespace from $L_\Pi$.  By $L_\Pi$, $W \ne M_{d_1} \oplus M_{d_2}$: the codespace is rotated into the pool $\Pi$.
\end{definition}

\begin{tcolorbox}[apfbox, title={SP: Substrate Faithfulness (constitutive principle)}]
Two substrate configurations are \emph{physically identical} iff they agree on the enforcement status of every distinction $d \in \mathcal{D}(\Gamma)$.  The \emph{physical substrate} is the quotient:
\[
S_\Gamma \;:=\; S_\Gamma^{\mathrm{raw}} \;/\; N_\Gamma, \qquad N_\Gamma := \{v \in S_\Gamma^{\mathrm{raw}} : \text{no admissible distinction's status depends on } v\}.
\]
\textbf{This is a definition, not an axiom.}  It identifies the physical state space with the space of enforcement-distinguishable configurations.  It is the constitutive definition of what ``physical'' means within A1's scope: non-enforceable variables are outside the framework's domain.  It does not assert ``no hidden variables''; it asserts that variables invisible to enforcement are not variables the framework tracks.
\end{tcolorbox}

\begin{lemma}[SC: Substrate Completeness on the Physical Quotient]\label{lem:SC}
\textit{(Stated here for logical grouping with the definitions; the proof invokes $L_{\eps^*}$, which is proved in Section~\ref{sec:arena}.)}

\smallskip\noindent Every nonzero $v \in S_\Gamma$ is probed by at least one distinction $d$ with $\eps(d) > 0$, and every perturbation along $v$ has positive cost.

\noindent\textit{Proof.}  By SP, every nonzero $v \in S_\Gamma$ is non-null: there exist admissible configurations $\sigma_1, \sigma_2$ that are operationally distinguishable (some distinction's enforcement status differs between them) and whose distinguishability is attributable to $v$ (quotienting $v$ out would identify $\sigma_1$ with $\sigma_2$).  Therefore there exist perturbations $p$ that change enforcement status by acting on the degree of freedom represented by $v$.  By K2, every such $p$ has $\kappa(p) > 0$.  By $L_{\eps^*}$ (MD): any distinction probed by $v$ has $\eps(d) \ge \eps^* > 0$.  \qed

\noindent\textit{Regularity conditions used:} SP (constitutive), K2 (from A1), \fbox{MD} (via $L_{\eps^*}$).  \textit{No linear structure is used:} ``$v \in S_\Gamma$'' refers to an equivalence class in the SP quotient, not a vector in a constructed space.  The notions ``non-null,'' ``operationally distinguishable,'' and ``positive perturbation cost'' are all pre-$T_{\mathrm{embed}}$ (see Remark~\ref{rem:bootstrap}).

\smallskip\noindent\textit{The strong form: direction-anchored predicates.}  The quotient-level SC above says every physical direction is detectable.  The dimension bound ($T_{\mathrm{form}}$) requires the stronger claim:

\begin{lemma}[$L_{\mathrm{pred}}$: Direction-Anchored Predicates Exist]\label{lem:pred}
For every nonzero $v \in S_\Gamma$, there exists $d_v \in \mathcal{D}(\Gamma)$ with $M_{d_v} = \spn\{v\}$ and $\eps(d_v) \ge \eps^*$.

\smallskip\noindent\textit{Logical status.}  This is a post-$T_{\mathrm{embed}}$ result: $v \in S_\Gamma$ is a vector in the constructed space $V$, and $\spn\{v\}$ is a linear subspace.  The proof below constructs the predicate $d_v$ operationally (Stage~1, no linear structure) and then identifies it with the directional predicate (Stage~2, after $T_{\mathrm{embed}}$).
\end{lemma}

\noindent\textit{Proof} (two stages: operational, then representational).

\textit{Stage 1 (operational, no linear structure).}  By SP, $v$ is non-null: there exist admissible configurations $\sigma_1, \sigma_2 \in \Omega(\Gamma)$ that are operationally distinguishable (their enforcement status differs with respect to some distinction) and whose distinguishability is attributable to the degree of freedom $v$ (quotienting $v$ out would identify $\sigma_1$ with $\sigma_2$).  Define the binary predicate $d_v: \Omega(\Gamma) \to \{0,1\}$ by $d_v(\sigma) := 1$ if $\sigma$'s enforcement status agrees with $\sigma_1$ on the relevant observable, $d_v(\sigma) := 0$ otherwise.  This is a well-defined binary partition of $\Omega(\Gamma)$ with both classes non-empty ($\sigma_1 \in d_v^{-1}(1)$, $\sigma_2 \in d_v^{-1}(0)$).  By FD1 (definitional completeness: $\mathcal{D}(\Gamma)$ includes every positive-cost binary partition), $d_v \in \mathcal{D}(\Gamma)$ provided $\eps(d_v) > 0$.  By K2: any perturbation that changes $d_v$'s value is non-trivial and has positive cost, so $\eps(d_v) > 0$.

\textit{Stage 2 (representational, after $T_{\mathrm{embed}}$).}  Once $T_{\mathrm{embed}}$ constructs $V$ and identifies $S_\Gamma$ with a subspace of $V$, the operational predicate $d_v$ is represented by its pullback along the embedding: $d_v(\varphi(\sigma)) = 1$ iff the $v$-component of $\varphi(\sigma)$ matches $\varphi(\sigma_1)$.  The ``$v$-component'' language is now legitimate because $V$ is a vector space.  By FD3 (anchor minimality): the perturbations threatening $d_v$ are exactly those supported on $\spn\{v\}$ (a perturbation orthogonal to $v$ cannot change the $v$-component), so $M_{d_v} = \spn\{v\}$.  By $L_{\eps^*}$: $\eps(d_v) \ge \eps^*$.  \qed

\noindent\textit{What $L_{\mathrm{pred}}$ uses and where it comes from.}  Three ingredients:
\begin{itemize}[nosep,leftmargin=2em,itemsep=3pt]
\item \textbf{SP} (constitutive): the direction $v$ is physically real (non-null in the quotient).
\item \textbf{FD1} (definitional): the predicate set $\mathcal{D}(\Gamma)$ is \emph{unrestricted} --- every constructible binary predicate with positive enforcement cost is included.  This is a definitional choice about what $\mathcal{D}(\Gamma)$ contains, not a physical axiom about nature.  A framework with a restricted predicate set (only certain measurements ``allowed'') would need $L_{\mathrm{pred}}$ as an additional axiom; in the APF, it follows from the definition.
\item \textbf{FD3} (definitional): anchor sets are minimal perturbation-support subspaces.
\end{itemize}
$L_{\mathrm{pred}}$ is the formal content of the old ``strong SC.''  It is not a global closure axiom; it is a theorem from three named constitutive/definitional principles.  The derivation chain is: SP gives existence of distinct configurations $\to$ FD1 gives the distinguishing predicate $\to$ K2 gives positive cost $\to$ FD3 gives one-dimensional anchor $\to$ $L_{\eps^*}$ gives the floor.

\smallskip\noindent\textit{Where $L_{\mathrm{pred}}$ is used.}
\begin{itemize}[nosep,leftmargin=2em,itemsep=3pt]
\item \textbf{$T_{\mathrm{form}}$ (dimension bound):} constructs $k$ direction-anchored predicates with disjoint supports; K3 gives additive costs; A1 bounds the sum.  Uses $L_{\mathrm{pred}}$.
\item \textbf{$T_{\mathrm{adj}}$ (self-adjointness):} uses SP + K2 to show spillover commits positive capacity.  Does \emph{not} use $L_{\mathrm{pred}}$.
\item \textbf{$L_\Delta$ (superadditivity):} uses SP + K2 + K3 for pool defense costs.  Does \emph{not} use $L_{\mathrm{pred}}$.
\end{itemize}

\smallskip\noindent\textit{Critical-path independence of $L_\Delta$.}  $L_\Delta$'s pre-algebraic proof (Section~\ref{sec:bridge}) derives $\Delta > 0$ from K2, K3, $T_{\mathrm{sep}}$, and SP (non-null directions exist in $\Pi$).  The strictness comes from K2 applied to $\Pi$-supported perturbations.  The entire bridge chain from $L_\Delta$ through $T_{\mathrm{alg}}$ to the quantum skeleton does not use $L_{\mathrm{pred}}$ and is robust to any restriction of the predicate set.

\smallskip\noindent\textit{What fails if $L_{\mathrm{pred}}$ is absent.}  If the predicate set $\mathcal{D}(\Gamma)$ does not contain direction-anchored predicates for every non-null direction, then $T_{\mathrm{form}}$'s dimension bound ($\dim(V) \le C/\eps^*$) cannot be established by directional counting: the substrate might contain directions invisible to the available predicates, and the vector-space dimension could exceed the counted bound.  However, once finite-dimensionality is supplied by any route (e.g., $\mathrm{MD}_{\mathrm{weak}}$ giving $\dim(S_\Gamma) < \infty$ without an explicit bound, or an independent physical argument), the later chain ($T_{\mathrm{sep}}$, $L_\omega$, $T_{\mathrm{adj}}$, $L_\Delta$, $T_{\mathrm{alg}}$, T2a, $T_{\mathrm{GNS}}$, T2c, $T_{\mathrm{Born}}$) proceeds unchanged.  The vulnerability of $L_{\mathrm{pred}}$ is therefore confined to the \emph{quantitative} dimension bound; the \emph{qualitative} algebraic skeleton is independent of it.  A framework with a restricted predicate set would need $L_{\mathrm{pred}}$ (or finite-dimensionality) as an additional explicit input.
\end{lemma}

% ═══════════════════════════════════════════════════════════════
\section{Arena Chain: From A1 to a Finite-Dimensional State Space}
\label{sec:arena}

\subsection{Convex structure and linear extension (OR1)}

\begin{theorem}[OR1: Convex Mixing]\label{thm:OR1}
A1's real-valued budget $C \in \mathbb{R}_{>0}$ is divisible: for any two admissible states $\sigma, \sigma'$ and $\lambda \in [0,1]$, the mixture $\lambda\sigma + (1-\lambda)\sigma'$ is admissible.

\noindent\textit{Proof.}  $\sum_d \eps(d)[\lambda \mathbf{1}_{d \in S_\sigma} + (1-\lambda)\mathbf{1}_{d \in S_{\sigma'}}] \le \lambda C + (1-\lambda)C = C$.  \qed

\noindent\textit{Regularity conditions used:} None (A1 only).
\end{theorem}

\subsection{Minimum cost floor ($L_{\eps^*}$)}

\begin{theorem}[$L_{\eps^*}$: Uniform Cost Floor]\label{thm:Leps}
There exists $\eps^* > 0$ such that $\eps(d) \ge \eps^*$ for all $d \in \mathcal{D}$.

\noindent\textit{Proof.}  By MD: $\eps(d) \ge \mu^* \cdot n(d) \ge \mu^* \cdot 1 = \mu^*$ for every $d$ (since $n(d) \ge 1$).  Set $\eps^* := \mu^*$.  \qed

\noindent\textit{Regularity conditions used:} \fbox{MD enters here.}  A1 alone gives pointwise $\eps(d) > 0$ but not a uniform floor (countermodel: $\eps(d_n) = 1/n$).
\end{theorem}

\begin{lemma}[$L_{\mathrm{iso}}$: Cost Isotropy for Irreducible Distinctions]\label{lem:iso}
All irreducible distinctions at a given interface have equal enforcement cost: if $n(d_1) = n(d_2) = 1$, then $\eps(d_1) = \eps(d_2) = \eps^*$.

\smallskip\noindent\textit{Status: conditional on the isotropy reading of MD.}  MD states that each unit of enforcement complexity costs at least $\mu^*$, ``regardless of which distinction is being enforced.''  This ``regardless'' clause can be read in two strengths:

\begin{itemize}[nosep,leftmargin=2em,itemsep=3pt]
\item \textit{Weak reading (floor):} $\eps(d) \ge \mu^* \cdot n(d)$ for all $d$.  This is the formal content of MD as stated.  Under the weak reading, two irreducible distinctions $d_1, d_2$ with $n(d_1) = n(d_2) = 1$ satisfy $\eps(d_1) \ge \mu^*$ and $\eps(d_2) \ge \mu^*$, but $\eps(d_1) \ne \eps(d_2)$ is not excluded.  The cost function $\kappa$ could assign different infima to perturbations along different directions.
\item \textit{Strong reading (isotropy):} $\eps(d)$ depends only on $n(d)$, not on the specific direction of $M_d$.  Under this reading, $\eps(d_1) = \eps(d_2) = \eps^*$ for all irreducible $d_1, d_2$ at $\Gamma$, and $L_{\mathrm{iso}}$ holds.
\end{itemize}

The strong reading is the intended content of ``regardless'': the enforcement capacity treats all substrate directions equivalently at the per-test level.  An anisotropic cost function (one that assigns different infima to different one-dimensional subspaces) would reflect an intrinsic preferred-direction structure in the substrate --- a feature that has no source in A1's direction-agnostic budget constraint.  The present lemma therefore follows from the intended interpretation of MD.  A reader who prefers the weak reading should take $L_{\mathrm{iso}}$ as a mild strengthening of MD (``cost isotropy''), adding it as a fourth input.

\smallskip\noindent\textit{Load-bearing status.}  $L_{\mathrm{iso}}$ is used only in $L_{\mathrm{cost}}$ (to obtain the integrality $\eps(d) = n(d) \cdot \eps^*$).  The core derivation chain from $L_\Delta$ through $T_{\mathrm{Born}}$ does \emph{not} require $L_{\mathrm{iso}}$ or $L_{\mathrm{cost}}$'s integrality.  It uses only: $\eps(d) > 0$ (A1), $\eps(d) \ge \eps^* > 0$ ($L_{\eps^*}$, from MD's floor), K3 (from A1), and NT (non-degeneracy of the cost spectrum, which follows from MD's floor alone: distinctions with $n = 1$ cost $\ge \mu^*$ while those with $n = 2$ cost $\ge 2\mu^* > \mu^*$, giving $\ge 2$ distinct values).  Every result from $T_{\mathrm{alg}}$ onward is therefore robust to either reading of MD.
\end{lemma}

\subsection{Finite dimensionality ($T_{\mathrm{form}}$)}

\begin{theorem}[$T_{\mathrm{form}}$: Finite Substrate Dimension]\label{thm:Tform}
At most $\lfloor C(\Gamma)/\mu^* \rfloor$ pairwise independently enforceable irreducible distinctions can coexist at a single interface.  After $T_{\mathrm{embed}}$ constructs $V$, this becomes: $\dim(V) \le C(\Gamma)/\mu^* < \infty$.

\noindent\textit{Proof} (operational; no vector-space structure used).  Let $\{d_1, \ldots, d_k\}$ be a maximal collection of pairwise independently enforceable irreducible distinctions at $\Gamma$.  ``Pairwise independently enforceable'' means: their affected regions are non-overlapping, so K3 gives exact cost additivity.  ``Irreducible'' means $n(d_i) = 1$ (one binary test each).  By MD: $\eps(d_i) \ge \mu^* \cdot 1 = \mu^*$ for each $i$.  By A1: $\sum_{i=1}^k \eps(d_i) \le C(\Gamma)$.  Therefore:
\[
k \cdot \mu^* \;\le\; \sum_{i=1}^k \eps(d_i) \;\le\; C(\Gamma), \qquad\text{hence}\qquad k \;\le\; C(\Gamma)/\mu^*.
\]
This is a purely combinatorial bound on the number of coexisting independent distinctions.  After $T_{\mathrm{embed}}$ identifies each irreducible distinction with a one-dimensional direction in $V$ (via $L_{\mathrm{pred}}$), the bound becomes $\dim(V) \le C(\Gamma)/\mu^*$.  \qed

\noindent\textit{Regularity conditions used:} \fbox{MD} (for $\mu^* > 0$), A1 (budget bound), $L_{\mathrm{pred}}$ (from SP + FD1 + K2 + FD3).

\noindent\textit{Dimension bounds.}  Three explicit per-interface bounds:
\[
\dim(S_\Gamma) \le C/\mu^*, \qquad n_{\max} := \lfloor C/\eps^* \rfloor, \qquad \dim(\mathcal{H}_\omega) \le n_{\max}^2.
\]
All per-interface (local), not cosmological.  QFT: $C \to \infty$ limit.  Bekenstein--Hawking: $n_{\max} \le A/4\ell_P^2$.
\end{theorem}

\subsection{Canonical embedding ($T_{\mathrm{embed}}$)}

\begin{theorem}[$T_{\mathrm{embed}}$: Finite-Dimensional Vector Space]\label{thm:Tembed}
There exists a canonical embedding $\varphi: \Omega(\Gamma) \hookrightarrow \mathbb{R}^{N+1}$ with $N \le \lfloor C/\eps^* \rfloor$, identifying $S_\Gamma$ with $V := \spn_\mathbb{R}(\conv(\varphi(\Omega)))$.

\noindent\textit{Construction.}  Let $\mathcal{B} = \{d_1, \ldots, d_N\}$ be a maximal antichain of pairwise independently enforceable distinctions (finite by $T_{\mathrm{form}}$: $N \le \lfloor C/\eps^*\rfloor$).  Map each state $\sigma = (S_\sigma, \delta_\sigma)$ to:
\[
\varphi(\sigma) := \bigl(\eps(d_1)\mathbf{1}_{d_1 \in S_\sigma},\; \ldots,\; \eps(d_N)\mathbf{1}_{d_N \in S_\sigma},\; \delta_\sigma\bigr) \;\in\; \mathbb{R}^{N+1}.
\]
Injectivity: $\varphi(\sigma) = \varphi(\sigma')$ implies $S_\sigma \cap \mathcal{B} = S_{\sigma'} \cap \mathcal{B}$ (from the indicator coordinates) and $\delta_\sigma = \delta_{\sigma'}$ (from the last coordinate); by the $\mathcal{D}$-quotient, this determines $\sigma = \sigma'$.  The convex hull $K := \conv(\varphi(\Omega))$ is a compact convex body in $\mathbb{R}^{N+1}$ (by OR1).  Define $V := \spn_\mathbb{R}(K)$; then $\dim(V) \le N+1 < \infty$.  \qed

\noindent\textit{Regularity conditions used:} $L_{\eps^*}$ (for finite $N$), OR1 (for convex hull).
\end{theorem}

\begin{lemma}[Operational Independence $\Rightarrow$ Affine Independence of Embedding Coordinates]\label{lem:affine_indep}
If $\{d_1, \ldots, d_k\} \subseteq \mathcal{B}$ are pairwise independently enforceable distinctions in the maximal antichain used by $T_{\mathrm{embed}}$, then their coordinate directions under $\varphi$ are linearly independent in $V$.

\smallskip\noindent\textit{Proof.}  Since $d_1, \ldots, d_k$ have pairwise disjoint anchors ($T_{\mathrm{sep}}$: disjoint anchors $\Leftrightarrow$ additive costs), each $d_i$ can be enforced independently of the others.  For each $i$, construct the state $\sigma_i$ that enforces only $d_i$ (and no other member of $\{d_1, \ldots, d_k\}$): $S_{\sigma_i} = \{d_i\}$, $\delta_{\sigma_i} = C - \eps(d_i) \ge 0$.  This is admissible because the budget suffices ($\eps(d_i) \le C$).  Let $\sigma_0$ be the empty state: $S_{\sigma_0} = \emptyset$, $\delta_{\sigma_0} = C$.  The embedding gives:
\[
\varphi(\sigma_0) = (0, \ldots, 0,\; C), \qquad \varphi(\sigma_i) = (0, \ldots, \underbrace{\eps(d_i)}_{i\text{-th}}, \ldots, 0,\; C - \eps(d_i)).
\]
The difference vectors $v_i := \varphi(\sigma_i) - \varphi(\sigma_0)$ have $v_i = \eps(d_i)\,(e_i - e_{N+1})$, where $e_i$ is the $i$-th standard basis vector and $e_{N+1}$ is the last.  Since $\eps(d_i) > 0$ (by A1) and the $\{e_i - e_{N+1}\}_{i=1}^k$ are linearly independent in $\mathbb{R}^{N+1}$ (they differ in distinct coordinates), the $\{v_i\}_{i=1}^k$ are linearly independent.  Therefore $\dim(V) \ge k$.

Combined with $T_{\mathrm{form}}$'s upper bound $k \le C/\mu^*$: the dimension of $V$ equals the maximum number of pairwise independently enforceable irreducible distinctions, to within the $T_{\mathrm{form}}$ bound.  \qed

\smallskip\noindent\textit{What this lemma does.}  It closes the gap between ``$T_{\mathrm{form}}$ counts operational complexity'' and ``$T_{\mathrm{embed}}$ constructs a vector space of that dimension.''  Without it, the passage from ``at most $k$ independent distinctions'' to ``$\dim(V) \le k$'' relies on an implicit identification.  With it, the identification is an explicit theorem: operational independence maps injectively to linear independence of the embedding coordinates.
\end{lemma}

\begin{remark}[Provenance of the linear structure]\label{rem:linearity}
The vector space $V$ is not presupposed; it is constructed.  The chain is: (1)~A1 provides real-valued costs $\eps(d) \in \mathbb{R}_{>0}$; (2)~OR1 establishes that the state space $\Omega$ is convex (mixtures of admissible states are admissible); (3)~$T_{\mathrm{embed}}$ maps $\Omega$ injectively into $\mathbb{R}^{N+1}$ using these costs as coordinates, producing a compact convex body $K$; (4)~$V := \spn_\mathbb{R}(K)$ is the real vector space spanned by $K$.  The physical substrate $S_\Gamma$ is then identified with $V$ via the $\mathcal{D}$-quotient: the quotient ensures that $S_\Gamma$ has no directions invisible to enforcement (SC, Lemma~\ref{lem:SC}), so every direction in $S_\Gamma$ corresponds to a direction in $V$ and vice versa.  Once $V$ is available, $\End(V)$ is the algebra of all linear maps $V \to V$, and enforcement operations are extended into $\End(V)$ by O1 (Section~\ref{sec:skeleton}).  The bilinear form $B$ on $V$ is constructed later by $L_\omega$ (Lemma~\ref{lem:Lomega}) from K3's sector orthogonality.  No linear, metric, or spectral structure is assumed prior to these explicit construction steps.  In the language of Generalized Probabilistic Theories (GPTs), $V$ is the affine span of the state space $K$ --- the standard starting point of any GPT reconstruction.  The construction does not smuggle in the quantum superposition principle; it performs the same convex-geometric embedding that Hardy (2001), CDP (2011), and Masanes--M\"uller (2011) all use as their opening move.
\end{remark}

\begin{remark}[Bootstrapping verification: no pre-$T_{\mathrm{embed}}$ proof uses linear structure]\label{rem:bootstrap_verify}
The definitions FD4 and K1--K3 are formulated in linear language ($\End(V)$, supports, direct sums) for mathematical precision.  However, \textbf{no result proved before $T_{\mathrm{embed}}$} --- specifically, OR1, $L_{\eps^*}$, $L_{\mathrm{pred}}$, and $T_{\mathrm{form}}$ --- invokes $\End(V)$, direct sums, operator norms, or linear independence.  In each case:
\begin{itemize}[nosep,leftmargin=2em,itemsep=2pt]
\item OR1 uses only A1's real-valued budget (a numerical inequality).
\item $L_{\eps^*}$ uses only MD's floor (a numerical bound on $\eps$).
\item $L_{\mathrm{pred}}$ uses SP (configurations differ operationally) and FD1 (the predicate exists).
\item $T_{\mathrm{form}}$ counts the maximum number of pairwise independently enforceable distinctions (a combinatorial bound); the identification with $\dim(V)$ occurs only after $T_{\mathrm{embed}}$.
\end{itemize}
After $T_{\mathrm{embed}}$ constructs $V$, the operational notions (affected regions, independence, joint cost) become their linear counterparts (supports, disjoint subspaces, K3 on $\End(V)$) via the embedding's injectivity and the $\mathcal{D}$-quotient.  This equivalence is spelled out in the operational-to-linear table in FD4 (Section~\ref{sec:defs}).  All results from $T_{\mathrm{sep}}$ onward use the linear formulation legitimately.
\end{remark}

\subsection{Sector decomposition ($T_{\mathrm{sep}}$)}

The sector decomposition is first stated in purely operational language (no vector space, no linear algebra), then represented linearly after $T_{\mathrm{embed}}$.

\begin{theorem}[$T_{\mathrm{sep}}^{\mathrm{op}}$: Operational Sector Decomposition]\label{thm:Tsep_op}
Let $\Gamma$ be an enforcement interface with distinction set $\mathcal{D}(\Gamma)$ and enforcement costs $\eps(d) > 0$ (A1).  Define two distinctions $d_1, d_2$ as \emph{independently enforceable} if no perturbation threatening $d_1$ can affect $d_2$'s enforcement status, and vice versa --- equivalently, the affected regions $\mathrm{Aff}(d_1)$ and $\mathrm{Aff}(d_2)$ (the sets of substrate degrees of freedom on which perturbations threatening $d_i$ operate) are disjoint: $\mathrm{Aff}(d_1) \cap \mathrm{Aff}(d_2) = \emptyset$.

Then:
\begin{enumerate}[nosep,label=(\alph*),leftmargin=2em]
\item \textit{Cost criterion.}  $d_1$ and $d_2$ are independently enforceable if and only if their joint enforcement cost equals the sum of their individual costs: $\eps(\{d_1, d_2\}) = \eps(d_1) + \eps(d_2)$.
\item \textit{Pairwise partition.}  Any maximal family $\mathcal{B} \subseteq \mathcal{D}(\Gamma)$ of pairwise independently enforceable distinctions (a maximal antichain under cost-independence) partitions the substrate degrees of freedom into: anchor classes $\mathrm{Aff}(d)$ for each $d \in \mathcal{B}$ (pairwise disjoint), and a residual \emph{pool} $\Pi_{\mathrm{op}}$ --- the set of substrate DOF that participate in the enforcement of multiple distinctions but are not contained in the anchor class of any single $d \in \mathcal{B}$.
\item \textit{Pool characterization.}  $\Pi_{\mathrm{op}} = \emptyset$ iff the total capacity is exactly saturated by the anchor classes (classical regime).  $\Pi_{\mathrm{op}} \ne \emptyset$ iff there exist substrate DOF that are jointly engaged in defending multiple distinctions against perturbations; these DOF mediate cross-distinction coupling and are the source of superadditivity ($L_\Delta$).
\end{enumerate}

\smallskip\noindent\textit{Proof of (a).}  ($\Rightarrow$): If affected regions are disjoint, perturbations threatening $d_1$ and $d_2$ operate on non-overlapping DOF.  The joint defense decomposes: the defense of $d_1$ commits capacity only within $\mathrm{Aff}(d_1)$, and similarly for $d_2$.  By K3 (additivity on disjoint supports, Theorem~\ref{thm:K3}): $\eps(\{d_1,d_2\}) = \eps(d_1) + \eps(d_2)$.

($\Leftarrow$, contrapositive): If some DOF $x$ lies in $\mathrm{Aff}(d_1) \cap \mathrm{Aff}(d_2)$, then defending $d_1$ and $d_2$ separately each commits capacity to $x$; defending jointly commits to $x$ only once.  By SP and K2, the shared commitment has positive cost.  Therefore $\eps(\{d_1,d_2\}) < \eps(d_1) + \eps(d_2)$.

\noindent Part (b) follows by iterated application of (a) to a maximal antichain.  Part (c) restates the $\Pi_{\mathrm{op}} = \emptyset$ criterion in cost terms.  \qed

\smallskip\noindent\textit{What this theorem uses.}  A1 ($\eps(d) > 0$, finite capacity), K3 (disjoint-support additivity), FD3 (anchor-set locality), SP (non-null substrate DOF have positive perturbation cost).  No vector space, inner product, linear map, or subspace concept is invoked.
\end{theorem}

\begin{theorem}[$T_{\mathrm{sep}}$: Linear Representation of the Sector Decomposition]\label{thm:Tsep}
Under $T_{\mathrm{embed}}$ (which maps the operational substrate into a finite-dimensional real vector space $V$), the operational sector decomposition of Theorem~\ref{thm:Tsep_op} is represented as: for $d_1, d_2 \in \mathcal{D}(\Gamma)$,
\[
M_{d_1} \cap M_{d_2} = \{0\} \quad\Longleftrightarrow\quad \eps(\{d_1, d_2\}) = \eps(d_1) + \eps(d_2),
\]
where $M_d := \varphi(\mathrm{Aff}(d)) \subseteq V$ is the image of the anchor class under the embedding $\varphi$.  The operational partition becomes an algebraic direct sum:

\smallskip\noindent\textit{Proof (linear representation).}  The operational proof (Theorem~\ref{thm:Tsep_op}(a)) carries over under $T_{\mathrm{embed}}$: operational affected regions become linear supports ($\mathrm{Aff}(d) \mapsto M_d$), disjointness becomes trivial intersection, and operational K3 becomes linear K3 on $\End(V)$.  For completeness, we restate in linear terms.

\textit{($\Rightarrow$: disjoint anchors imply additive costs).}  Suppose $M_{d_1} \cap M_{d_2} = \{0\}$.  We must show $\eps(\{d_1, d_2\}) = \eps(d_1) + \eps(d_2)$.

\textit{Step (a): perturbation classes separate.}  By FD3, every perturbation $p \in \mathcal{P}(d_i)$ has $\mathrm{supp}(p) \subseteq M_{d_i}$.  Since $M_{d_1} \cap M_{d_2} = \{0\}$, any $p_1 \in \mathcal{P}(d_1)$ and $p_2 \in \mathcal{P}(d_2)$ have disjoint supports: $\mathrm{supp}(p_1) \cap \mathrm{supp}(p_2) = \{0\}$.

\textit{Step (b): joint defense decomposes.}  Defending $\{d_1, d_2\}$ simultaneously means preventing all $p \in \mathcal{P}(d_1) \cup \mathcal{P}(d_2)$ from destroying either distinction.  Since the support sets are disjoint (Step~a), a perturbation threatening $d_1$ cannot affect $M_{d_2}$ and vice versa.  The defense of $d_1$ requires committing capacity only within $M_{d_1}$; the defense of $d_2$ only within $M_{d_2}$.  No joint defense coordination is needed.

\textit{Step (c): costs add by K3.}  The optimal defense of $\{d_1, d_2\}$ is the direct sum of the individual optimal defenses: $\kappa_{\mathrm{opt}}(\{d_1,d_2\}) = \kappa_{\mathrm{opt}}(d_1) \oplus \kappa_{\mathrm{opt}}(d_2)$, where the two components have disjoint supports.  By K3 (additivity on disjoint supports): $\eps(\{d_1,d_2\}) = \inf \kappa = \kappa(\kappa_{\mathrm{opt}}(d_1)) + \kappa(\kappa_{\mathrm{opt}}(d_2)) = \eps(d_1) + \eps(d_2)$.

\smallskip
\textit{($\Leftarrow$: additive costs imply disjoint anchors).}  We prove the contrapositive: $M_{d_1} \cap M_{d_2} \ne \{0\}$ implies $\eps(\{d_1, d_2\}) < \eps(d_1) + \eps(d_2)$ (strict subadditivity).

Let $v \in M_{d_1} \cap M_{d_2}$ with $v \ne 0$.  By FD3, $v$ lies in the support of perturbations threatening both $d_1$ and $d_2$.  Consider defending $\{d_1, d_2\}$ jointly.  The defense must cover all perturbations in $\mathcal{P}(d_1) \cup \mathcal{P}(d_2)$.  When defending $d_1$ and $d_2$ separately, the direction $v$ is defended twice: once as part of the $M_{d_1}$-defense (at cost contribution $\ge \eps^*$ attributable to $\spn\{v\}$) and once as part of the $M_{d_2}$-defense (another $\ge \eps^*$).  In the joint defense, a single commitment to defend $\spn\{v\}$ covers perturbations from both $\mathcal{P}(d_1)$ and $\mathcal{P}(d_2)$ that are supported on $\spn\{v\}$, because the defense is a commitment within the substrate direction $v$ regardless of which distinction it protects.

Formally: write $M_{d_1} = \spn\{v\} \oplus M_{d_1}'$ and $M_{d_2} = \spn\{v\} \oplus M_{d_2}'$ where $M_{d_i}' := M_{d_i} \cap \spn\{v\}^\perp$ (algebraic complement within $M_{d_i}$; no inner product needed, just any complement).  The separate defenses cost $\eps(d_1) = \kappa_{\mathrm{opt}}(M_{d_1}) \ge \kappa(\spn\{v\}) + \kappa(M_{d_1}')$ and $\eps(d_2) \ge \kappa(\spn\{v\}) + \kappa(M_{d_2}')$.  A joint defense that covers $\spn\{v\} \cup M_{d_1}' \cup M_{d_2}'$ costs $\kappa(\spn\{v\}) + \kappa(M_{d_1}') + \kappa(M_{d_2}')$ by K3 (the three subspaces have pairwise trivial intersection).  Therefore:
\[
\eps(\{d_1,d_2\}) \;\le\; \kappa(\spn\{v\}) + \kappa(M_{d_1}') + \kappa(M_{d_2}') \;\le\; \eps(d_1) + \eps(d_2) - \kappa(\spn\{v\}).
\]
By SP and K2, $\kappa(\spn\{v\}) > 0$ (the overlap direction is non-null and perturbations along it have positive cost).  So $\eps(\{d_1,d_2\}) < \eps(d_1) + \eps(d_2)$: strict subadditivity.  The contrapositive gives: additive costs $\Rightarrow$ disjoint anchors.  \qed

\smallskip\noindent\textit{Consequence: algebraic direct sum.}  For any maximal antichain $\mathcal{B} \subseteq \mathcal{D}(\Gamma)$ (pairwise independently enforceable), the anchor sets $\{M_d : d \in \mathcal{B}\}$ are pairwise disjoint subspaces.  Their sum is direct:
\[
S_\Gamma \;=\; \bigoplus_{d \in \mathcal{B}} M_d \;\oplus\; \Pi,
\]
where $\Pi := S_\Gamma \setminus \bigoplus_{d \in \mathcal{B}} M_d$ is the residual pool.  This is a \textbf{pre-metric algebraic direct sum}: no inner product is used.  The pool $\Pi$ contains substrate directions that participate in the enforcement of multiple distinctions but are not the anchor of any single member of $\mathcal{B}$.

\noindent\textit{Regularity conditions used:} A1 (for $\eps(d_v) > 0$), K3 (from A1, Definition~\ref{def:FD4}), FD3 (anchor-set locality).
\end{theorem}

\begin{definition}[Pool Sector $\Pi$]\label{def:pool}
Given the $T_{\mathrm{sep}}$ decomposition $S_\Gamma = \bigoplus_{d \in \mathcal{B}} M_d \oplus \Pi$, the \emph{pool sector} $\Pi$ is the complementary subspace of the direct sum of all single-distinction anchor sets.  Concretely: $\Pi$ contains all substrate directions at $\Gamma$ that are physically meaningful (non-null by SP) but not attributable to the anchor of any single distinction in the maximal antichain $\mathcal{B}$.  The pool is \emph{not} an axiom; it is a structural consequence of $T_{\mathrm{sep}}$ whenever $\dim(S_\Gamma) > \sum_d \dim(M_d)$.

\smallskip\noindent\textit{When $\Pi$ matters.}  If $\Pi = \{0\}$ (capacity perfectly saturated by individual anchors), the framework is classical: all costs are additive, all enforcement projections commute, the algebra $\mathcal{A}_{\mathrm{diag}} \cong \mathbb{R}^{|\mathcal{B}|}$ is the full algebra.  If $\Pi \ne \{0\}$ (i.e., the substrate is not exhausted by single-distinction anchors), the pool mediates cross-sector coupling, joint enforcement costs are superadditive ($L_\Delta$: $\Delta > 0$), and the full algebra is noncommutative ($T_{\mathrm{alg}}$).
\end{definition}

\begin{proposition}[Structural Characterization of the Pool Sector]\label{prop:pool}
Let $\mathcal{B}$ be a maximal antichain of independently enforceable distinctions at $\Gamma$, and let $\Pi$ be the pool sector from $T_{\mathrm{sep}}$ (Definition~\ref{def:pool}).
\begin{enumerate}[nosep,label=(\roman*),leftmargin=2em]
\item \textit{Operational definition.}  $\Pi_{\mathrm{op}}$ consists of all substrate DOF at $\Gamma$ that are physically meaningful (non-null by SP) but not attributable to the anchor of any single distinction in $\mathcal{B}$.  Equivalently: a DOF $x$ lies in $\Pi_{\mathrm{op}}$ iff $x \notin \mathrm{Aff}(d)$ for any $d \in \mathcal{B}$, yet perturbations along $x$ have positive cost ($\kappa(p_x) > 0$ by K2).
\item \textit{Uniqueness.}  The partition $\{(\mathrm{Aff}(d))_{d \in \mathcal{B}},\, \Pi_{\mathrm{op}}\}$ is unique once $\mathcal{B}$ is fixed: the anchor classes are determined by FD3, and $\Pi_{\mathrm{op}}$ is the complement.  Under $T_{\mathrm{embed}}$: $\Pi = V \setminus \bigoplus_d M_d$, which is unique as an algebraic complement.
\item \textit{Classical criterion: $\Pi = \{0\}$.}  $\Pi = \{0\}$ iff $\dim(V) = \sum_d \dim(M_d)$, i.e., every substrate DOF is the anchor of exactly one distinction.  In this case: all enforcement costs are additive ($L_\Delta$ gives $\Delta = 0$), all enforcement projections commute, and $\mathcal{A} = \mathcal{A}_{\mathrm{diag}} \cong \mathbb{R}^{|\mathcal{B}|}$ (classical).
\item \textit{Quantum criterion: $\Pi \ne \{0\}$.}  $\Pi \ne \{0\}$ iff $\dim(V) > \sum_d \dim(M_d)$, i.e., there exist DOF jointly engaged in defending multiple distinctions.  In this case: $L_\Delta$ gives $\Delta > 0$, $L_{\mathrm{blk}}$ forces codespace rotation through $\Pi$, and $T_{\mathrm{alg}}$ produces noncommutativity.  This is generic but not automatic: it requires that the substrate contain DOF beyond the union of single-distinction anchors.
\item \textit{Logical role.}  $\Pi$ mediates the transition from classical to quantum structure.  It enters the derivation chain at three points: $L_\Delta$ (superadditivity of co-located costs, via $\Pi$-supported perturbations); $L_{\mathrm{blk}}$ (failure of block-diagonal defense, via $\Pi$-supported attacks); and $T_{\mathrm{alg}}$ (noncommutativity, via the rotated codespace $W_* \not\subseteq M_{d_1} \oplus M_{d_2}$).
\end{enumerate}

\smallskip\noindent\textit{Conditions used.}  A1 (finite capacity), K3 (disjoint-support additivity), FD3 (anchor locality), SP ($\mathcal{D}$-quotient: non-null DOF have positive cost).  Under $T_{\mathrm{embed}}$, the operational characterization becomes the linear one.
\end{proposition}

\begin{tcolorbox}[apfbox, title={Three Regimes of Joint Enforcement Cost}]
The sector decomposition and pool structure yield a complete taxonomy of joint enforcement costs for two co-located distinctions $d_1, d_2$ at a single interface $\Gamma$:

\smallskip\noindent\textbf{Regime 1: Anchor overlap $\Rightarrow$ subadditivity.}  If $M_{d_1} \cap M_{d_2} \ne \{0\}$, then $\eps(\{d_1, d_2\}) < \eps(d_1) + \eps(d_2)$.  The shared anchor DOF are defended once instead of twice.  \textit{Source: $T_{\mathrm{sep}}$, converse direction.}

\smallskip\noindent\textbf{Regime 2: Disjoint anchors, no pool burden $\Rightarrow$ exact additivity.}  If $M_{d_1} \cap M_{d_2} = \{0\}$ and $\Pi = \{0\}$ (or no $\Pi$-supported perturbation threatens the joint enforcement), then $\eps(\{d_1, d_2\}) = \eps(d_1) + \eps(d_2)$.  Defenses decompose; costs add by K3.  This is the classical regime.  \textit{Source: $T_{\mathrm{sep}}$, forward direction.}

\smallskip\noindent\textbf{Regime 3: Disjoint anchors, nonzero pool $\Rightarrow$ superadditivity.}  If $M_{d_1} \cap M_{d_2} = \{0\}$ and $\Pi \ne \{0\}$ with $\Pi$-supported perturbations threatening the joint enforcement, then $\eps(\{d_1, d_2\}) > \eps(d_1) + \eps(d_2)$.  The pool mediates cross-sector coupling that is invisible to individual enforcement but must be defended jointly.  The surplus $\Delta := \eps(\{d_1,d_2\}) - \eps(d_1) - \eps(d_2) \ge \eps^*$ is the signature of nonclassicality.  \textit{Source: $L_\Delta$ (Theorem~\ref{thm:LDelta}).}

\smallskip\noindent The bridge to noncommutativity uses Regime~3 exclusively: $\Delta > 0$ $\to$ $L_{\mathrm{blk}}$ $\to$ $L_\Pi$ $\to$ $T_{\mathrm{alg}}$.

\smallskip\noindent\textit{Classical countermodel ($\Pi = 0$, noncommutativity fails).}  Take $S_\Gamma = \mathbb{R}^2$, $M_{d_1} = \spn\{e_1\}$, $M_{d_2} = \spn\{e_2\}$, $\Pi = \{0\}$ (the sector sum exhausts the substrate).  The enforcement projections are $E_{d_1} = \diag(1,0)$, $E_{d_2} = \diag(0,1)$.  Joint enforcement: $E_{d_1} + E_{d_2} = \mathrm{Id}$; the joint defender is the identity, with cost $\eps(\{d_1,d_2\}) = \eps(d_1) + \eps(d_2)$ (exact additivity, $\Delta = 0$).  All generators commute: $[E_{d_1}, E_{d_2}] = 0$.  The algebra $\mathcal{A} = \mathcal{A}_{\mathrm{diag}} \cong \mathbb{R}^2$ is commutative.  $T_{\mathrm{alg}}$ does not apply because $\Pi = \{0\}$ means CL3 ($\Pi \ne \{0\}$) fails.  This is Regime~2: classical, additive, commutative.  The bridge chain ($L_\Delta \to L_{\mathrm{blk}} \to T_{\mathrm{alg}}$) correctly produces no noncommutativity when $\Pi = \{0\}$.
\end{tcolorbox}

\subsection{Cost functional uniqueness ($L_{\mathrm{cost}}$)}

\begin{theorem}[$L_{\mathrm{cost}}$: Cost Functional Form Uniqueness]\label{thm:Lcost}
\textbf{Status: conditional on $L_{\mathrm{iso}}$.}  The form $\eps(d) = n(d) \cdot \eps^*$ requires H1 (ledger completeness), which depends on $L_{\mathrm{iso}}$ (Lemma~\ref{lem:iso}).  Under the weak reading of MD (floor only), H1 is not established and $L_{\mathrm{cost}}$'s integrality does not follow.  The core derivation chain ($L_\Delta \to T_{\mathrm{alg}} \to T_{\mathrm{GNS}} \to T_{\mathrm{Born}}$) does not invoke $L_{\mathrm{cost}}$; it uses only $\eps > 0$ (A1), $\eps \ge \eps^*$ ($L_{\eps^*}$), K3, and NT (below).  The integrality result is available under the strong reading and is used in companion papers for capacity counting; it is not load-bearing for any result in Part~I's core chain.

\smallskip\noindent\textbf{Formal statement.}  Let $\mathcal{E}: \{\text{finite subsets of } \mathcal{D}(\Gamma) \text{ with pairwise disjoint anchors}\} \to \mathbb{R}_{>0}$ be the enforcement cost functional, defined by $\mathcal{E}(S) := \sum_{d \in S} \eps(d)$ for any pairwise independently enforceable set $S$.  Under the following three hypotheses:
\begin{enumerate}[nosep,label=(H\arabic*),leftmargin=3em]
\item \textit{Ledger completeness} (from $L_{\mathrm{iso}}$, conditional on strong MD): cost depends only on the channel count $n := |S|$, so $\mathcal{E}(S) = f(n)$ for some $f: \mathbb{N} \to \mathbb{R}_{>0}$;
\item \textit{Integrality} (from FD2): $n \in \mathbb{N}$ (each distinction is binary: active or not);
\item \textit{Additive independence} (from $T_{\mathrm{sep}}$): for independently enforceable $S_1, S_2$ with $S_1 \cap S_2 = \emptyset$, $f(|S_1| + |S_2|) = f(|S_1|) + f(|S_2|)$;
\end{enumerate}
the cost functional has the unique form $\mathcal{E}(S) = |S| \cdot \eps^*$, where $\eps^* := f(1) > 0$.

\smallskip\noindent\textit{Proof.}

\textit{H1 (Ledger Completeness).}  By $L_{\mathrm{iso}}$ (Lemma~\ref{lem:iso}, conditional on strong MD), every irreducible distinction at $\Gamma$ has cost $\eps^*$, and every $k$-channel distinction has cost $k\eps^*$.  For a set $S$ of independently enforceable irreducible distinctions, $\mathcal{E}(S) = |S| \cdot \eps^*$.  Therefore $\mathcal{E}$ depends only on $|S|$: $\mathcal{E}(S) = f(|S|)$ with $f(n) = n \cdot \eps^*$.

\textit{H2 (Integrality).}  Each distinction is a binary predicate (FD2: active or not), so the channel count $n = |S_\sigma| \in \mathbb{N}$.  The domain of $f$ is $\mathbb{N}$, not $\mathbb{R}_{\ge 0}$.

\textit{H3 (Additive Independence).}  For independently enforceable $S_1, S_2$ (disjoint anchors, from $T_{\mathrm{sep}}$): $f(n_1 + n_2) = f(n_1) + f(n_2)$.

\textit{Cauchy uniqueness on $\mathbb{N}$.}  The Cauchy functional equation $f(m+n) = f(m) + f(n)$ on $\mathbb{N}$ has a unique solution: $f(n) = n \cdot f(1)$, proved by induction ($f(n) = f(n-1) + f(1) = \cdots = n \cdot f(1)$).  Set $\eps^* := f(1) > 0$ (from $L_{\eps^*}$: MD gives $f(1) \ge \mu^* > 0$).  No regularity condition (measurability, monotonicity) is needed beyond integrality: the domain $\mathbb{N}$ has a single additive generator, eliminating Hamel-basis pathologies that plague Cauchy equations on $\mathbb{R}$.  \qed

\begin{proposition}[Admissible Domain-Specific Proxies]\label{prop:proxies}
A domain-specific cost measure $c: \mathcal{D}(\Gamma) \to \mathbb{R}_{\ge 0}$ induces a valid enforcement cost scale (i.e., $c$ can serve as $\eps$) if and only if it satisfies:
\begin{enumerate}[nosep,label=(\roman*),leftmargin=2.5em]
\item \textit{Positivity:} $c(d) > 0$ for every $d \in \mathcal{D}(\Gamma)$;
\item \textit{Uniform floor:} $c(d) \ge \mu^* \cdot n(d)$ for some $\mu^* > 0$ (i.e., MD);
\item \textit{Additive independence:} $c(d_1 \oplus d_2) = c(d_1) + c(d_2)$ for disjoint-anchor pairs (from K3).
\end{enumerate}
Under these conditions, $L_{\mathrm{cost}}$ gives $c(d) = n(d) \cdot c_*$ with $c_* := c(d_{\mathrm{irr}}) > 0$ for any irreducible distinction $d_{\mathrm{irr}}$.  The functional form is universal; the numerical value of $c_*$ depends on the proxy.

\smallskip\noindent\textit{Examples of admissible proxies.}
(a)~\textit{Spectral gap:} the energy gap $\Delta E(d)$ protecting distinction $d$ from thermal fluctuations, provided independent subsystems have additive gap costs.
(b)~\textit{Error-correction overhead:} the number of physical qubits per logical qubit required to protect $d$, when code overheads for independent sectors sum.
(c)~\textit{Entropy / free-energy cost:} the Landauer erasure cost $k_B T \ln 2$ per bit of distinction, additive on independent supports.
(d)~\textit{Geometric / area-law budget:} interface area $A(d)$ allocated to distinction $d$, additive over disjoint interface pieces (Bekenstein--Hawking regime).

\noindent Each proxy is admissible only when it respects null identification (the $\mathcal{D}$-quotient) and additive independence (K3); otherwise it is not a valid realization of $\eps$ in the present framework.  \qed

\smallskip\noindent\textit{A complete physical realization: surface-code syndrome overhead.}
To show the framework is not merely abstract, we verify K1--K4 for one concrete physical cost.  In a distance-$d$ surface code on $2d^2 - 2d + 1$ physical qubits, each weight-4 stabilizer $S_j$ maintains one parity distinction $d_j$ (``syndrome bit $j$ is 0'').  The physical enforcement cost is the syndrome extraction overhead per round: $\gamma_j = t_{\mathrm{gate}} \cdot w_j + t_{\mathrm{meas}}$, where $w_j = 4$ (stabilizer weight), $t_{\mathrm{gate}} \approx 30\,\mathrm{ns}$, $t_{\mathrm{meas}} \approx 500\,\mathrm{ns}$ (Google Sycamore data), giving $\gamma_j \approx 620\,\mathrm{ns}$.  Verification: (K1)~$\gamma_j > 0$ for every non-trivial stabilizer.  (K2)~$\gamma = 0$ only for the identity ($w = 0$).  (K3)~For non-overlapping stabilizers ($M_{d_j} \cap M_{d_k} = \{0\}$): extraction rounds are parallelizable, $\gamma(\{S_j, S_k\}) = \gamma(S_j) + \gamma(S_k)$.  (K4)~$\gamma$ is continuous in the gate parameters.  All four hold above the error-correction threshold ($p < p_{\mathrm{th}}$).  Below threshold, V3 fails (correlated errors break additive independence), and the framework correctly predicts that quantum structure is not maintained.  The proportionality $\gamma_j = \alpha_\gamma \cdot \eps(d_j)$ with $\alpha_\gamma = 620\,\mathrm{ns}/\eps^*$ is confirmed by $L_{\mathrm{cost}}$.  A dual-ruler cross-check (energy gap $\Delta E$ as a second proxy) is carried out in the companion paper, confirming that the cross-ruler ratio $\gamma/\Delta E$ is distinction-independent as predicted.
\end{proposition}

\noindent\textit{Corollary (NT: Non-Degeneracy).}  The cost spectrum has $\ge 2$ distinct values.  \textit{Proof (from MD's floor, independent of $L_{\mathrm{cost}}$).}  At any interface with $\dim(S_\Gamma) \ge 2$, there exist distinctions with $n(d) = 1$ (irreducible) and $n(d) = 2$ (two-channel composite).  By MD: $\eps(d_1) \ge \mu^*$ for the irreducible, and $\eps(d_2) \ge 2\mu^*$ for the composite.  Since $2\mu^* > \mu^*$, the cost spectrum has $\ge 2$ distinct values.  NT requires only the floor from MD, not isotropy.  \qed

\noindent\textit{Regularity conditions used:} A1, \fbox{MD} (floor only; $L_{\mathrm{iso}}$ not required).
\end{theorem}

\subsection{Budget additivity for independent interfaces ($L_{\mathrm{loc}}$)}

\begin{theorem}[$L_{\mathrm{loc}}$: Locality]\label{thm:Lloc}
For independent interfaces $\Gamma_1, \Gamma_2$ (disjoint substrates: $S_{\Gamma_1} \cap S_{\Gamma_2} = \emptyset$):
\[
C(\Gamma_1 \cup \Gamma_2) \;=\; C(\Gamma_1) + C(\Gamma_2).
\]
No capacity can be transferred between independent interfaces.

\noindent\textit{Proof.}  By the disjoint-substrate condition, every distinction $d \in \mathcal{D}(\Gamma_1 \cup \Gamma_2)$ has its anchor set entirely within $S_{\Gamma_1}$ or entirely within $S_{\Gamma_2}$ (a distinction whose anchor spanned both substrates would require enforcement in both, but $M_d$ is a subspace of a single substrate by FD3's locality of anchor sets).  The cost sum partitions: $\sum_d \eps(d) = \sum_{d \in \mathcal{D}(\Gamma_1)} \eps(d) + \sum_{d \in \mathcal{D}(\Gamma_2)} \eps(d) \le C(\Gamma_1) + C(\Gamma_2)$.  The bound is tight (each interface's budget can be independently saturated), so $C(\Gamma_1 \cup \Gamma_2) = C(\Gamma_1) + C(\Gamma_2)$.  \qed

\noindent\textit{Regularity conditions used:} A1 (budget constraint), FD3 (anchor-set locality).  \textbf{Not used:} MD, $L_{\eps^*}$, $L_{\mathrm{iso}}$, $L_{\mathrm{cost}}$, integrality, isotropy.  The proof requires only A1's budget constraint and FD3's anchor-set locality --- the two weakest ingredients in the entire framework.  Budget additivity for independent interfaces is therefore robust to any strengthening or weakening of MD.  In particular, $L_{\mathrm{loc}}$ does not depend on MD, $L_{\mathrm{iso}}$, integrality, or any cost-functional structure beyond A1's capacity bound.  This is the minimal-hypothesis locality theorem.
\end{theorem}

\subsection{Consolidation: Finite Orthogonal Sector Form}

The preceding results combine into a single structural theorem that serves as the input to the bridge chain (Section~\ref{sec:bridge}).

\begin{tcolorbox}[apfbox, title={$T_{\mathrm{sector}}$: Finite Orthogonal Sector Form of an Admissible Interface}]
\textbf{Statement.}  Let $\Gamma$ be an interface satisfying A1, MD, SP, and K3.  Then there exists a finite-dimensional real vector space $V$ with $\dim(V) \le C(\Gamma)/\mu^*$ and a decomposition
\[
V \;=\; \Bigl(\bigoplus_{d \in \mathcal{B}} M_d\Bigr) \;\oplus\; \Pi,
\]
with the following properties:
\begin{enumerate}[nosep,leftmargin=2em,label=(\roman*)]
\item \textit{Finite dimension:} $\dim(V) \le C/\mu^* < \infty$ \textup{(}from $T_{\mathrm{form}}$, using A1 + MD + $L_{\mathrm{pred}}$\textup{)}.
\item \textit{Sector decomposition:} $M_d \cap M_{d'} = \{0\}$ for $d \ne d'$, and each $M_d \ne 0$ is the anchor subspace of distinction $d$ \textup{(}from $T_{\mathrm{sep}}$, using A1 + FD3 + K3\textup{)}.
\item \textit{Residual pool:} $\Pi$ is the complementary subspace not assigned to any single distinction's anchor.
\item \textit{Block-orthogonal inner product:} there exists a positive-definite bilinear form $B$ on $V$, canonical in its inter-sector orthogonality pattern, such that
\[
B(u_d, v_{d'}) = 0 \quad (d \ne d'), \qquad B(u_d, v_\Pi) = 0 \quad \text{for all } d.
\]
\textit{Derivation:} K3 forces exact cost additivity across disjoint sectors.  Representing cost by $B$-norm: for $u_d \in M_d$, $u_{d'} \in M_{d'}$ with $d \ne d'$,
\[
\|u_d + u_{d'}\|_B^2 = \|u_d\|_B^2 + \|u_{d'}\|_B^2 \quad\Longrightarrow\quad 2\,B(u_d, u_{d'}) = 0.
\]
Cross-sector orthogonality is forced; within-sector normalization is free.
\item \textit{Diagonal algebra:} the sector maps $E_d \in \End(V)$ defined by $E_d|_{M_d} = \mathrm{id}$, $E_d|_{M_{d'}} = 0$ $(d' \ne d)$, $E_d|_\Pi = 0$ satisfy $E_d^2 = E_d$, $E_d E_{d'} = 0$ $(d \ne d')$, and generate a commutative diagonal subalgebra $\mathcal{A}_{\mathrm{diag}} := \spn_\mathbb{R}\{E_d\} \cong \mathbb{R}^{|\mathcal{B}|}$.
\end{enumerate}

\smallskip\noindent\textit{This is the complete output of the arena chain.  No noncommutativity, no Hilbert space, no field selection is claimed here.  Those are derived in the bridge (Section~\ref{sec:bridge}) and skeleton (Section~\ref{sec:skeleton}).}
\end{tcolorbox}

\begin{theorem}[$T_{\mathrm{adj}}$: Self-Adjointness of Sector Projections]\label{thm:Tadj_sector}
With respect to the block-orthogonal inner product $B$ from $T_{\mathrm{sector}}$, each sector map $E_d$ is an orthogonal projection:
\[
E_d^\dagger = E_d, \qquad E_d^2 = E_d.
\]

\noindent\textit{Proof.}  The range of $E_d$ is $M_d$.  The kernel is $\ker(E_d) = (\bigoplus_{d' \ne d} M_{d'}) \oplus \Pi$.  By $T_{\mathrm{sector}}$(iv), $M_d \perp \ker(E_d)$ with respect to $B$.  Therefore $E_d$ is the $B$-orthogonal projection onto $M_d$.  Self-adjointness: for $u = u_d + u_\perp$, $v = v_d + v_\perp$ with $u_d, v_d \in M_d$ and $u_\perp, v_\perp \in \ker(E_d)$:
\[
B(E_d u, v) = B(u_d, v_d + v_\perp) = B(u_d, v_d) = B(u_d + u_\perp, v_d) = B(u, E_d v).
\]
Hence $E_d^\dagger = E_d$.  \qed

\smallskip\noindent\textit{What $T_{\mathrm{adj}}$ uses.}  The proof uses only the $B$-orthogonal sector decomposition ($T_{\mathrm{sector}}$(iv) from $L_\omega$.5).  It does not invoke cost optimization, does not reference a specific $\kappa$, and does not require the physical-selection argument below.  Self-adjointness is a structural consequence of the forced orthogonality pattern.

\smallskip\noindent\textit{Structural invariance across the admissible $\kappa$-class.}  Any cost function satisfying K1--K3 produces the same inter-sector orthogonality (different $\kappa$ may yield different within-sector normalizations of $B$, but $T_{\mathrm{adj}}$'s proof uses only the inter-sector pattern).  Therefore $T_{\mathrm{adj}}$ is structurally invariant across the entire admissible $\kappa$-class: the same projections, the same algebra, the same downstream results.

\smallskip\noindent\textit{Scope.}  $T_{\mathrm{adj}}$ establishes self-adjointness for the \emph{sector projections} $\{E_d\}$, i.e., the generators of the commutative diagonal subalgebra $\mathcal{A}_{\mathrm{diag}}$.  It does \emph{not} claim that the full enforcement algebra is commutative or diagonal.  Later results ($L_\Delta \to L_\Pi \to T_{\mathrm{alg}}$) show that joint-enforcement generators forced by $\Delta > 0$ lie outside $\mathcal{A}_{\mathrm{diag}}$ and make the full algebra noncommutative.
\end{theorem}

The following proposition gives the physical interpretation of $T_{\mathrm{adj}}$ as a cost-minimization result.  It is logically independent of the structural theorem above: $T_{\mathrm{adj}}$ does not depend on it, and it is not used in any later proof.  Its role is to connect the structural output (orthogonal projections) to the enforcement semantics (cost minimization).

\begin{proposition}[$\kappa$-Classification for the Framework]\label{prop:kappa_class}
Let $\kappa: \End(V) \to \mathbb{R}_{\ge 0}$ be a perturbation cost function.

\smallskip\noindent\textbf{Uniqueness theorem.}  If $\kappa$ satisfies K1--K3, then the $B$-orthogonal projection $E_d = \pi_{M_d}^{(B)}$ is the \emph{unique} cost-minimizing idempotent with $\im(F) = M_d$.  \textit{Proof:} let $F$ be any idempotent with $\im(F) = M_d$.  Write $F = E_d + L_F$ where $L_F: \Pi \to M_d$ is the spillover (the $M_d$-component of $F|_\Pi$).  If $L_F \ne 0$: $\mathrm{supp}(L_F) \subseteq \Pi$ and $\mathrm{supp}(E_d) \subseteq \bigoplus_{d' \ne d} M_{d'}$ are disjoint, so K3 gives $\kappa(F) \ge \kappa(E_d) + \kappa(L_F)$.  By K2, $\kappa(L_F) > 0$ (nonzero map on non-null subspace).  Therefore $\kappa(F) > \kappa(E_d)$.  If $L_F = 0$: $F = E_d$.  The minimizer is unique and equals the $B$-orthogonal projection.

\smallskip\noindent\textbf{The precise condition is K3.}  Uniqueness requires disjoint-support additivity (K3) and positive perturbation cost (K2).  K1 (nonnegativity) and K4 (continuity) are structural but do not contribute to the uniqueness argument.  K3 is both necessary and sufficient: without K3, the spillover cost $\kappa(L_F)$ need not add to $\kappa(E_d)$, and non-orthogonal idempotents may achieve equal or lower cost.

\smallskip\noindent\textbf{Classification of examples.}
\begin{enumerate}[nosep,label=(\roman*),leftmargin=2em]
\item $\kappa_{\mathrm{sq}}(p) = \|p - \mathrm{id}\|_F^2$ (squared Frobenius norm): satisfies K1--K4.  K3 holds because $\|A \oplus B\|_F^2 = \|A\|_F^2 + \|B\|_F^2$.  \textbf{Minimizer is the $B$-orthogonal projection.  Full derivation chain proceeds.}  This is the $\kappa$ used in the worked models (\S\ref{sec:worked}, \S\ref{sec:qubit_model}).
\item $\kappa_{\max}(p) = \|p - \mathrm{id}\|_{\mathrm{op}}$ (operator norm): satisfies K1--K2, K4.  \textbf{Fails K3} (subadditive: $\|A \oplus B\|_{\mathrm{op}} = \max(\|A\|_{\mathrm{op}}, \|B\|_{\mathrm{op}}) \le \|A\|_{\mathrm{op}} + \|B\|_{\mathrm{op}}$).  The spillover argument fails: $\kappa_{\max}(F)$ may equal $\kappa_{\max}(E_d)$ even with $L_F \ne 0$.  \textit{Explicit $2 \times 2$ counterexample:} take $V = \mathbb{R}^2$, $M_d = \spn\{e_1\}$, $\Pi = \spn\{e_2\}$.  The orthogonal projection is $E_d = \diag(1,0)$ with $\kappa_{\max}(E_d) = \|E_d - I\|_{\mathrm{op}} = 1$.  The non-orthogonal idempotent $F = \bigl(\begin{smallmatrix} 1 & \alpha \\ 0 & 0 \end{smallmatrix}\bigr)$ (spillover $\alpha \ne 0$) has $\kappa_{\max}(F) = \|F - I\|_{\mathrm{op}} = \|\bigl(\begin{smallmatrix} 0 & \alpha \\ 0 & -1 \end{smallmatrix}\bigr)\|_{\mathrm{op}}$.  For $|\alpha| \le 1$: $\kappa_{\max}(F) = 1 = \kappa_{\max}(E_d)$; the non-orthogonal $F$ is equally cheap.  Uniqueness fails.  The minimizer is not necessarily orthogonal.  \textbf{Derivation chain does not proceed} --- this is the correct outcome, since $\kappa_{\max}$ models a budget semantics that is not exact (completeness S4 fails).
\item $\kappa_{\mathrm{Fro}}(p) = \|p - \mathrm{id}\|_F$ (Frobenius norm, unsquared): satisfies K1--K2, K4.  \textbf{Fails K3} (subadditive by the triangle inequality).  Same conclusion as (ii).
\item $\kappa_{\mathrm{gap}}(d) = \Delta E(d)$ (energy gap): satisfies K1--K3 when gaps are operationally resolvable and additive on disjoint subsystems (above the error-correction threshold).  \textbf{Minimizer is the $B$-orthogonal projection in the qualifying regime.}
\item $\kappa_{\gamma}(d) = \gamma(d)$ (syndrome extraction overhead): satisfies K1--K3 above threshold, as verified in the surface-code example.  \textbf{Same conclusion as (iv).}
\end{enumerate}

\smallskip\noindent\textbf{Relation to $T_{\mathrm{adj}}$.}  The structural theorem ($T_{\mathrm{adj}}$) does not use this proposition.  The uniqueness result here gives the physical interpretation of why the orthogonal projector is selected; the structural proof of self-adjointness uses only block orthogonality.  This proposition is classificatory, not load-bearing.

\smallskip\noindent\textbf{Approximate K3} (cross-talk $\delta$): if $|\kappa(x \oplus y) - \kappa(x) - \kappa(y)| \le \delta$, then $E_d$ is an approximate orthogonal projection ($\|E_d - E_d^\dagger\| = O(\delta)$), and the discrete algebraic structure is exactly preserved for $\delta < \delta_{\mathrm{crit}}$ (Proposition~\ref{prop:robustness}).
\end{proposition}

\begin{corollary}[Agreement of Structural and Physical Projections]\label{cor:Tadj_agree}
The $B$-orthogonal projection ($T_{\mathrm{adj}}$, structural) and the cost-minimizing idempotent (Proposition~\ref{prop:kappa_class}, physical) agree: both yield $E_d = \pi_{M_d}^{(B)}$.  The structural proof uses only block orthogonality; the physical proof uses K3 + K2.  Neither depends on the other; they independently select the same operator.

\smallskip\noindent\textit{When uniqueness may fail.}  If K3 holds only approximately ($|\kappa(x \oplus y) - \kappa(x) - \kappa(y)| \le \delta$), the cost-minimizing idempotent is $O(\delta)$-close to the $B$-orthogonal projector but may not be exactly self-adjoint.  The discrete algebraic structure (Wedderburn block decomposition, field selection) is exactly preserved for $\delta < \delta_{\mathrm{crit}}$; continuous quantities (bilinear form, probabilities) degrade as $O(\delta)$.
\end{corollary}

% ═══════════════════════════════════════════════════════════════
\section{The Bridge: From Order-Dependence to Noncommutativity}
\label{sec:bridge}

\noindent\textit{Bridge chain at a glance.}  The arena (Section~\ref{sec:arena}) delivers a finite-dimensional real vector space with sector decomposition, block-orthogonal inner product, and self-adjoint sector projections generating a commutative diagonal algebra.  The bridge shows this diagonal algebra is \emph{not} the full enforcement algebra:
\[
\underbrace{L_\Delta}_{\Delta > 0} \;\to\; \underbrace{L_{\mathrm{blk}}}_{\text{diagonal} \Rightarrow \text{additive cost: contradiction}} \;\to\; \underbrace{L_\Pi}_{\text{codespace rotation}} \;\to\; \underbrace{T_{\mathrm{alg}}}_{[E_{d_1},\, F_\Pi] \ne 0}
\]
Every step uses only pre-algebraic or sector-level structure; the full noncommutative $*$-algebra, GNS representation, and field selection are derived in the skeleton (Section~\ref{sec:skeleton}).  A complete numerical instantiation of every step from $\kappa$ through $T_{\mathrm{Born}}$ is given in the worked examples (Sections~\ref{sec:worked} and~\ref{sec:qubit_model}).

\begin{tcolorbox}[apfbox, title={How the pre-algebraic structure arises from A1+MD+BW (no forward references)}]
The vector space, inner product, and adjoint are constructed in the following order.  Each step uses only results already established; no later result is presupposed.

\smallskip
\noindent\textbf{Step 1: Real numbers.}  A1 assigns $C(\Gamma) \in \mathbb{R}_{>0}$ and $\eps(d) \in \mathbb{R}_{>0}$.  The real number field is present from A1's statement, not imported later.

\smallskip
\noindent\textbf{Step 2: Convex state space.}  OR1 (Theorem~\ref{thm:OR1}): the budget constraint $\sum \eps(d) \le C$ is linear in mixture weights, so $\Omega(\Gamma)$ is convex.  \textit{Input: A1 only.}

\smallskip
\noindent\textbf{Step 3: Finite dimensionality.}  $L_{\eps^*}$ (from MD) gives a uniform cost floor $\eps^* > 0$.  $T_{\mathrm{form}}$ (Theorem~\ref{thm:Tform}): $\dim(S_\Gamma) \le C/\mu^* < \infty$.  \textit{Input: A1 + MD.}

\smallskip
\noindent\textbf{Step 4: Finite-dimensional vector space $V$.}  $T_{\mathrm{embed}}$ (Theorem~\ref{thm:Tembed}): map each state $\sigma$ to its cost-coordinate vector $\varphi(\sigma) \in \mathbb{R}^{N+1}$.  Define $V := \spn_\mathbb{R}(\conv(\varphi(\Omega)))$.  The $\mathcal{D}$-quotient identifies $S_\Gamma$ with $V$ (Remark~\ref{rem:linearity}).  \textit{Input: A1 + MD + $\mathcal{D}$-quotient.  Output: a finite-dimensional real vector space $V$ with $\dim(V) \le N+1$.}

\smallskip
\noindent\textbf{Step 5: Sector decomposition.}  $T_{\mathrm{sep}}$ (Theorem~\ref{thm:Tsep}): disjoint anchor sets give a pre-metric algebraic direct sum $S_\Gamma = \bigoplus_d M_d \oplus \Pi$.  No inner product used.  \textit{Input: A1 + FD3 + K3.}

\smallskip
\noindent\textbf{Step 6: Sector-orthogonal bilinear form $B$.}  $L_\omega$ (Lemma~\ref{lem:Lomega}): K3 forces inter-sector orthogonality of any admissible bilinear form.  Choose positive-definite $B_d$ on each sector; take the direct sum.  The inter-sector orthogonality is forced; the within-sector forms are free (normalization convention).  \textit{Input: $T_{\mathrm{sep}}$ + K3.  Output: a positive-definite bilinear form $B$ on $V$ with $B(u_d, v_{d'}) = 0$ for $d \ne d'$.}

\smallskip
\noindent\textbf{Step 7: Adjoint.}  Define the $B$-adjoint: $B(Au, v) = B(u, A^*v)$ for all $u, v \in V$.  This is not presupposed; it is the adjoint induced by $B$ on $\End(V)$.  The $*$-involution exists because $B$ is a non-degenerate bilinear form on a finite-dimensional space.  \textit{Input: $B$ from Step~6.}

\smallskip
\noindent\textbf{Step 8: Self-adjoint enforcement projections.}  $T_{\mathrm{adj}}$ (Theorem~\ref{thm:Tadj_sector}): by $T_{\mathrm{sector}}$(iv), $M_d \perp \ker(E_d)$ with respect to $B$, so $E_d$ is the $B$-orthogonal projection onto $M_d$.  Self-adjointness $E_d = E_d^*$ follows.  \textbf{This result is structurally invariant across the admissible $\kappa$-class}: any cost function satisfying K1--K3 produces the same projection structure (Proposition~\ref{prop:kappa_class}).  \textit{Input: $T_{\mathrm{sector}}$ + K3.  Output: $E_d \in \End(V)$, $E_d = E_d^*$, $E_d^2 = E_d$.}

\smallskip
\noindent At this point --- and only at this point --- the algebra $\mathcal{A} := \mathrm{Alg}_\mathbb{R}(\{E_d\} \cup \{E_{\{d_1,d_2\}}\}) \subseteq \End(V)$ is generated (O3), and noncommutativity is claimed ($T_{\mathrm{alg}}$).  The Hilbert space is not constructed until Section~\ref{sec:skeleton} (GNS, $T_{\mathrm{GNS}}$).
\end{tcolorbox}

\subsection{Superadditivity ($L_\Delta$)}

\begin{remark}[Joint enforcement cost: why sharing an interface costs more]\label{rem:joint_cost}
At separate interfaces $\Gamma_1, \Gamma_2$ with $S_{\Gamma_1} = M_{d_1}$, $S_{\Gamma_2} = M_{d_2}$: each interface has its own budget, and the total enforcement cost is $\eps(d_1) + \eps(d_2)$ (additive by $L_{\mathrm{loc}}$).  There is no pool.

At a shared interface $\Gamma$ with $S_\Gamma = M_{d_1} \oplus M_{d_2} \oplus \Pi$ and $\Pi \ne \{0\}$: the enforcement system is responsible for the \emph{entire} substrate $S_\Gamma$ under a single budget $C(\Gamma)$.  The pool $\Pi$ exists at $\Gamma$ whether or not any named distinction is anchored there.  The joint enforcement cost $\eps(\{d_1, d_2\})$ is the minimum capacity committed at $\Gamma$ such that an admissible state containing $d_1$ and $d_2$ is maintained against all perturbations at $\Gamma$ --- including those supported in $\Pi$.

The surplus $\Delta := \eps(\{d_1,d_2\}) - \eps(d_1) - \eps(d_2)$ measures the cost of sharing an interface: the capacity overhead of co-location beyond what independent enforcement at separate interfaces would require.  This overhead is zero when $\Pi = \{0\}$ (the classical limit) and strictly positive when $\Pi \ne \{0\}$.
\end{remark}

\noindent\textit{Summary.}  $\Pi = \{0\}$ is the classical boundary; $\Pi \ne \{0\}$ is the source of pool-supported joint-defense overhead and ultimately of noncommutativity.  The pool is not an additional assumption; it is a consequence of the sector decomposition at any sufficiently rich interface.  See Proposition~\ref{prop:pool} (Section~\ref{sec:arena}) for the complete structural characterization.

\begin{theorem}[$L_\Delta$: Superadditivity of Co-Located Enforcement]\label{thm:LDelta}
For co-located distinctions $d_1, d_2$ at a single interface $\Gamma$ with $M_{d_1} \cap M_{d_2} = \{0\}$ and a non-empty shared pool $\Pi \ne \{0\}$:
\[
\Delta(d_1, d_2) \;:=\; \eps(\{d_1, d_2\}) - \eps(d_1) - \eps(d_2) \;>\; 0.
\]

\noindent\textit{Definition: joint perturbation class} (independent of $\Pi$).  $\mathcal{P}(\{d_1, d_2\}) := \{p \in \End(S_\Gamma) : p$ can cause the simultaneous enforcement of $d_1$ and $d_2$ at $\Gamma$ to fail --- either by directly destroying $d_1$ or $d_2$, or by rendering the enforcement state at $\Gamma$ inadmissible (total committed cost $> C$)$\}$.  This class is defined by its effect on the pair, not by support restrictions.  The co-located joint enforcement cost is $\eps(\{d_1, d_2\}) := \inf\{\text{total capacity committed to resist all } p \in \mathcal{P}(\{d_1,d_2\})\}$.

\noindent\textit{Co-location conditions} (CL1--CL3, all derived).  CL1: $M_{d_1} \cap M_{d_2} = \{0\}$ (from $T_{\mathrm{sep}}$).  CL2: both at same $\Gamma$.  CL3: $\Pi \ne \{0\}$, equivalently $S_\Gamma \ne \bigoplus_{d \in \mathcal{B}} M_d$ (Proposition~\ref{prop:pool}(iii)--(iv)).  This is generic at sufficiently rich interfaces but not automatic from $\dim(S_\Gamma) \ge 3$ alone: a three-sector classical model ($S_\Gamma = M_{d_1} \oplus M_{d_2} \oplus M_{d_3}$, each one-dimensional) has $\dim = 3$ and $\Pi = \{0\}$.

\smallskip\noindent\textit{Proof} (five steps; pre-algebraic).

\textit{Step 1: the joint perturbation class is strictly larger than the union of individual classes.}
Define $\mathcal{P}(d_i)$ as the perturbations that can destroy $d_i$ individually.  Clearly $\mathcal{P}(\{d_1,d_2\}) \supseteq \mathcal{P}(d_1) \cup \mathcal{P}(d_2)$ (destroying either distinction causes joint enforcement to fail).  By Proposition~\ref{prop:pool}(iii), $\Pi \ne \{0\}$ places the interface in the non-classical regime: perturbations with $\mathrm{supp}(p) \subseteq \Pi$ are in $\mathcal{P}(\{d_1, d_2\})$ but not in $\mathcal{P}(d_1) \cup \mathcal{P}(d_2)$.  \textit{Why:} a $\Pi$-supported perturbation does not directly destroy $d_1$ or $d_2$ (their anchors are disjoint from $\Pi$ by $T_{\mathrm{sep}}$), but it changes physically meaningful degrees of freedom at $\Gamma$ (by SP: $\Pi$ has no null directions).  A1's budget covers the entire interface $\Gamma$; an unresisted change to physically meaningful DOF at $\Gamma$ can increase the total enforcement cost at $\Gamma$ beyond $C$, rendering the state inadmissible.

\textit{Step 2: $\Pi$-supported perturbations cannot be neutralized by defenses confined to $M_{d_1} \oplus M_{d_2}$.}
A defense $D$ with $\mathrm{supp}(D) \subseteq M_{d_1} \oplus M_{d_2}$ acts as the identity on $\Pi$ (by definition of support: $D|_{\Pi} = \mathrm{id}$).  It therefore has no effect on $\Pi$-supported perturbations.  To resist a perturbation $p_v$ with $\mathrm{supp}(p_v) \subseteq \Pi$ and $\kappa(p_v) > 0$ (K2), the defense must have nonzero support in $\Pi$.

\textit{Step 3: the $\Pi$-defense has strictly positive cost.}
By Step~2, any admissible joint defense must commit positive capacity in $\Pi$.  Let $\kappa_\Pi$ denote the minimum such commitment.  By K2, any perturbation $p_v$ along a non-null $\Pi$ direction has $\kappa(p_v) > 0$; resisting it requires committing positive capacity (otherwise $p_v$ proceeds unopposed at the defended state, altering physically meaningful DOF at positive cost).  Therefore $\kappa_\Pi > 0$.

\textit{Step 4: the $\Pi$-defense is supported disjointly from the sector defenses.}
By $T_{\mathrm{sep}}$, the three regions $M_{d_1}$, $M_{d_2}$, $\Pi$ are pairwise disjoint subspaces.  The sector defenses (cost $\eps(d_1)$ in $M_{d_1}$, cost $\eps(d_2)$ in $M_{d_2}$) and the pool defense (cost $\kappa_\Pi$ in $\Pi$) have pairwise disjoint supports.

\textit{Step 5: K3 gives exact additivity, hence $\Delta > 0$.}
By K3 (Theorem~\ref{thm:K3}), costs on disjoint supports add exactly:
\[
\eps(\{d_1, d_2\}) \;=\; \eps(d_1) + \eps(d_2) + \kappa_\Pi, \qquad \kappa_\Pi > 0.
\]
Therefore $\Delta(d_1, d_2) = \kappa_\Pi > 0$.  \qed

\noindent\textit{The mechanism in one sentence:} at a co-located interface, the joint perturbation class includes $\Pi$-supported threats absent at separate interfaces; defending against them requires positive $\Pi$-committed capacity, which adds to the sector costs by K3.

\noindent\textit{Regularity conditions used:} A1 (budget), K2 (positive perturbation cost, from A1), K3 (disjoint-support additivity, Theorem~\ref{thm:K3}), $T_{\mathrm{sep}}$ (sector decomposition, from A1), SP (no null directions in the physical quotient, constitutive).  \textbf{No algebraic structure, faithful state, or GNS machinery is used.}
\end{theorem}

\subsection{Order-dependence (T1)}

\begin{theorem}[T1: Order-Dependent Enforcement]\label{thm:T1}
There exist distinctions $d_1, d_2, d_3$ and a state $\sigma_0$ such that:
\[
\mathfrak{E}_{d_3} \circ \mathfrak{E}_{d_1}(\sigma_0) \;\ne\; \mathfrak{E}_{d_3} \circ \mathfrak{E}_{d_2}(\sigma_0).
\]

\noindent\textit{Proof.}  By NT (non-degeneracy of the cost spectrum, from MD's floor alone), $\exists\, d_1, d_2$ with $\eps(d_1) \ne \eps(d_2)$; WLOG $\eps(d_1) < \eps(d_2)$.  By \fbox{BW}, $\exists\, d_3$ with $C - \eps(d_1) \ge \eps(d_3) > C - \eps(d_2)$.  Starting from $\sigma_0 = (\emptyset, C)$:
\begin{align*}
\mathfrak{E}_{d_1}(\sigma_0) &= (\{d_1\}, C - \eps(d_1)), \\
\mathfrak{E}_{d_3}(\mathfrak{E}_{d_1}(\sigma_0)) &= (\{d_1, d_3\}, C - \eps(d_1) - \eps(d_3)) \;\in\; \Omega \quad (\text{budget $\ge 0$}), \\
\mathfrak{E}_{d_2}(\sigma_0) &= (\{d_2\}, C - \eps(d_2)), \\
\mathfrak{E}_{d_3}(\mathfrak{E}_{d_2}(\sigma_0)) &= \bot \quad (\text{budget } C - \eps(d_2) - \eps(d_3) < 0).
\end{align*}
The first composite is an admissible state; the second is the failure symbol $\bot$.  These are distinct.  \qed

\noindent\textit{Regularity conditions used:} NT (from MD's floor); \fbox{BW enters here}.

\noindent\textit{Classical countermodel (why T1 alone does not force noncommutativity).}  Bayesian conditioning on events $A$ then $B$ gives a different posterior than $B$ then $A$: this is order-dependent, but classical.  A classical probability simplex satisfies A1 with enforcement operations $E_d$ as face projections (diagonal matrices that zero out coordinates outside the face).  These commute: $[E_d, E_{d'}] = 0$ for all $d, d'$.  The bridge to noncommutativity requires $L_\Delta$ ($\Delta > 0$), which requires the co-location conditions CL1--CL3 to hold simultaneously.  In a classical simplex, distinctions sharing an interface have overlapping mechanisms ($M_{d_1} \cap M_{d_2} \ne \{0\}$): CL2 fails, the cross-talk coupling satisfies $\kappa \ge 1/2$, and $\Delta \le 0$.  The classical regime is precisely the regime where CL1--CL3 do not all hold.  The passage from order-dependence (T1) to noncommutativity ($T_{\mathrm{alg}}$) therefore requires the bridge chain $L_\Delta \to L_{\mathrm{blk}} \to L_\Pi$, which is operative only when all three co-location conditions are satisfied --- a regime that excludes classical models.
\end{theorem}

\begin{tcolorbox}[apfbox, title={Two Distinct Routes to Noncommutativity}]
The derivation chain reaches $T_{\mathrm{alg}}$ (noncommutativity) via two independent witness routes.  Both produce the same output ($[E_{d_1}, F_\Pi] \ne 0$), but they use different inputs and have different scopes.

\smallskip\noindent\textbf{Route A: General budget-window witness.}  NT (non-degeneracy of cost spectrum, from MD's floor) + BW (budget-window triple) $\to$ T1 (order-dependent enforcement).  T1 witnesses distinct enforcement outcomes under different orderings, establishing that the enforcement algebra cannot be represented by commuting operations.  BW is the richness premise that guarantees a witness triple exists at sufficiently complex interfaces ($n_{\max} \ge 3$).  \textit{Scope:} general interfaces with three or more distinct cost levels.

\smallskip\noindent\textbf{Route B: Direct pool-rotation witness.}  $M_{d_1} \cap M_{d_2} = \{0\}$, $\Pi \ne \{0\}$ (disjoint anchors with nonzero pool) $\to$ $L_\Delta$ ($\Delta > 0$) $\to$ $L_{\mathrm{blk}}$ (block-diagonal defense fails) $\to$ $L_\Pi$ (codespace rotation through $\Pi$) $\to$ $T_{\mathrm{alg}}$ ($[E_{d_1}, F_\Pi] \ne 0$).  This route does \emph{not} invoke BW or T1; it uses only the co-location conditions CL1--CL3 and the pool structure.  \textit{Scope:} any interface with two disjoint-anchor distinctions and a nonzero pool ($\dim(S_\Gamma) \ge 3$).

\smallskip\noindent\textbf{Relationship.}  Route A provides a general existence guarantee for order-dependent enforcement at rich interfaces; Route B provides the structural mechanism (pool rotation) that generates noncommutativity at any interface with $\Pi \ne \{0\}$.  The worked examples in Sections~\ref{sec:worked} and~\ref{sec:qubit_model} use Route~B exclusively: they construct $\Delta > 0$ and the rotated codespace directly, without invoking BW.  BW is sufficient but not necessary for noncommutativity at interfaces with $\dim(S_\Gamma) \ge 3$.
\end{tcolorbox}

\subsection{Sector-orthogonal bilinear form ($L_\omega$)}

\noindent The inter-sector orthogonality pattern of the bilinear form $B$ is forced by K3.  The within-sector forms are free (a normalization gauge).  All downstream results --- projections, algebra, GNS, field selection --- depend only on the forced pattern, not on the gauge choice.

\begin{lemma}[$L_\omega$: Sector-Orthogonal Bilinear Form]\label{lem:Lomega}
Given the sector decomposition $S_\Gamma = \bigoplus_d M_d \oplus \Pi$ from $T_{\mathrm{sep}}$ and K3 (Theorem~\ref{thm:K3}):
\begin{enumerate}[nosep,label=(\alph*),leftmargin=2em]
\item \textbf{Forced structure:} K3 forces $B(u_d, v_{d'}) = 0$ for $d \ne d'$ and $B(u_d, v_\Pi) = 0$ for all $d$.  This inter-sector orthogonality pattern is derived, not chosen ($L_\omega$.5 below).
\item \textbf{Constructed completion:} a symmetric positive-definite bilinear form $B: V \times V \to \mathbb{R}$ compatible with this pattern is constructed by choosing any positive-definite form $B_d$ on each sector $M_d$ and on $\Pi$, and defining $B$ as their direct sum with zero cross-terms.  The within-sector forms $B_d$ are \emph{not} derived from cost; they are free choices.
\item \textbf{Gauge invariance:} different within-sector choices $B_d \mapsto B_d'$ produce gauge-equivalent forms.  All downstream results --- the projection structure $E_d = \pi_{M_d}^{(B)}$, the algebra $\mathcal{A} \cong \bigoplus_k M_{n_k}(\mathbb{F}_k)$, the GNS Hilbert space $\mathcal{H}_\omega$, and the field selection (T2c) --- are invariant under this gauge.
\end{enumerate}
The construction and proofs follow.

\medskip\noindent
\textbf{$L_\omega$.1: Basis-level definition.}
Choose any basis $\{e_1^{(d)}, \ldots, e_{k_d}^{(d)}\}$ for each sector $M_d$ and any basis $\{f_1, \ldots, f_m\}$ for $\Pi$.  On each sector, choose any positive-definite bilinear form: e.g., the Euclidean form $B_d(e_i^{(d)}, e_j^{(d)}) := \delta_{ij}$.  Similarly $B_\Pi(f_i, f_j) := \delta_{ij}$.  Define $B$ on all of $V$ by:
\[
B(u, v) \;:=\; \sum_{d \in \mathcal{B}} B_d(u_d, v_d) \;+\; B_\Pi(u_\Pi, v_\Pi),
\]
where $u = \sum_d u_d + u_\Pi$ and $v = \sum_d v_d + v_\Pi$ are the $T_{\mathrm{sep}}$ decompositions.  Cross-sector terms are set to zero: $B(u_d, v_{d'}) := 0$ for $d \ne d'$, and $B(u_d, v_\Pi) := 0$ for all $d$.

\medskip\noindent
\textbf{$L_\omega$.2: Well-definedness.}
The definition depends on the $T_{\mathrm{sep}}$ decomposition but \emph{not} on the choice of basis within each sector: if $\{e'\}$ is another basis for $M_d$, the change of basis within $M_d$ produces a different $B_d$ (a gauge choice), but the inter-sector orthogonality pattern is unchanged.  All downstream results depend only on this pattern, not on the within-sector normalization.

\medskip\noindent
\textbf{$L_\omega$.3: Symmetry.}
$B(u, v) = B(v, u)$ for all $u, v \in V$.  \textit{Proof:} By the sector decomposition, $B(u, v) = \sum_d B_d(u_d, v_d) + B_\Pi(u_\Pi, v_\Pi)$.  Each $B_d$ is symmetric (the Euclidean form on each sector is symmetric by construction; any positive-definite bilinear form on a finite-dimensional real space can be chosen symmetric without loss, since the symmetric part $(B_d + B_d^\top)/2$ is still positive-definite).  Therefore $B$ is symmetric.

\medskip\noindent
\textbf{$L_\omega$.4: Positive-definiteness.}
$B(v, v) > 0$ for all nonzero $v \in V$.  \textit{Proof:} Decompose $v = \sum_d v_d + v_\Pi$.  Since $v \ne 0$, at least one component is nonzero.  All terms $B_d(v_d, v_d) \ge 0$ and $B_\Pi(v_\Pi, v_\Pi) \ge 0$ (positive-semidefiniteness of each block form).  If $v_d \ne 0$ for some $d$, then $B_d(v_d, v_d) > 0$ (positive-definiteness of $B_d$).  If only $v_\Pi \ne 0$, then $B_\Pi(v_\Pi, v_\Pi) > 0$.  In either case $B(v,v) > 0$.  In particular, $B$ is nondegenerate: $B(u,v) = 0$ for all $v$ implies $u = 0$ (by taking $v = u$).

\medskip\noindent
\textbf{$L_\omega$.5: K3 forces inter-sector orthogonality.}
The vanishing of all cross-sector terms is not a definition but a \emph{consequence} of K3.  \textit{Proof:} A bilinear form compatible with K3's disjoint-support additivity must satisfy: the norm of a configuration with components in sectors $M_d$ and $M_{d'}$ depends only on what happens within each sector, not on their relative alignment.  Formally: K3 requires $\kappa(p_d \oplus p_{d'}) = \kappa(p_d) + \kappa(p_{d'})$ for disjoint-support perturbations.  If $B$ had a nonzero cross-term $B(u_d, v_{d'}) \ne 0$, then $B(u_d + v_{d'}, u_d + v_{d'}) = B(u_d, u_d) + B(v_{d'}, v_{d'}) + 2B(u_d, v_{d'}) \ne B(u_d, u_d) + B(v_{d'}, v_{d'})$: the cost of the joint configuration would depend on the cross-sector alignment, violating K3's independence.  Therefore $B(u_d, v_{d'}) = 0$ for $d \ne d'$.  The same argument gives $B(u_d, v_\Pi) = 0$ for all $d$.

An equivalent one-line derivation (from $T_{\mathrm{sector}}$(iv)): for $u_d \in M_d$ and $u_{d'} \in M_{d'}$ with $d \ne d'$, K3 gives $\|u_d + u_{d'}\|_B^2 = \|u_d\|_B^2 + \|u_{d'}\|_B^2$, hence $2\,B(u_d, u_{d'}) = 0$.

\medskip\noindent
\textbf{Adjoint compatibility.}
The positive-definite symmetric bilinear form $B$ defines a $B$-adjoint on $\End(V)$: for any $T \in \End(V)$, $T^\dagger$ is the unique map satisfying $B(Tu, v) = B(u, T^\dagger v)$ for all $u, v \in V$.  Existence and uniqueness: standard for nondegenerate bilinear forms on finite-dimensional spaces (Riesz representation).  The $*$-involution on the enforcement algebra $\mathcal{A}$ is then $A^* := A^\dagger$, the $B$-adjoint.  This satisfies involutivity ($(A^*)^* = A$) and anti-multiplicativity ($(AB)^* = B^* A^*$).  For the sector projections: $E_d^* = E_d$ ($T_{\mathrm{adj}}$, which is now a one-line corollary of $L_\omega$.5 + sector orthogonality).

\medskip\noindent
\textbf{Uniqueness.}
The inter-sector orthogonality pattern is uniquely forced.  The within-sector forms $B_d$ are free (a normalization gauge).  Two forms in the same equivalence class produce the same projection structure, the same $*$-isomorphism class of $\mathcal{A}$, and the same downstream results (T2a, T2c, $T_{\mathrm{Born}}$).  \qed
\end{lemma}

\begin{tcolorbox}[apfbox, title={$L_\omega$ Summary: Forced vs.\ Gauge vs.\ Invariant}]
\textbf{Forced by K3:} inter-sector orthogonality ($B(u_d, v_{d'}) = 0$ for $d \ne d'$; $B(u_d, v_\Pi) = 0$).

\smallskip\noindent\textbf{Chosen up to gauge:} within-sector positive-definite geometry ($B_d$ on each $M_d$; $B_\Pi$ on $\Pi$).

\smallskip\noindent\textbf{Invariant downstream (independent of gauge choice):} the projection structure $E_d = \pi_{M_d}^{(B)}$; the $*$-isomorphism class of the algebra $\mathcal{A} \cong \bigoplus_k M_{n_k}(\mathbb{F}_k)$; the Wedderburn block sizes $\{n_k\}$ and block count; the GNS Hilbert space $\mathcal{H}_\omega$ up to unitary equivalence; the field selection (T2c); the Born-rule representation theorem $p(E) = \tr(\rho\,E)$ (the form is invariant; numerical values of $p(E)$ depend on the chosen faithful state $\omega$, which is fixed up to the gauge described in O4).
\end{tcolorbox}

\smallskip\noindent
\textit{Corollary: $N_d = M_d^{\perp_B}$.}  $L_\omega$.5 gives: for any $x_N \in N_d = \bigoplus_{d' \ne d} M_{d'} \oplus \Pi$ and any $m \in M_d$, $B(x_N, m) = 0$.  Conversely, if $x \in M_d^{\perp_B}$, decompose $x = x_d + x_N$; then $0 = B(x, m) = B_d(x_d, m)$ for all $m \in M_d$, giving $x_d = 0$.  Therefore $N_d = M_d^{\perp_B}$, confirming $T_{\mathrm{sector}}$(iv) from the explicit construction.

\smallskip\noindent
\textit{Worked verification ($\mathbb{R}^3$ model).}  $V = \mathbb{R}^3$, $M_{d_1} = \spn\{e_1\}$, $M_{d_2} = \spn\{e_2\}$, $\Pi = \spn\{e_3\}$.  Choose $B_d = B_\Pi = $ Euclidean form.  Then $B = I_3$: $B(e_i, e_j) = \delta_{ij}$.  Inter-sector orthogonality: $B(e_1, e_2) = B(e_1, e_3) = B(e_2, e_3) = 0$.  Positive-definiteness: $B(v,v) = \|v\|^2 > 0$ for $v \ne 0$.  Adjoint: $T^\dagger = T^\top$ (matrix transpose).  Self-adjointness of $E_{d_1} = \diag(1,0,0)$: $(E_{d_1})^\top = E_{d_1}$.  This instantiates all five $L_\omega$ properties concretely.

\smallskip\noindent
\textit{Cauchy--Schwarz inequality.}  Since $B$ is symmetric and positive-definite ($L_\omega$.3 + $L_\omega$.4), the Cauchy--Schwarz inequality holds:
\[
|B(u, v)|^2 \;\le\; B(u, u) \cdot B(v, v) \qquad \text{for all } u, v \in V,
\]
with equality iff $u$ and $v$ are $B$-proportional.  This is a standard consequence of positive-definiteness on finite-dimensional real inner-product spaces (expand $B(u + tv, u + tv) \ge 0$ for all $t \in \mathbb{R}$; the discriminant of the resulting quadratic in $t$ must be $\le 0$).  The inequality is used implicitly throughout: it underlies the well-definedness of $B$-angles between subspaces, the projection inequalities in $T_{\mathrm{adj}}$, and the norm bounds in the robustness proposition (Proposition~\ref{prop:robustness}).

\smallskip\noindent
\textit{Dependence on $\kappa$.}  $B$ depends on $\kappa$ only through the within-sector forms $B_d$.  Specifically:
\begin{itemize}[nosep,leftmargin=2em,itemsep=2pt]
\item \textit{What $\kappa$ determines:} $\kappa$ determines the sector costs $\eps(d)$ (via FD4: $\eps(d) = \inf \kappa$), which determine the within-sector normalization of $B_d$.  Different $\kappa$ satisfying K1--K3 may give different $\eps(d)$ values and hence different $B_d$.
\item \textit{What $\kappa$ does not determine:} the \emph{inter-sector orthogonality pattern} $B(u_d, v_{d'}) = 0$ is forced by K3 alone ($L_\omega$.5), independent of the specific $\kappa$.  The \emph{projection structure} $E_d = \pi_{M_d}^{(B)}$ depends only on the orthogonality pattern, not on the within-sector normalization.  The \emph{algebra structure} $\mathcal{A} \cong \bigoplus_k M_{n_k}(\mathbb{F}_k)$ and all downstream results (T2a, T2c, $T_{\mathrm{Born}}$) depend only on the projection structure, not on $B_d$.
\item \textit{Summary:} $B$ is unique up to within-sector gauge.  The gauge freedom is $B_d \mapsto B_d'$ (any positive-definite form on $M_d$).  All structural results are gauge-invariant.  The choice of $\kappa$ fixes the gauge; the structural conclusions do not depend on the gauge choice.  See Proposition~\ref{prop:kappa_class} for the explicit classification.
\end{itemize}

\subsection{Failure of direct-sum realizability ($L_{\mathrm{blk}}$)}

\begin{theorem}[$L_{\mathrm{blk}}$: Excess Cost $\Leftrightarrow$ Failure of Block-Diagonal Defense]\label{thm:Lblk}
Let $S_\Gamma = M_{d_1} \oplus M_{d_2} \oplus \Pi$ (from $T_{\mathrm{sep}}$, which is pre-metric: no inner product used).  The following are equivalent:
\begin{enumerate}[nosep,label=(\roman*)]
\item $\Delta(d_1, d_2) = 0$;
\item there exists an idempotent joint defender $P \in \End(V)$ (where $\End(V)$ is available from $T_{\mathrm{embed}}$, Step~4 of the construction box) with $\im(P) \subseteq M_{d_1} \oplus M_{d_2}$, $P|_{M_{d_i}} = \mathrm{id}$ for $i = 1,2$, and $P|_\Pi = 0$;
\item no positive-cost distinction with anchor in $\Pi$ needs to be defended for joint enforcement of $\{d_1, d_2\}$.
\end{enumerate}

\smallskip\noindent\textit{Note: what is used and where it comes from.}  This proof uses only: the $T_{\mathrm{sep}}$ direct sum (pre-metric, Section~3), the definition of $\End(V)$ (from $T_{\mathrm{embed}}$), idempotency ($E_d^2 = E_d$, from FD5a), and the cost results $L_\Delta$ and SP (non-null directions in the quotient).  Self-adjointness ($E_d = E_d^*$, from $T_{\mathrm{adj}}$) and $B$-orthogonality (from $L_\omega$) are \emph{not required} for the biconditional.  They enter only in the corollary below, where they are cited at point of use.

\smallskip\noindent\textit{Proof.}

\textit{Key bridge: diagonal algebras force additive cost.}  In $\mathcal{A}_{\mathrm{diag}} = \spn\{E_d\}$ ($T_{\mathrm{sector}}$), the sector projections commute ($E_{d_1} E_{d_2} = 0$ for $d_1 \ne d_2$) and no off-diagonal coupling term exists.  Any joint enforcement map expressible in $\mathcal{A}_{\mathrm{diag}}$ therefore has cost $\omega(E_{d_1} + E_{d_2}) = \omega(E_{d_1}) + \omega(E_{d_2})$: exact additivity.  Strict superadditivity ($\Delta > 0$ from $L_\Delta$) is incompatible with a diagonal joint defender.  This is the cost contradiction that drives the entire bridge.

\textit{(i)$\Leftrightarrow$(iii).}  If $\Delta > 0$, then by $L_\Delta$ the joint perturbation class includes $\Pi$-supported threats requiring positive defense cost in $\Pi$: (iii) fails.  Conversely, if no positive-cost distinction has its anchor in $\Pi$, then the joint defense of $\{d_1,d_2\}$ requires no $\Pi$-defense; the cost is $\eps(d_1) + \eps(d_2)$ (additive by $T_{\mathrm{sep}}$) and $\Delta = 0$: (i) holds.

\textit{(ii)$\Rightarrow$(iii) (contrapositive).}  Suppose a positive-cost distinction $d_\Pi$ exists with $M_{d_\Pi} \subseteq \Pi$ and $\eps(d_\Pi) > 0$.  Let $P$ be any idempotent with $\im(P) \subseteq M_{d_1} \oplus M_{d_2}$ and $P|_{M_{d_i}} = \mathrm{id}$.  Since $\im(P) \subseteq M_{d_1} \oplus M_{d_2}$ and $M_{d_1} \oplus M_{d_2}$ is a $T_{\mathrm{sep}}$ direct-sum complement of $\Pi$ (pre-metric: no $B$ needed), $P$ maps $\Pi$ into $M_{d_1} \oplus M_{d_2}$.  But $P^2 = P$ and $P|_{M_{d_i}} = \mathrm{id}$ force $P|_\Pi = 0$ (since $P(\Pi) \subseteq M_{d_1} \oplus M_{d_2}$, applying $P$ again: $P(Pv) = Pv$ for $v \in \Pi$, so $Pv \in M_{d_1} \oplus M_{d_2}$ with $P(Pv) = Pv$; but $P|_{M_{d_1} \oplus M_{d_2}} = \mathrm{id}$ (since $P|_{M_{d_i}} = \mathrm{id}$), so $Pv = Pv$ is consistent only if $Pv$ is already in $M_{d_1} \oplus M_{d_2}$ and fixed by $P$, giving $P|_\Pi: \Pi \to M_{d_1} \oplus M_{d_2}$; idempotency then forces $Pv = 0$ for $v \in \Pi$ because any nonzero $Pv \in M_{d_i}$ would require $P(Pv) = Pv$ with $Pv \in M_{d_i}$ already fixed, but the image $Pv$ must also satisfy $E_{d_i}(Pv) = Pv$, which is consistent only if $Pv$ is the zero vector or already in $M_{d_i}$; the direct-sum structure forces $Pv = 0$).  Therefore $P|_{M_{d_\Pi}} = 0$: $P$ cannot defend $d_\Pi$.  So no block-diagonal idempotent is an admissible joint defender: (ii) fails.

\textit{(iii)$\Rightarrow$(ii).}  If no positive-cost distinction has its anchor in $\Pi$, define $P_{\mathrm{ds}} := E_{d_1} + E_{d_2}$ (using FD5a: both are idempotents with $\im(E_{d_i}) = M_{d_i}$).  By the $T_{\mathrm{sep}}$ direct sum, $E_{d_1}$ and $E_{d_2}$ have disjoint images, so $E_{d_1} E_{d_2} = 0$ (this uses only the pre-metric direct sum, not $B$-orthogonality).  Therefore $P_{\mathrm{ds}}^2 = E_{d_1}^2 + E_{d_2}^2 = E_{d_1} + E_{d_2} = P_{\mathrm{ds}}$: an idempotent with $\im(P_{\mathrm{ds}}) = M_{d_1} \oplus M_{d_2}$.  It defends $d_1$ and $d_2$ by construction.  No pool defense needed (hypothesis).  So $P_{\mathrm{ds}}$ is admissible: (ii) holds.

\textit{Closing.}  (i)$\Leftrightarrow$(iii) by the first step.  (iii)$\Rightarrow$(ii) by the third.  (ii)$\Rightarrow$(iii) by the second (contrapositive).  Therefore (i)$\Leftrightarrow$(ii)$\Leftrightarrow$(iii).  \qed

\begin{corollary}[Nonzero Commutator Witness]\label{cor:commutator}
If $\Delta(d_1,d_2) > 0$, then for any admissible joint defender $P$:
\begin{enumerate}[nosep,label=(\alph*)]
\item $F_\Pi := P - E_{d_1} - E_{d_2} \ne 0$.
\item There exists $v \in M_{d_1}$ with $Pv = v + \pi_v$, $\pi_v \in \Pi \setminus \{0\}$.
\item $[E_{d_1}, P]\,v = -\pi_v \ne 0$.
\end{enumerate}
\end{corollary}

\noindent\textit{What is used.}  Part (a) uses only $L_{\mathrm{blk}}$ (block-diagonal defense fails).  Part (b) uses $P$'s nontrivial $\Pi$-support plus idempotency (no inner product needed).  Part (c) uses $E_{d_1}$'s kernel: $E_{d_1}(\pi_v) = 0$ because $\pi_v \in \Pi$ and $\ker(E_{d_1}) \supseteq \Pi$.  This follows from $T_{\mathrm{sector}}$(v): $E_{d_1}|_\Pi = 0$ by definition of the sector maps.  This is the \emph{only} point in the bridge where $T_{\mathrm{adj}}$'s output enters.

\smallskip
\textit{Proof of (b).}  $P$ has nontrivial $\Pi$-action ($L_{\mathrm{blk}}$: no block-diagonal defender exists when $\Delta > 0$) and defends $d_1$ ($P|_{M_{d_1}}$ is injective: $Pv = 0$ for $v \in M_{d_1}$ would mean $d_1$ is undefended).  If $Pv \in M_{d_1}$ for all $v \in M_{d_1}$, then $P|_{M_{d_1}} = E_{d_1}|_{M_{d_1}}$, $P|_{M_{d_2}} = E_{d_2}|_{M_{d_2}}$, and $P|_\Pi = 0$ by idempotency --- contradicting nontrivial $\Pi$-action.  So some $v \in M_{d_1}$ has $\pi_v \ne 0$.

\textit{Proof of (c).}  $[E_{d_1}, P]\,v = E_{d_1}(Pv) - P(E_{d_1}v) = E_{d_1}(v + \pi_v) - P(v) = (v + E_{d_1}\pi_v) - (v + \pi_v)$.  Now $E_{d_1}\pi_v = 0$ because $\pi_v \in \Pi \subseteq \ker(E_{d_1})$ (from $T_{\mathrm{sector}}$(v): $E_{d_1}|_\Pi = 0$).  Therefore $[E_{d_1}, P]\,v = -\pi_v \ne 0$.  \qed

\noindent\textit{Regularity conditions used:} $T_{\mathrm{sep}}$ (pre-metric), FD5a (idempotents), SP (from constitutive quotient), $L_\Delta$ (from A1+K2+K3).  $T_{\mathrm{adj}}$/$L_\omega$ enter only at Corollary part~(c).

\noindent\textit{Provenance of projections.}  All projections in $L_{\mathrm{blk}}$ are constructed from prior results, not assumed.  The individual enforcement maps $E_{d_1}, E_{d_2}$ are idempotents from FD5a, upgraded to $B$-orthogonal projections by $T_{\mathrm{adj}}$ (using $L_\omega$'s bilinear form, constructed from K3).  The joint defender $P = \pi_{W_*}$ is the $B$-orthogonal projection onto the minimum-cost codespace $W_*$, whose existence is guaranteed by $L_\Pi$ (compactness of the Grassmannian $\Gr(k, V)$, from OR3's finite dimensionality).  No spectral theory or pre-existing Hilbert space is invoked.
\end{theorem}

\begin{lemma}[$L_\omega$-mono: Cost Monotonicity in Codespace Dimension]\label{lem:omega_mono}
Let $\rho := \sum_d (\eps(d)/C)\,E_d + \delta E_\Pi$ with $\delta > 0$ (the positive-definite density operator from O4), and $\omega(\pi_W) := \tr(\rho\,\pi_W)/\tr(\rho)$.  For any subspaces $W \subsetneq W'$ of $S_\Gamma$:
\[
\omega(\pi_{W'}) \;>\; \omega(\pi_W).
\]

\noindent\textit{Proof.}  Since $W \subsetneq W'$, the difference $\pi_{W'} - \pi_W$ is the orthogonal projection onto $W' \cap W^\perp$, which is a nonzero subspace (of dimension $\dim W' - \dim W \ge 1$).  Therefore $\pi_{W'} - \pi_W$ is a nonzero positive semidefinite operator.  Since $\rho$ is positive-definite (all eigenvalues $> 0$):
\[
\tr\bigl(\rho\,(\pi_{W'} - \pi_W)\bigr) = \sum_i \lambda_i \|(\pi_{W'} - \pi_W)\,e_i\|^2 > 0,
\]
where $\rho = \sum_i \lambda_i |e_i\rangle\langle e_i|$ with $\lambda_i > 0$.  The strict inequality holds because at least one eigenvector $e_i$ has nonzero component in $W' \cap W^\perp$ (otherwise $W' \cap W^\perp = \{0\}$, contradicting $W \subsetneq W'$).  Dividing by $\tr(\rho) > 0$ gives the result.  \qed

\noindent\textit{Regularity conditions used:} O4 ($\rho > 0$), finite dimensionality (OR3).
\end{lemma}

\begin{lemma}[$L_\Pi$: Codespace Rotation]\label{lem:LPi}
Under the hypotheses of $L_{\mathrm{blk}}$ with $\Delta > 0$, the minimum-cost admissible joint defender $\pi_{W_*}$ satisfies: (a)~$W_*$ exists (infimum attained on the compact Grassmannian $\Gr(k, S_\Gamma)$, from OR3); (b)~$W_* \ne M_{d_1} \oplus M_{d_2}$ (by $L_{\mathrm{blk}}$); (c)~$\dim(W_*) = \dim(M_{d_1} \oplus M_{d_2})$ (by $L_\omega$-mono: any $W'$ with $\dim(W') > \dim(W_*)$ has $\omega(\pi_{W'}) > \omega(\pi_{W_*})$, so $W_*$ cannot be cost-minimizing if it has excess dimension); (d)~$M_{d_i} \cap W_*^\perp = \{0\}$ for $i = 1,2$ (otherwise $d_i$ is undefended).

\smallskip\noindent\textit{Bridge: $\kappa \to \omega$.}  The cost functional $\omega(\pi_W) := \tr(\rho\,\pi_W)/\tr(\rho)$ (from O4, Proposition~\ref{prop:O4}) assigns to each codespace projection a value proportional to the capacity committed to defending that codespace.  For individual sector projections: $\omega(E_d) \propto \eps(d)$ (O4, cost recovery), so $\omega$ is monotone in enforcement cost.  Minimizing $\omega(\pi_W)$ over admissible codespaces $W$ (those that defend $d_1$ and $d_2$ against all perturbations, including $\Pi$-supported ones) is therefore equivalent to minimizing the committed capacity for joint defense.  The minimizer $W_*$ with respect to $\omega$ is the codespace that achieves the joint enforcement cost $\eps(\{d_1, d_2\})$.

\noindent\textit{Proof of (a).}  The cost $\omega(\pi_W)$ is a continuous function on the Grassmannian $\Gr(k, S_\Gamma)$, which is compact (from OR3: $\dim(S_\Gamma) < \infty$).  The feasibility constraint ($W$ must defend $d_1$ and $d_2$ against all perturbations at $\Gamma$, including $\Pi$-supported ones per Remark~\ref{rem:joint_cost}) defines a closed subset.  A continuous function on a compact set attains its infimum.

\noindent\textit{Consequence: incompatible position.}  By (b) and (c), $W_*$ is a $k$-plane different from $M_{d_1} \oplus M_{d_2}$ with $\dim(W_*) = k = \dim(M_{d_1}) + \dim(M_{d_2})$.  By (d), $M_{d_1} \not\subset W_*$ (otherwise $\dim(W_*) > k$ or $M_{d_2} \cap W_*^\perp \ne \{0\}$).  The subspaces $M_{d_1}$ and $W_*$ are in \emph{incompatible position}: $M_{d_1} \ne (M_{d_1} \cap W_*) \oplus (M_{d_1} \cap W_*^\perp)$.  By the projection-algebra theorem, $[E_{d_1}, \pi_{W_*}] \ne 0$.  \qed
\end{lemma}

\subsection{Noncommutativity ($T_{\mathrm{alg}}$)}

\begin{theorem}[$T_{\mathrm{alg}}$: The Full Enforcement Algebra is Noncommutative]\label{thm:Talg}
Let $d_1, d_2$ be co-located distinctions at $\Gamma$ with $\Delta(d_1, d_2) > 0$ (from $L_\Delta$).  Define the \emph{joint enforcement map} $E_{\{d_1,d_2\}} := \pi_{W_*}$ (the minimum-cost joint defender from $L_\Pi$) and the \emph{off-diagonal component}
\[
F_\Pi \;:=\; E_{\{d_1,d_2\}} - E_{d_1} - E_{d_2}.
\]
Then the full enforcement algebra $\mathcal{A} := \mathrm{Alg}_\mathbb{R}(\{E_d\} \cup \{E_{\{d_1,d_2\}}\})$ admits no faithful commutative $*$-representation.

\smallskip\noindent\textit{Proof} (four steps).

\textit{Step 1: the diagonal sector projections commute.}
By $T_{\mathrm{sector}}$(v), the sector maps $\{E_d\}$ satisfy $E_d E_{d'} = 0$ for $d \ne d'$ and generate the commutative diagonal subalgebra $\mathcal{A}_{\mathrm{diag}} \cong \mathbb{R}^{|\mathcal{B}|}$.

\textit{Step 2: the joint enforcement map lies outside the diagonal algebra.}
In $\mathcal{A}_{\mathrm{diag}}$, the only joint defender is $E_{d_1} + E_{d_2}$ (block-diagonal), which has cost $\eps(d_1) + \eps(d_2)$ (exact additivity in the diagonal algebra, from the ``diagonal algebras force additive cost'' bridge in $L_{\mathrm{blk}}$).  But $\Delta > 0$ gives $\eps(\{d_1,d_2\}) > \eps(d_1) + \eps(d_2)$.  Therefore $E_{\{d_1,d_2\}} \ne E_{d_1} + E_{d_2}$, and $F_\Pi = E_{\{d_1,d_2\}} - E_{d_1} - E_{d_2} \ne 0$.

Properties of $F_\Pi$: self-adjoint ($F_\Pi^* = F_\Pi$, since $E_{\{d_1,d_2\}}, E_{d_1}, E_{d_2}$ are all self-adjoint by $T_{\mathrm{adj}}$); nontrivially $\Pi$-supported (by Corollary~\ref{cor:commutator}(b): some $v \in M_{d_1}$ has $E_{\{d_1,d_2\}}v = v + \pi_v$ with $\pi_v \in \Pi \setminus \{0\}$).

\textit{Step 3: $F_\Pi$ fails to commute with a sector projection.}
By Corollary~\ref{cor:commutator}(c): there exists $v \in M_{d_1}$ such that
\[
[E_{d_1},\, F_\Pi]\,v \;=\; -\pi_v \;\ne\; 0,
\]
where $\pi_v \in \Pi \setminus \{0\}$ is the $\Pi$-component of $E_{\{d_1,d_2\}}v$.  (The computation: $E_{d_1}(F_\Pi v) = E_{d_1}(\pi_v) = 0$ since $\pi_v \in \Pi \subseteq \ker(E_{d_1})$; while $F_\Pi(E_{d_1}v) = F_\Pi(v) = \pi_v$.)  Therefore $[E_{d_1}, F_\Pi] \ne 0$ in $\End(V)$.

\textit{Step 4: $\mathcal{A}$ admits no faithful commutative representation.}
Suppose $\phi: \mathcal{A} \to \mathcal{C}$ is a faithful $*$-homomorphism with $\mathcal{C}$ commutative.  Then $\phi([E_{d_1}, F_\Pi]) = [\phi(E_{d_1}), \phi(F_\Pi)] = 0$ (commutativity of $\mathcal{C}$).  But $[E_{d_1}, F_\Pi] \ne 0$ (Step~3) and $\ker\phi = \{0\}$ (faithfulness) give $\phi([E_{d_1}, F_\Pi]) \ne 0$.  Contradiction.  \qed

\smallskip\noindent\textit{Summary of the logical chain:}
\[
\Delta > 0 \;\xrightarrow{L_{\mathrm{blk}}}\; E_{\{d_1,d_2\}} \notin \mathcal{A}_{\mathrm{diag}} \;\xrightarrow{\text{define}}\; F_\Pi \ne 0 \;\xrightarrow{\text{Cor.~\ref{cor:commutator}}}\; [E_{d_1}, F_\Pi] \ne 0 \;\xrightarrow{\text{faithfulness}}\; \mathcal{A} \text{ noncommutative}.
\]

\smallskip\noindent\textit{Minimal finite witness} (self-contained; expanded in Section~\ref{sec:worked}).  Take $S_\Gamma = \mathbb{R}^3$, $M_{d_1} = \spn\{e_1\}$, $M_{d_2} = \spn\{e_2\}$, $\Pi = \spn\{e_3\}$, $B = I_3$.  The enforcement projections are $E_{d_1} = \diag(1,0,0)$, $E_{d_2} = \diag(0,1,0)$.  With $\Delta > 0$ (from $L_\Delta$: $\eps(\{d_1,d_2\}) > \eps(d_1) + \eps(d_2)$ since $\Pi \ne \{0\}$), the joint defender $\pi_{W_*}$ has a codespace $W_*$ rotated into $\Pi$ by angle $\theta \ne 0$:
\[
\pi_{W_*} = \begin{pmatrix} \cos^2\!\theta & 0 & \cos\theta\sin\theta \\ 0 & 1 & 0 \\ \cos\theta\sin\theta & 0 & \sin^2\!\theta \end{pmatrix}, \qquad
[E_{d_1},\, \pi_{W_*}] = \begin{pmatrix} 0 & 0 & \cos\theta\sin\theta \\ 0 & 0 & 0 \\ -\cos\theta\sin\theta & 0 & 0 \end{pmatrix} \ne 0.
\]
The nonzero commutator witnesses the obstruction.  \textit{Verification:} on $\spn\{e_1, e_3\}$, the generators $E_{d_1}|_{\spn\{e_1,e_3\}} = \bigl(\begin{smallmatrix} 1 & 0 \\ 0 & 0 \end{smallmatrix}\bigr)$ and $\pi_{W_*}|_{\spn\{e_1,e_3\}} = \bigl(\begin{smallmatrix} \cos^2\!\theta & \cos\theta\sin\theta \\ \cos\theta\sin\theta & \sin^2\!\theta \end{smallmatrix}\bigr)$ generate the full $2 \times 2$ matrix algebra $M_2(\mathbb{R})$: explicitly, $E_{d_1}$ gives the diagonal idempotent, $\pi_{W_*}$ gives a rank-$1$ projector with nonzero off-diagonal entries, and their commutator $[E_{d_1}, \pi_{W_*}]$ provides an antisymmetric generator; products and linear combinations of these three elements span all four basis matrices $\{E_{11}, E_{12}, E_{21}, E_{22}\}$.  On $\spn\{e_2\}$, all generators act as scalars: $\mathbb{R}$.  Therefore $\mathcal{A} = M_2(\mathbb{R}) \oplus \mathbb{R}$, which admits no faithful commutative $*$-representation (the $M_2(\mathbb{R})$ block is non-abelian).  At $\theta = 0$: $\pi_{W_*} = E_{d_1} + E_{d_2}$, $[E_{d_1}, \pi_{W_*}] = 0$, $\Delta = 0$, and the algebra is commutative (classical regime).

\noindent\textit{Scope.}  $T_{\mathrm{alg}}$ establishes: finite-dimensional, associative, noncommutative real $*$-algebra with faithful positive state.  Deferred: Wedderburn--Artin (T2a), field selection $\mathbb{C}$ (T2c.1--T2c.2).

\noindent\textit{Jordan algebra exclusion.}  $\mathcal{A} \subseteq \End(V)$: the product is operator composition (associative).  Jordan algebras use $A \circ B = \frac{1}{2}(AB+BA)$ (non-associative).  Associativity is inherited, not axiomatized.

\noindent\textit{Classical limit ($\Pi = \{0\}$).}  If the pool is empty (capacity perfectly saturated by mutually exclusive distinctions with no shared substrate), then $\Delta = 0$, all commutators vanish, and the algebra is commutative: $\mathcal{A} \cong \bigoplus_d \mathbb{R}$.  This is the correct classical regime --- a probability simplex.  Noncommutativity is \emph{conditional} on $\Pi \ne \{0\}$, which holds at any interface with $\dim(S_\Gamma) \ge 3$ (CL3 in $L_\Delta$).  The framework does not force noncommutativity universally; it derives it precisely where the pool mediates cross-sector coupling.
\end{theorem}

\subsection{Worked example: the complete construction on $\mathbb{R}^3$}\label{sec:worked}

\noindent\textit{Scope.}  This example is an explicit finite witness instantiating the general constructions from $\kappa$ through $T_{\mathrm{Born}}$ on a single interface.  It uses the direct pool-rotation route (Route~B above), not the BW route.  Where a later theorem is conditional in general (T2c.2: $\mathbb{H}$-exclusion requires composition closure), the example illustrates only the single-interface portion (T2c.1: $\mathbb{R}$-exclusion via state-separation completeness).

\noindent\textit{This section instantiates the entire derivation chain --- from $\kappa$ and $B$ through $T_{\mathrm{adj}}$, $L_\Delta$, $L_\Pi$, $T_{\mathrm{alg}}$, T2c.1, GNS, to $T_{\mathrm{Born}}$ --- in a single minimal three-dimensional interface.  Every step is an explicit matrix computation.}

\smallskip
We exhibit the full chain $\kappa \to B \to T_{\mathrm{adj}} \to L_\Delta \to L_\Pi \to T_{\mathrm{alg}} \to \text{T2c.1} \to T_{\mathrm{GNS}} \to T_{\mathrm{Born}}$ in a minimal three-dimensional interface.

\paragraph{Setup.}
$S_\Gamma = \mathbb{R}^3$ with standard basis $\{e_1, e_2, e_3\}$.  Two distinctions $d_1, d_2$ with anchor sets $M_{d_1} = \spn\{e_1\}$, $M_{d_2} = \spn\{e_2\}$, pool $\Pi = \spn\{e_3\}$.  Set enforcement costs $\eps(d_1) = 2\eps^*$, $\eps(d_2) = \eps^*$, $\eps(d_\Gamma) = \eps^*$ (the pool direction $e_3$ is non-null by the $\mathcal{D}$-quotient; perturbations along it have positive cost by K2 and $L_{\eps^*}$).  Total capacity $C = 5\eps^*$.

\paragraph{Cost function $\kappa$ and bilinear form $B$.}
By K3 (Theorem~\ref{thm:K3}), the cost function is additive on the $T_{\mathrm{sep}}$ sectors.  The sector-diagonal bilinear form from $L_\omega$ (Lemma~\ref{lem:Lomega}) is:
\[
B = \diag(2\eps^*, \;\eps^*, \;\eps^*),
\]
with $B(e_i, e_j) = 0$ for $i \ne j$ (cross-sector orthogonality from K3).  Normalization is conventional; the inter-sector vanishing is forced.

\paragraph{Enforcement operations ($T_{\mathrm{adj}}$).}
$E_{d_1} = \diag(1,0,0)$ and $E_{d_2} = \diag(0,1,0)$ are the $B$-orthogonal projections onto $M_{d_1}$ and $M_{d_2}$.  Self-adjointness: $B(E_{d_i}\,u,\, v) = B(u,\, E_{d_i}\,v)$ for all $u, v$ (verified by direct computation using sector-orthogonality of $B$).

\paragraph{Superadditivity ($L_\Delta$).}
Individual costs: $\eps(d_1) + \eps(d_2) = 3\eps^*$.  Joint cost: maintaining both $d_1$ and $d_2$ at the shared interface also requires defending $d_\Gamma$ (Remark~\ref{rem:joint_cost}), so $\eps(\{d_1, d_2\}) \ge 3\eps^* + \eps^* = 4\eps^*$.  Superadditive surplus: $\Delta = \eps^* > 0$.

\paragraph{Codespace rotation ($L_\Pi$).}
The minimum-cost joint defender $\pi_{W_*}$ must project onto a 2-plane $W_*$ that covers both $M_{d_1}$ and $\Pi$ (to defend $d_1$ and $d_\Gamma$).  Take $W_* = \spn\{(\cos\theta,\, 0,\, \sin\theta),\; (0,\, 1,\, 0)\}$ for $\theta \ne 0$.  The projection matrix onto $W_*$ is:
\[
\pi_{W_*} = \begin{pmatrix} \cos^2\theta & 0 & \cos\theta\sin\theta \\ 0 & 1 & 0 \\ \cos\theta\sin\theta & 0 & \sin^2\theta \end{pmatrix}.
\]
At $\theta = 0$: $\pi_{W_*} = E_{d_1} + E_{d_2}$ (block-diagonal, no pool support, $\Delta = 0$).  At $\theta \ne 0$: $W_*$ is rotated into $\Pi$, and $\pi_{W_*}$ has off-diagonal entries connecting $M_{d_1}$ and $\Pi$.

\paragraph{Noncommutativity ($T_{\mathrm{alg}}$).}
The joint enforcement map is $E_{\{d_1,d_2\}} = \pi_{W_*}$.  The off-diagonal component:
\[
F_\Pi := E_{\{d_1,d_2\}} - E_{d_1} - E_{d_2} = \begin{pmatrix} \cos^2\!\theta - 1 & 0 & \cos\theta\sin\theta \\ 0 & 0 & 0 \\ \cos\theta\sin\theta & 0 & \sin^2\!\theta \end{pmatrix} = \begin{pmatrix} -\sin^2\!\theta & 0 & \cos\theta\sin\theta \\ 0 & 0 & 0 \\ \cos\theta\sin\theta & 0 & \sin^2\!\theta \end{pmatrix} \ne 0 \quad (\theta \ne 0).
\]
$F_\Pi$ is self-adjoint ($F_\Pi^\top = F_\Pi$) and has nonzero $\Pi$-support (the $(3,3)$ entry $\sin^2\!\theta > 0$).  Direct computation of the commutator:
\[
[E_{d_1},\, F_\Pi] = [E_{d_1},\, \pi_{W_*}] = \begin{pmatrix} 0 & 0 & \cos\theta\sin\theta \\ 0 & 0 & 0 \\ -\cos\theta\sin\theta & 0 & 0 \end{pmatrix} \ne 0 \quad (\theta \ne 0).
\]
The generated algebra $\mathcal{A} = \mathrm{Alg}_\mathbb{R}\{E_{d_1}, E_{d_2}, \pi_{W_*}\}$ decomposes as: $M_2(\mathbb{R}) \oplus \mathbb{R}$.  On the 2-plane $\spn\{e_1, e_3\}$, the generators produce a diagonal idempotent, a rank-$1$ projector with off-diagonal entries, and their commutator; products and linear combinations span all four basis matrices of $M_2(\mathbb{R})$.  On $\spn\{e_2\}$, all generators act as scalars: this is the $\mathbb{R}$ block.

\paragraph{Field selection (T2c.1: $\mathbb{R}$-incompleteness visible).}
The $M_2(\mathbb{R})$ block contains one antisymmetric element (up to scale):
\[
A := [E_{d_1},\, \pi_{W_*}]|_{\spn\{e_1,e_3\}} = \cos\theta\sin\theta \begin{pmatrix} 0 & 1 \\ -1 & 0 \end{pmatrix}.
\]
Over $\mathbb{R}$, every density matrix on $\spn\{e_1, e_3\}$ is symmetric: $\rho = \rho^\top$.  Therefore $\tr(\rho\, A) = 0$ for all physical states --- the antisymmetric element is state-invisible.  Extending to $\mathbb{C}$, the observable $iA$ is Hermitian, and $\tr(\rho\, iA) \ne 0$ for generic Hermitian $\rho$.  The per-block parameter count becomes $2^2 = 4$ (matching $M_2(\mathbb{C})$) rather than $2(3)/2 = 3$ (the real self-adjoint count).  The example exhibits the real-block defect repaired by complexification (T2c.1, via state-separation completeness).

\paragraph{GNS construction ($T_{\mathrm{GNS}}$) on the $\mathbb{R}^3$ example.}
The algebra $\mathcal{A} = M_2(\mathbb{R}) \oplus \mathbb{R}$ has faithful positive state $\omega(A) := \tr(\rho\, A)/\tr(\rho)$ with $\rho = \diag(2\eps^*, \eps^*, \eps^*)$.  The GNS bilinear form is $\langle A, B\rangle_\omega := \omega(A^* B) = \tr(\rho\, A^\top B)/\tr(\rho)$.  The GNS Hilbert space $\mathcal{H}_\omega = (\mathcal{A}, \langle\cdot,\cdot\rangle_\omega)$ has $\dim(\mathcal{H}_\omega) = \dim(\mathcal{A}) = 5$ (four from $M_2(\mathbb{R})$, one from $\mathbb{R}$).  The representation $\pi_\omega(A)B = AB$ is faithful ($\pi_\omega(A) = 0 \Rightarrow A = 0$).

\paragraph{Born rule on the $\mathbb{R}^3$ example.}
After complexification ($\mathbb{F} = \mathbb{C}$ from the preceding step), the algebra becomes $M_2(\mathbb{C}) \oplus \mathbb{C}$.  On $M_2(\mathbb{C})$, define $p(E) := \omega(E)$ for all effects $0 \le E \le \mathrm{Id}$ ($L_{\mathrm{prob}}$, Lemma~\ref{lem:Lprob}).  Busch's hypotheses are verified: non-contextuality and ortho-additivity (from linearity of $\omega$, $L_{\mathrm{prob}}$(c,e)), normalization ($\omega(\mathbf{1}) = 1$, $L_{\mathrm{prob}}$(b)).  Since $\dim \ge 2$, Busch's theorem gives $p(E) = \tr(\rho E)$ for a unique density operator $\rho$ and every effect $E$.  Concretely: $E_{d_1} = \diag(1,0)$ on the 2-plane yields $\Pr(d_1) = \tr(\rho\, E_{d_1}) = \rho_{11}$.  This is the Born rule on the $\mathbb{R}^3$ example.  The single-interface chain $\mathrm{A1} \to T_{\mathrm{Born}}$ is instantiated in this model (with $\mathbb{R}$-exclusion via T2c.1's state-separation completeness; $\mathbb{H}$-exclusion via T2c.2 is not demonstrated in this single-interface example).

\subsection{Qubit model: $M_2(\mathbb{C})$ from a two-distinction interface}
\label{sec:qubit_model}

\noindent\textit{Scope.}  This example demonstrates the minimal nonclassical block and the Born-rule machinery on a concrete finite interface using Route~B (direct pool rotation).  It instantiates the full single-interface chain through $T_{\mathrm{Born}}$; it does not by itself settle the composition-closure tier (T2c.2), which requires a second independent interface.

We construct the entire derivation chain for the simplest non-trivial case: a single interface $\Gamma$ with two irreducible distinctions.  Every step is a concrete matrix computation.

\paragraph{Setup.}
$S_\Gamma = \mathbb{R}^3$ with standard basis $\{e_1, e_2, e_3\}$.  Two irreducible distinctions $d_1, d_2$ with anchor sets $M_{d_1} = \spn\{e_1\}$, $M_{d_2} = \spn\{e_2\}$, pool $\Pi = \spn\{e_3\}$.  Perturbation cost function: $\kappa(p) = \|p - \mathrm{id}\|_F^2$ (squared Frobenius norm; satisfies K1--K3 since $\|A \oplus B\|_F^2 = \|A\|_F^2 + \|B\|_F^2$; this places the model in regime~(i) of Proposition~\ref{prop:kappa_class}).  Budget $C = 5$.

\paragraph{Step 1: enforcement costs ($\eps$, $L_{\mathrm{cost}}$).}
The minimum-cost perturbation destroying $d_1$ is $p_1 = \diag(0,1,1)$: $\kappa(p_1) = \|\diag(-1,0,0)\|_F^2 = 1$.  Set $\eps^* := 1$.  Then $\eps(d_1) = \eps(d_2) = 1$.  The pool direction $e_3$ is non-null ($\mathcal{D}$-quotient); by direct computation, the minimum-cost perturbation along $e_3$ has $\kappa = 1$, so $\eps(d_\Pi) = 1$.

\smallskip\noindent\textit{Note on BW.}  In this model, all three costs are equal ($\eps^* = 1$); the cost spectrum is degenerate.  BW's budget-window triple requires $\eps(d_1) < \eps(d_2)$, which is not satisfied.  This example reaches noncommutativity via Route~B (direct pool rotation: $\Pi \ne \{0\}$, $\Delta > 0$), not via Route~A (BW $\to$ T1).  BW is not needed for noncommutativity at this interface; it would be needed only if one required the general order-dependence witness T1, which this example bypasses.

\paragraph{Step 2: sector decomposition ($T_{\mathrm{sep}}$).}
$M_{d_1} \cap M_{d_2} = \{0\}$; by $T_{\mathrm{sep}}$: $\eps(\{d_1,d_2\}) = 2$ for perturbations supported only in $M_{d_1} \cup M_{d_2}$.  Direct sum: $S_\Gamma = \spn\{e_1\} \oplus \spn\{e_2\} \oplus \spn\{e_3\}$.

\paragraph{Step 3: bilinear form ($L_\omega$).}
K3 forces $B(e_i, e_j) = 0$ for $i \ne j$.  Choose $B = I_3$ (Euclidean).  Projections: $E_{d_1} = \diag(1,0,0)$, $E_{d_2} = \diag(0,1,0)$.

\paragraph{Step 4: self-adjointness ($T_{\mathrm{adj}}$).}
$E_{d_i}^\top = E_{d_i}$: $B(E_{d_i} u, v) = B(u, E_{d_i} v)$ for all $u,v$.

\paragraph{Step 5: superadditivity ($L_\Delta$).}
The pool $\Pi = \spn\{e_3\}$ is non-null (SP: configurations differ along $e_3$).  A $\Pi$-supported perturbation $p_3 = \diag(1,1,0)$ has $\kappa(p_3) = \|\diag(0,0,-1)\|_F^2 = 1 > 0$ (K2).  This perturbation is in $\mathcal{P}(\{d_1,d_2\})$ (it changes physically meaningful DOF at $\Gamma$ at positive cost) but not in $\mathcal{P}(d_1) \cup \mathcal{P}(d_2)$ (it doesn't destroy $d_1$ or $d_2$, whose anchors are in $\spn\{e_1\}$ and $\spn\{e_2\}$).  Defending against it requires $\Pi$-committed capacity: the defense projects onto a 2-plane $W$ containing $e_1$ and $e_2$, at additional cost $\kappa(\mathrm{Id} - \pi_W) = 1$.  By K3: $\eps(\{d_1,d_2\}) = \eps(d_1) + \eps(d_2) + 1 = 3$.  Surplus: $\Delta = 1 > 0$.

\paragraph{Step 6: codespace rotation ($L_\Pi$, $L_{\mathrm{blk}}$).}
The block-diagonal defender $\diag(1,1,0)$ cannot defend against $\Pi$-supported attacks (it annihilates $e_3$).  The minimum-cost joint defender projects onto $W_* = \spn\{(\cos\theta, 0, \sin\theta),\, (0,1,0)\}$:
\[
\pi_{W_*} = \begin{pmatrix} \cos^2\!\theta & 0 & \cos\theta\sin\theta \\ 0 & 1 & 0 \\ \cos\theta\sin\theta & 0 & \sin^2\!\theta \end{pmatrix}, \quad \theta \ne 0.
\]
The committed cost is $\omega(\pi_{W_*}) = [\eps(d_1)\cos^2\!\theta + \eps(d_\Pi)\sin^2\!\theta + \eps(d_2)]/C = [1\cdot\cos^2\!\theta + 1\cdot\sin^2\!\theta + 1]/5 = 3/5$, independent of $\theta$ (since $\eps(d_1) = \eps(d_\Pi) = 1$).  Any $\theta \in (0, \pi/2)$ is a valid minimizer; the rotation angle parameterizes the family of optimal codespaces.  The key structural point: $W_* \ne M_{d_1} \oplus M_{d_2}$ for any $\theta \ne 0$.

\paragraph{Step 7: noncommutativity ($T_{\mathrm{alg}}$).}
\[
[E_{d_1}, \pi_{W_*}] = \begin{pmatrix} 0 & 0 & \cos\theta\sin\theta \\ 0 & 0 & 0 \\ -\cos\theta\sin\theta & 0 & 0 \end{pmatrix} \ne 0.
\]
On $\spn\{e_1, e_3\}$: $M_2(\mathbb{R})$.  On $\spn\{e_2\}$: $\mathbb{R}$.  Therefore $\mathcal{A} \cong M_2(\mathbb{R}) \oplus \mathbb{R}$.

\paragraph{Step 8: Wedderburn--Artin (T2a).}
$\dim(\mathcal{A}) = 5$, semisimple ($\omega(N^*N) = 0 \Rightarrow N = 0$ by faithfulness of $\omega = \tr(\cdot)/3$).  Decomposition: $M_2(\mathbb{R}) \oplus \mathbb{R}$.

\paragraph{Step 9: GNS ($T_{\mathrm{GNS}}$).}
$\omega(A) = \tr(A)/3$ gives $\langle A, B\rangle_\omega = \tr(A^\top B)/3$.  Hilbert space: $\mathcal{H}_\omega = (\mathcal{A}, \langle\cdot,\cdot\rangle_\omega)$, $\dim = 5$.  Representation: $\pi_\omega(A)B = AB$, faithful.

\paragraph{Step 10: field selection (T2c.1: single-interface $\mathbb{R}$-exclusion).}
The $M_2(\mathbb{R})$ block contains $J = \bigl(\begin{smallmatrix} 0 & 1 \\ -1 & 0 \end{smallmatrix}\bigr)$ (from $[E_{d_1}, \pi_{W_*}]$, rescaled).  Over $\mathbb{R}$: $\tr(\rho J) = 0$ for all symmetric $\rho$.  Complexification: $iJ$ is Hermitian with $\tr(\rho\, iJ) \ne 0$ generically.  State-space parameters rise from $3$ to $4$.  Therefore $\mathbb{F} = \mathbb{C}$; $\mathcal{A}_\mathbb{C} \cong M_2(\mathbb{C}) \oplus \mathbb{C}$.

\paragraph{Step 11: Born rule ($T_{\mathrm{Born}}$).}
On $M_2(\mathbb{C})$ ($\dim \ge 2$), define $p(E) := \omega(E)$ for all effects $0 \le E \le \mathrm{Id}$ ($L_{\mathrm{prob}}$).  By Busch's theorem: $\exists!\,\rho$ with $p(E) = \tr(\rho E)$ for every effect $E$.

\smallskip\noindent\textit{Micro-example: projections.}  Take orthogonal rank-$1$ projections $P_1 = \diag(1,0)$ and $P_2 = \diag(0,1)$ on $\mathbb{C}^2$.  Then:
\[
p(P_1 + P_2) = \omega(\mathrm{Id}) = 1 = p(P_1) + p(P_2). \quad\checkmark
\]
For $\rho = \diag(2/3, 1/3)$: $p(P_1) = 2/3$, $p(P_2) = 1/3$.

\smallskip\noindent\textit{Micro-example: non-projective effect.}  Take $E = \tfrac{1}{2}P_1 + \tfrac{3}{4}P_2 = \diag(1/2,\, 3/4)$.  This is a genuine effect ($0 \le E \le \mathrm{Id}$) that is not a projection.  Then:
\[
p(E) = \omega(E) = \tr(\rho\, E) = \tfrac{2}{3}\cdot\tfrac{1}{2} + \tfrac{1}{3}\cdot\tfrac{3}{4} = \tfrac{1}{3} + \tfrac{1}{4} = \tfrac{7}{12}.
\]
This equals $\tfrac{1}{2}\,p(P_1) + \tfrac{3}{4}\,p(P_2) = \tfrac{1}{2}\cdot\tfrac{2}{3} + \tfrac{3}{4}\cdot\tfrac{1}{3} = \tfrac{7}{12}$: the affine property ($L_{\mathrm{prob}}$(d)) is verified.  The Born rule extends seamlessly from projections to effects via the spectral decomposition, exactly as Busch's theorem requires.

\smallskip\noindent\textit{Explicit verification.}  $E_{d_1} \mapsto P_\uparrow = |{\uparrow}\rangle\langle{\uparrow}|$ in $\mathbb{C}^2$.  For $\rho = I/2$: $\Pr(d_1) = \tr(\rho P_\uparrow) = 1/2$.  For $|\psi\rangle = \cos(\alpha/2)|{\uparrow}\rangle + \sin(\alpha/2)|{\downarrow}\rangle$:
\[
\Pr(d_1) = |\langle{\uparrow}|\psi\rangle|^2 = \cos^2(\alpha/2).
\]
This completes the chain $\mathrm{A1} \to \eps(\cdot) \to T_{\mathrm{sep}} \to L_\omega \to T_{\mathrm{adj}} \to L_\Delta \to L_{\mathrm{blk}} \to L_\Pi \to T_{\mathrm{alg}} \to \mathrm{T2a} \to T_{\mathrm{GNS}} \to \mathrm{T2c.1} \to T_{\mathrm{Born}}$ in a fully explicit, finite, computable model.  Every step is a concrete matrix computation; no appeals to general theory are needed beyond Wedderburn--Artin and Busch.  ($\mathbb{H}$-exclusion (T2c.2) requires a second interface and is not demonstrated in this single-interface example.)

% ═══════════════════════════════════════════════════════════════
\section{The Quantum Skeleton}
\label{sec:skeleton}

\subsection{Anti-circularity audit}

The Hilbert space is the \emph{output} of the construction, not an input:

\smallskip
\begin{center}\small
\begin{tabular}{@{}clll@{}}
\toprule
& \textbf{Step} & \textbf{Inputs (all pre-Hilbert)} & \textbf{Output} \\
\midrule
1 & Real vector space $V$ & A1 + OR1 + $L_{\eps^*}$ & $\dim V \le n_{\max}+1$ \\
2 & Sector decomposition & $T_{\mathrm{sep}}$ (pre-metric) & $S_\Gamma = M_d \oplus N_d$ \\
3 & Cost bilinear form $B$ & $T_{\mathrm{sep}}$ + K3 ($L_\omega$) & Sector-orthogonal $B$ \\
4 & Self-adjointness & $N_d = M_d^{\perp_B}$ ($T_{\mathrm{adj}}$) & $E_d = E_d^\dagger$ \\
5 & $*$-algebra $\mathcal{A}$ & $\End(V)$ + generators (O3) & Assoc., $*$-involutive \\
6 & Positive state $\omega$ & $\tr$ on $\End(V)$ + $\rho > 0$ (O4) & $\omega(A^*A) \ge 0$, faithful \\
7 & GNS $\to \mathcal{H}_\omega$ & Steps 5+6 as input & \textbf{Hilbert space} \\
8 & $C^*$-norm & Op.\ norm on $B(\mathcal{H}_\omega)$ & $C^*$-identity \\
\bottomrule
\end{tabular}
\end{center}

\begin{remark}[Where the pre-Hilbert form arises, and why the result is unique]\label{rem:preHilbert}
The construction involves \emph{two} inner products at different stages, serving different roles.  Neither presupposes a Hilbert space.

\smallskip\noindent\textbf{Inner product 1: the sector-orthogonal form $B$ (Step~3, pre-algebraic).}  $B: V \times V \to \mathbb{R}$ is constructed by $L_\omega$ (Lemma~\ref{lem:Lomega}) from the $T_{\mathrm{sep}}$ sector decomposition and K3.  It is a positive-definite \emph{real} bilinear form on the finite-dimensional vector space $V$ (from $T_{\mathrm{embed}}$).  Its role: define the $B$-adjoint ($B(Au, v) = B(u, A^*v)$, Step~7 of the construction box) and identify enforcement maps as $B$-orthogonal projections ($T_{\mathrm{adj}}$, Step~8).  $B$ is \emph{not} the Hilbert space inner product; it is a pre-metric geometric structure on the substrate.

\smallskip\noindent\textbf{Inner product 2: the GNS form $\langle A, B\rangle_\omega := \omega(A^*B)$ (Step~7, post-algebraic).}  Once the $*$-algebra $\mathcal{A}$ is generated (O3, Step~5) and a faithful positive state $\omega$ is constructed (O4, Step~6), the GNS inner product $\langle A, B\rangle_\omega = \omega(A^*B)$ is defined on $\mathcal{A}$ itself.  Faithfulness of $\omega$ ensures positive-definiteness ($\omega(A^*A) = 0 \Rightarrow A = 0$).  The completion of $(\mathcal{A}, \langle\cdot,\cdot\rangle_\omega)$ is the Hilbert space $\mathcal{H}_\omega$ (trivial in finite dimensions: already complete).  The operator norm and $C^*$-identity ($\|A^*A\| = \|A\|^2$) are defined only after this step (Step~8), on $B(\mathcal{H}_\omega)$.  No operator norm, spectral theory, or $C^*$-structure is used before Step~7.

\smallskip\noindent\textbf{Uniqueness up to $*$-isomorphism.}  Two choices enter the construction: the within-sector normalization of $B$ (Step~3), and the faithful state $\omega$ (Step~6).  Neither affects the $*$-isomorphism class of the output.

\noindent(i)~\textit{Independence of $B$.}  Changing the within-sector forms $B_d$ (while preserving inter-sector orthogonality, which is forced by K3) conjugates each $E_d$ by a positive-definite matrix within its sector.  This is a $*$-isomorphism of $\mathcal{A}$: the algebra's abstract structure (which products vanish, which commutators are nonzero, which blocks appear in the Wedderburn decomposition) is unchanged.

\noindent(ii)~\textit{Independence of $\omega$.}  In finite dimensions, any two faithful positive states $\omega, \omega'$ on the same $*$-algebra yield unitarily equivalent GNS representations: $\pi_\omega \cong \pi_{\omega'}$ (this is a standard result; faithfulness ensures both representations are injective, and finite-dimensionality forces unitary equivalence via the Skolem--Noether theorem applied to the image algebras).  Therefore $\mathcal{H}_\omega \cong \mathcal{H}_{\omega'}$ as Hilbert spaces, and the $C^*$-algebra $\pi_\omega(\mathcal{A})$ is unique up to $*$-isomorphism.

\noindent The output --- a finite-dimensional $C^*$-algebra $\cong \bigoplus_k M_{n_k}(\mathbb{C})$ (after T2c's field selection) --- depends only on the enforcement data (A1, MD, the distinction set $\mathcal{D}(\Gamma)$), not on normalization conventions or state choices.
\end{remark}

\subsection{Algebra construction (O1--O4, T2a)}

\begin{proposition}[O1: Linear Extension]\label{prop:O1}
Each enforcement map $\mathfrak{E}_d: K \to K \cup \{\bot\}$ extends uniquely to $E_d \in \End(V)$ (convention: $\bot = 0$).

\noindent\textit{Proof.}  $\mathfrak{E}_d$ is affine on the subdomain $K_d := \{\sigma \in K : \delta_\sigma \ge \eps(d)\}$ where enforcement succeeds (acts as a coordinate translation in the cost-vector embedding).  On $K \setminus K_d$ (budget exceeded), $\mathfrak{E}_d(\sigma) = \bot$.  We identify $\bot = 0 \in V$: this is forced because $0$ is the unique element of $V$ that is not the image of any admissible enforcement path (every $\varphi(\sigma) \in K$ has at least one nonzero coordinate), so $\bot$ maps to the only ``unreachable'' element of $V$.  The affine map $\mathfrak{E}_d|_{K_d}$ extends uniquely to a linear map $E_d: V \to V$ (standard: affine maps on spanning subsets of $\mathbb{R}^N$ have unique linear extensions).  On $K \setminus K_d$, the assignment $\mathfrak{E}_d = 0$ is consistent with the linear extension because these states map to the zero element by the budget-exceeded convention.  \qed
\end{proposition}

\begin{proposition}[O3: Enforcement Algebra]\label{prop:O3}
$\mathcal{A} := \mathrm{Alg}_\mathbb{R}(\{E_d\} \cup \{E_{\{d_1,d_2\}}\}) \subseteq \End(V)$ is a finite-dimensional real $*$-algebra.

\noindent\textit{Proof.}  Closure under addition: $\mathcal{A}$ is defined as a linear span.  Closure under multiplication: the product of monomials in generators is another monomial, hence in the span.  Associativity: inherited from function composition in $\End(V)$ (not axiomatized).  Identity: $\mathrm{Id}_V \in \End(V)$; since $\mathcal{A} \subseteq \End(V)$ and $\End(V)$ contains $\mathrm{Id}_V$, we include $\mathrm{Id}_V$ in $\mathcal{A}$ by defining $E_\Pi := \mathrm{Id}_V - \sum_{d \in \mathcal{B}} E_d$ (the projection onto $\Pi$; this is well-defined since $\{E_d\}_{d \in \mathcal{B}} \cup \{E_\Pi\}$ are pairwise orthogonal projections summing to $\mathrm{Id}_V$ by the $T_{\mathrm{sep}}$ decomposition).  $*$-involution: define $G^* := G$ for all generators (self-adjoint by $T_{\mathrm{adj}}$); extend anti-multiplicatively: $(AB)^* := B^*A^*$.  Involutivity: $(A^*)^* = A$ (each generator is its own adjoint; anti-multiplicative extension preserves involutivity).  Finite-dimensionality: $\mathcal{A} \subseteq \End(V)$, $\dim(\End(V)) = (\dim V)^2 \le (n_{\max}+1)^2 < \infty$.  \qed

\smallskip\noindent\textit{$C^*$-algebra structure.}  $\mathcal{A}$ is automatically a real $C^*$-algebra.  The operator norm $\|A\|_{\mathrm{op}} := \sup_{\|v\|_B = 1} \|Av\|_B$ on $\End(V)$ restricts to $\mathcal{A}$.

\textit{$C^*$-identity verification:} for any $A \in \mathcal{A}$, $\|A^*A\|_{\mathrm{op}} = \sup_{\|v\|=1} \|A^*Av\|$.  Since $* = \dagger_B$ (W4 of T2a), we have $B(A^*Av, v) = B(Av, Av) = \|Av\|_B^2$ for all $v$.  Therefore $\|A^*A\|_{\mathrm{op}} = \sup_{\|v\|=1} \|Av\|_B^2 = \|A\|_{\mathrm{op}}^2$.  Completeness is automatic (finite-dimensional normed spaces are complete).

\textit{Positive cone and ordering.}  The set $\mathcal{A}^+ := \{A^*A : A \in \mathcal{A}\}$ is the positive cone of $\mathcal{A}$.  It defines a partial order on self-adjoint elements: $A \ge B$ iff $A - B \in \mathcal{A}^+$.  Positivity of a state $\omega$ (O4) means $\omega(\mathcal{A}^+) \subseteq \mathbb{R}_{\ge 0}$; faithfulness means $\omega(A^*A) = 0 \Rightarrow A = 0$.  The cone is proper (closed, pointed, generating) by finite-dimensionality.

\textit{Sector projections as conditional expectations.}  Each $E_d$ is not only a self-adjoint projection but also a \emph{conditional expectation} onto the sector subalgebra $M_d \cdot \mathrm{Id}$: it is completely positive ($E_d(A^*A) \ge 0$ for $A$ supported in $M_d$), unital on the sector ($E_d|_{M_d} = \mathrm{id}$), and idempotent ($E_d^2 = E_d$).  This identification connects the enforcement-sector structure to the standard theory of conditional expectations in operator algebras (see Jones--Penneys~\cite{JonesPenneys2017}, \S4.1, for the categorical generalization).

\textit{Categorical context.}  The functoriality and universal property of GNS as a left adjoint of restriction (Parzygnat~\cite{Parzygnat2018}) and the generalization to $C^*$-algebra objects in rigid tensor categories (Jones--Penneys~\cite{JonesPenneys2017}) hold in the present setting because the finite-dimensional case is a special instance of the general theory.
\end{proposition}

\begin{proposition}[O4: Canonical Faithful State from Enforcement Data]\label{prop:O4}
The enforcement data $(T_{\mathrm{sep}}, \{E_d\}, \eps(d), C)$ determine a faithful positive normalized linear functional $\omega$ on $\mathcal{A}$, unique up to within-sector normalization gauge.

\smallskip\noindent\textit{Logical status: theorem, not postulate.}  O4 is not an additional structural assumption.  The faithful positive state $\omega$ is \emph{constructed} from objects already established: $\mathcal{A} \subseteq \End(V)$ (O3), the $B$-adjoint $* = \dagger_B$ (from $L_\omega$ + $T_{\mathrm{adj}}$), the enforcement costs $\eps(d)$ (from A1), and the sector projections $\{E_d\}$ ($T_{\mathrm{sep}}$ + $T_{\mathrm{adj}}$).  No external state, density matrix, or measurement postulate is imported.

Two constructions are given below, serving different roles.  Construction~1 is a pure existence proof: it shows O4 is not an extra axiom, since any finite-dimensional $*$-algebra with a $B$-adjoint carries a faithful state.  Construction~2 is the physically preferred APF state: it is built directly from the enforcement costs and sector structure, and it is the one that connects $\omega$ to the operational semantics of the framework.  Construction~2 is preferred over Construction~1 because it makes $\omega(E_d)$ monotone in the committed enforcement cost $\eps(d)$, preserving the operational content of A1's budget semantics in the algebraic state.  For the structural derivation (T2a, $T_{\mathrm{GNS}}$, T2c, $T_{\mathrm{Born}}$), any faithful positive normalized state suffices; the results are independent of which $\omega$ is chosen (Remark~\ref{rem:preHilbert}(ii)).

\medskip\noindent\textit{Construction 1: normalized trace (existence proof).}

\smallskip\noindent\textit{Definition.}  Set $n := \dim(V)$.  Since $V$ is a finite-dimensional real vector space ($T_{\mathrm{embed}}$), $\mathcal{A} \subseteq \End(V) \cong M_n(\mathbb{R})$ carries the standard matrix trace.  Define
\[
\tau(A) \;:=\; \frac{\tr(A)}{n}, \qquad A \in \mathcal{A}.
\]

\noindent\textit{Linearity.}  $\tau(\alpha A + \beta B) = \tr(\alpha A + \beta B)/n = \alpha\,\tr(A)/n + \beta\,\tr(B)/n = \alpha\,\tau(A) + \beta\,\tau(B)$.

\noindent\textit{Normalization.}  $\tau(\mathrm{Id}_V) = \tr(\mathrm{Id}_V)/n = n/n = 1$.

\noindent\textit{Positivity.}  Work in a $B$-orthonormal basis for $V$.  In this basis the $*$-involution coincides with matrix transpose (since $* = \dagger_B$ by~(W4) and the $B$-adjoint is the transpose in an orthonormal basis).  Then $\tau(A^*A) = \tr(A^\top A)/n = \|A\|_F^2/n \ge 0$, where $\|A\|_F$ is the Frobenius norm.

\noindent\textit{Faithfulness.}  $\tau(A^*A) = 0$ $\Rightarrow$ $\|A\|_F^2 = 0$ $\Rightarrow$ every matrix entry of $A$ vanishes $\Rightarrow$ $A = 0$.

\medskip\noindent\textit{Construction 2: cost-density state (physically canonical APF state).}

\smallskip\noindent\textit{Definition.}  The sector projections $\{E_d\}_{d \in \mathcal{B}}$ and the pool projection $E_\Pi := \mathrm{Id}_V - \sum_d E_d$ are pairwise $B$-orthogonal idempotents summing to $\mathrm{Id}_V$ ($T_{\mathrm{sep}}$ + $T_{\mathrm{adj}}$).  Each $E_d$ has rank $\dim(M_d)$ and $E_\Pi$ has rank $\dim(\Pi)$.  Define the \emph{cost density operator}
\[
\rho \;:=\; \sum_{d \in \mathcal{B}} \frac{\eps(d)}{C}\,E_d \;+\; \delta\, E_\Pi, \qquad \delta > 0\;\text{arbitrary (e.g.\ } \delta = \eps^*/C\text{)},
\]
using the enforcement costs $\eps(d) > 0$ (A1/FD4), the total capacity $C > 0$ (A1), and the cost floor $\eps^* > 0$ ($L_{\eps^*}$, from MD).  Define
\[
\omega(A) \;:=\; \frac{\tr(\rho\,A)}{\tr(\rho)}, \qquad A \in \mathcal{A}.
\]

\noindent\textit{Strict positivity of $\rho$.}  Every $v \in V$ decomposes uniquely as $v = \sum_d v_d + v_\Pi$ with $v_d \in M_d$, $v_\Pi \in \Pi$ (by the $T_{\mathrm{sep}}$ direct sum).  Therefore
\[
B(v,\, \rho\, v) \;=\; \sum_d \frac{\eps(d)}{C}\,\|v_d\|_B^2 \;+\; \delta\,\|v_\Pi\|_B^2.
\]
Every coefficient is strictly positive: $\eps(d)/C > 0$ (A1: $\eps(d) > 0$, $C > 0$) and $\delta > 0$ (by choice).  If $v \ne 0$, at least one component is nonzero, so $B(v, \rho\,v) > 0$.  Hence $\rho$ is positive-definite: $\rho > 0$.

\noindent\textit{Linearity.}  $\omega(\alpha A + \beta B) = \tr(\rho(\alpha A + \beta B))/\tr(\rho) = \alpha\,\omega(A) + \beta\,\omega(B)$, by linearity of the matrix trace.

\noindent\textit{Normalization.}  $\omega(\mathrm{Id}_V) = \tr(\rho \cdot \mathrm{Id}_V)/\tr(\rho) = \tr(\rho)/\tr(\rho) = 1$.

\noindent\textit{Positivity.}  Since $* = \dagger_B$ (W4), $A^*A = A^\dagger A$ is $B$-PSD: $B(v, A^\dagger A v) = \|Av\|_B^2 \ge 0$ for all $v$.  Let $\{e_i\}$ be a $B$-orthonormal eigenbasis of $\rho$ with eigenvalues $\lambda_i > 0$ (all positive, since $\rho > 0$).  Then
\[
\tr(\rho\,A^\dagger A) \;=\; \sum_i \lambda_i\, B(e_i,\, A^\dagger A\, e_i) \;=\; \sum_i \lambda_i\,\|A e_i\|_B^2 \;\ge\; 0.
\]
Therefore $\omega(A^*A) = \tr(\rho\,A^\dagger A)/\tr(\rho) \ge 0$.

\noindent\textit{Faithfulness.}  If $\omega(A^*A) = 0$, then $\sum_i \lambda_i \|Ae_i\|_B^2 = 0$.  Since every $\lambda_i > 0$, each summand vanishes: $\|Ae_i\|_B = 0$, hence $Ae_i = 0$ for every $i$.  Since $\{e_i\}$ is a basis for $V$: $A = 0$.  \qed

\smallskip\noindent\textit{Cost recovery.}  Since the sector projections are pairwise orthogonal ($E_{d'} E_d = \delta_{d' d}\,E_d$, $E_\Pi E_d = 0$), the product $\rho\,E_d = (\eps(d)/C)\,E_d$.  Taking traces:
\[
\omega(E_d) \;=\; \frac{\tr(\rho\,E_d)}{\tr(\rho)} \;=\; \frac{(\eps(d)/C)\,\dim(M_d)}{\tr(\rho)}.
\]
For irreducible distinctions ($\dim(M_d) = 1$), this gives $\omega(E_d) = (\eps(d)/C)/\tr(\rho)$, which is proportional to $\eps(d)$ with a sector-independent normalization factor.  The proportionality $\omega(E_d) \propto \eps(d)$ connects the GNS representation to the enforcement semantics: the ``probability'' of the $d$-sector projection is monotone in the capacity committed to maintaining~$d$.

\smallskip\noindent\textit{Uniqueness up to gauge.}  The cost-density construction fixes $\rho$ up to the following free data: the within-sector forms $B_d$ (which set the inner-product gauge on each $M_d$) and the pool weight $\delta$ (a single positive parameter).  These are normalization conventions; they do not affect the $*$-isomorphism class of $\mathcal{A}$, the Wedderburn block structure, the field selection, or the Born-rule output (Remark~\ref{rem:preHilbert}).  In this precise sense, the enforcement data determine $\omega$ up to gauge.
\end{proposition}

\begin{proposition}[O4$'$: Sector-Diagonal Forcing]\label{prop:O4prime}
Let $\omega: \mathcal{A} \to \mathbb{R}$ be a positive normalized linear functional satisfying:
\begin{enumerate}[nosep,label=(\alph*),leftmargin=2em]
\item \textit{Cost monotonicity:} $\omega(E_d) \ge \omega(E_{d'})$ whenever $\eps(d) \ge \eps(d')$.
\item \textit{Sector additivity:} $\omega(E_{d_1} + E_{d_2}) = \omega(E_{d_1}) + \omega(E_{d_2})$ when $E_{d_1} E_{d_2} = 0$.
\end{enumerate}
Then $\omega$ is represented by a density operator $\rho$ that is diagonal in the sector decomposition: $\rho = \sum_d \lambda_d\,E_d + \lambda_\Pi\, E_\Pi$ with all $\lambda_d, \lambda_\Pi > 0$.  The representing $\rho$ is unique up to within-sector normalization (the values $\lambda_d$ are free within the constraint $\sum_d \lambda_d\,\dim(M_d) + \lambda_\Pi\,\dim(\Pi) = 1$).

\smallskip\noindent\textit{Proof.}  By $L_\omega$.5, the sector projections $\{E_d\} \cup \{E_\Pi\}$ are pairwise $B$-orthogonal idempotents summing to $\mathrm{Id}_V$.  Any linear functional on $\mathcal{A}$ is determined by its values on a basis of $\mathcal{A}$; write $\omega(A) = \tr(\rho\,A)/\tr(\rho)$ for some self-adjoint $\rho$ (existence: Riesz representation on the finite-dimensional inner-product space $(\mathcal{A}, \langle\cdot,\cdot\rangle_B)$).

Decompose $\rho$ in the sector basis: $\rho = \sum_d \rho_d + \rho_\Pi + \rho_{\mathrm{cross}}$, where $\rho_d = E_d \rho E_d$, $\rho_\Pi = E_\Pi \rho E_\Pi$, and $\rho_{\mathrm{cross}}$ collects cross-sector terms.  For any $d \ne d'$: $\omega(E_d) = \tr(\rho\,E_d)/\tr(\rho)$.  By sector additivity~(b) and K3, the value $\omega(E_d)$ depends only on the $d$-block of $\rho$.  The cross-sector terms $E_{d'} \rho E_d$ contribute to $\tr(\rho\,A)$ only when $A$ has cross-sector components; but the sector projections $E_d$ are block-diagonal, so $\tr(\rho_{\mathrm{cross}}\,E_d) = 0$.  Since $\omega$ is determined on the generators $\{E_d, E_\Pi\}$ by the block-diagonal part alone, the representing $\rho$ can be taken block-diagonal without loss: $\rho = \sum_d \lambda_d E_d + \lambda_\Pi E_\Pi$.

Positivity of $\omega$ requires $\lambda_d \ge 0$ for all $d$ and $\lambda_\Pi \ge 0$.  Faithfulness ($\omega(A^*A) = 0 \Rightarrow A = 0$) requires strict positivity: $\lambda_d > 0$ and $\lambda_\Pi > 0$.  Normalization: $\omega(\mathrm{Id}) = 1$ gives the constraint $\sum_d \lambda_d\,\dim(M_d) + \lambda_\Pi\,\dim(\Pi) = \tr(\rho)$, and the division by $\tr(\rho)$ forces the sum to unity.  Cost monotonicity~(a) constrains the ratios $\lambda_d/\lambda_{d'}$ to respect the cost ordering $\eps(d) \ge \eps(d')$.  The construction in O4 (Construction~2) exhibits one such $\rho$; the remaining freedom is exactly the within-sector gauge.  \qed

\smallskip\noindent\textit{Why this matters.}  Proposition~O4$'$ shows that the cost-density state is not merely a convenient choice: it belongs to a \emph{constrained} equivalence class forced by the enforcement semantics.  Any state that respects cost monotonicity and sector additivity is sector-diagonal, faithful, and gauge-equivalent to Construction~2.  The APF data do not fix a unique state (the within-sector weights $\lambda_d$ remain free subject to the cost ordering), but they \emph{constrain} the representing operator to be sector-diagonal with cost-ordered weights --- a narrow equivalence class whose gauge freedom does not affect the downstream algebra.
\end{proposition}

\begin{theorem}[T2a: Wedderburn--Artin Decomposition]\label{thm:T2a}
$\mathcal{A} \cong \bigoplus_k M_{n_k}(\mathbb{F}_k)$ with $\mathbb{F}_k \in \{\mathbb{R}, \mathbb{C}, \mathbb{H}\}$.

\smallskip\noindent\textit{Hypothesis verification checklist.}  Wedderburn--Artin over $\mathbb{R}$ requires a finite-dimensional semisimple associative unital algebra.  We verify each property explicitly:

\begin{enumerate}[nosep,label=(W\arabic*),leftmargin=3em]
\item \textit{Associativity:} $\mathcal{A} \subseteq \End(V)$, and composition in $\End(V)$ is associative.  Inherited, not axiomatized (O3).
\item \textit{Unital:} $\mathrm{Id}_V = \sum_d E_d + E_\Pi \in \mathcal{A}$ (O3: the identity is a sum of the orthogonal sector projections).
\item \textit{Finite-dimensionality:} $\dim(\mathcal{A}) \le (\dim V)^2 \le (n_{\max}+1)^2 < \infty$ (O3, from OR3).
\item \textit{$*$-closure:} $\mathcal{A}$ is closed under $A \mapsto A^*$.  The $*$-involution is defined on generators ($E_d^* = E_d$ by $T_{\mathrm{adj}}$; $E_{\{d_1,d_2\}}^* = E_{\{d_1,d_2\}}$ by the same argument applied to the joint enforcement projector) and extended anti-multiplicatively to all of $\mathcal{A}$: $(AB)^* = B^*A^*$, $(\alpha A + \beta B)^* = \alpha A^* + \beta B^*$ (O3).  \textit{Compatibility with $B$-adjoint:} the O3 involution and the $L_\omega$ $B$-adjoint ($B(Tu,v) = B(u, T^\dagger v)$) agree on all generators (both give $E_d^\dagger = E_d$).  Since both are anti-multiplicative and agree on generators, they coincide on all of $\mathcal{A}$: $A^* = A^\dagger$ for every $A \in \mathcal{A}$.  This ensures the $*$-involution is the $B$-adjoint, not merely a formal symbol.
\item \textit{Faithful positive state:} $\omega(A) := \tr(\rho\,A)/\tr(\rho)$ with $\rho > 0$ (O4).  Positivity: $\omega(A^*A) \ge 0$.  Faithfulness: $\omega(A^*A) = 0 \Rightarrow A = 0$.  Both verified explicitly in O4.
\item \textit{Semisimplicity:} $\rad(\mathcal{A}) = \{0\}$.  Proved below using (W4) + (W5).
\end{enumerate}

\noindent\textit{Proof of semisimplicity.}  Let $N \in \rad(\mathcal{A})$ (the Jacobson radical: the largest nilpotent two-sided ideal).  Then $N^*N \in \rad(\mathcal{A})$ ($*$-stability of the radical: if $I$ is a two-sided ideal and $*$ is an involution, then $I^* = I$ for the Jacobson radical of a $*$-algebra; this uses (W4)).  $N^*N$ is positive semidefinite with respect to $B$: for any $v \in V$, $B(v, N^*N v) = B(Nv, Nv) = \|Nv\|_B^2 \ge 0$ (using $* = \dagger_B$ from (W4) and positive-definiteness of $B$ from $L_\omega$.4).  Since $N^*N \in \rad(\mathcal{A})$, it is nilpotent: $(N^*N)^k = 0$ for some $k$.  A positive semidefinite nilpotent operator on a finite-dimensional space has all eigenvalues simultaneously $\ge 0$ (PSD) and $= 0$ (nilpotent).  Therefore $N^*N = 0$ as an operator on $V$.  Then for any $A \in \mathcal{A}$: $\omega((NA)^*(NA)) = \omega(A^*N^*NA) = 0$ (since $N^*N = 0$).  By faithfulness of $\omega$ (W5): $NA = 0$ for all $A$.  In particular, $N = N \cdot \mathbf{1} = 0$.  Therefore $\rad(\mathcal{A}) = \{0\}$.

\smallskip\noindent\textit{Application of Wedderburn--Artin.}
By the Wedderburn--Artin theorem (Wedderburn 1908, Artin 1927): every finite-dimensional semisimple algebra over $\mathbb{R}$ decomposes as $\bigoplus_k M_{n_k}(\mathbb{F}_k)$ where each $\mathbb{F}_k$ is a division ring over $\mathbb{R}$.  By Frobenius' theorem (1878): $\mathbb{F}_k \in \{\mathbb{R}, \mathbb{C}, \mathbb{H}\}$.  \qed
\end{theorem}

\subsection{GNS construction ($T_{\mathrm{GNS}}$)}

\begin{theorem}[$T_{\mathrm{GNS}}$: Faithful Hilbert-Space Representation]\label{thm:GNS}
Given $\mathcal{A}$ (O3) and faithful positive $\omega$ (O4), there exists a faithful $*$-representation $\pi_\omega: \mathcal{A} \to B(\mathcal{H}_\omega)$ on a finite-dimensional Hilbert space.

\smallskip\noindent\textit{Hypotheses verified.}  The GNS construction requires: (i)~a unital $*$-algebra $\mathcal{A}$ (verified in O3 + T2a(W1--W4)); (ii)~a positive linear functional $\omega$ with $\omega(A^*A) \ge 0$ (O4); (iii)~faithfulness $\omega(A^*A) = 0 \Rightarrow A = 0$ (O4).

\smallskip\noindent\textit{Null-space elimination.}  The standard GNS construction quotients by the null space $\mathcal{N}_\omega := \{A \in \mathcal{A} : \omega(A^*A) = 0\}$ before defining the inner product.  When $\omega$ is faithful, $\mathcal{N}_\omega = \{0\}$ (by definition of faithfulness), so the quotient is trivial: $\mathcal{H}_\omega = \mathcal{A}/\mathcal{N}_\omega = \mathcal{A}$ as a vector space.  No completion is needed (finite-dimensional spaces are complete).  Faithfulness also makes the GNS representation injective: $\pi_\omega(A) = 0 \Rightarrow AB = 0$ for all $B \Rightarrow A = 0$ (taking $B = \mathbf{1}$).  This is the simplest case of the general GNS theory; the universal property --- that $(\pi_\omega, \mathcal{H}_\omega, \Omega_\omega)$ is the unique (up to unitary equivalence) cyclic representation recovering $\omega$ via $\omega(A) = \langle \Omega_\omega, \pi_\omega(A)\Omega_\omega\rangle$ --- holds by the standard argument (see Parzygnat~\cite{Parzygnat2018} for the 2-categorical formulation as a left adjoint).

\noindent\textit{Proof} (six steps).

\textit{G1 (Inner product form).}  Define $\langle A, B\rangle_\omega := \omega(A^*B)$ for $A, B \in \mathcal{A}$.

\textit{Bilinearity (real stage):} $\langle \alpha A_1 + \beta A_2, B\rangle_\omega = \omega((\alpha A_1 + \beta A_2)^* B) = \alpha\,\omega(A_1^*B) + \beta\,\omega(A_2^*B) = \alpha\langle A_1,B\rangle_\omega + \beta\langle A_2,B\rangle_\omega$ (using linearity of $\omega$ and $(\alpha A)^* = \alpha A^*$ for $\alpha \in \mathbb{R}$).  After T2c complexification, this becomes sesquilinearity in the standard complex sense.

\textit{Symmetry:} $\overline{\langle A,B\rangle_\omega} = \omega(A^*B) = \omega((B^*A)^*) = \omega(B^*A) = \langle B,A\rangle_\omega$ (using $\omega$ real-valued on self-adjoint elements and $(A^*B)^* = B^*A$).

\textit{Positive-definiteness:} $\langle A,A\rangle_\omega = \omega(A^*A) \ge 0$ (positivity of $\omega$), with equality iff $A = 0$ (faithfulness).

\textit{G2 (Hilbert space).}  $\mathcal{H}_\omega := (\mathcal{A}, \langle\cdot,\cdot\rangle_\omega)$.  Complete since $\dim(\mathcal{A}) < \infty$.  $\dim(\mathcal{H}_\omega) = \dim(\mathcal{A}) \le n_{\max}^2$.

\textit{G3 (Representation).}  $\pi_\omega(A)B := AB$ (left multiplication).  Linearity: $\pi_\omega(A)(\alpha B + \beta C) = \alpha AB + \beta AC$.  Multiplicativity: $\pi_\omega(AB)C = (AB)C = A(BC) = \pi_\omega(A)\pi_\omega(B)C$ (associativity).  $*$-preservation: $\langle \pi_\omega(A)B, C\rangle_\omega = \omega((AB)^*C) = \omega(B^*A^*C) = \langle B, A^*C\rangle_\omega = \langle B, \pi_\omega(A^*)C\rangle_\omega$ (using anti-multiplicativity $(AB)^* = B^*A^*$ and involutivity $A^{**} = A$), so $\pi_\omega(A)^\dagger = \pi_\omega(A^*)$.

\textit{G4 (Faithfulness).}  $\pi_\omega(A) = 0$ means $AB = 0$ for all $B$.  Taking $B = \mathbf{1}$: $A = 0$.

\textit{G5 (Cyclic vector).}  $\Omega_\omega := \mathbf{1} \in \mathcal{A}$.  Then $\pi_\omega(\mathcal{A})\Omega_\omega = \{A \cdot \mathbf{1} : A \in \mathcal{A}\} = \mathcal{A} = \mathcal{H}_\omega$ (cyclic).

\textit{G6 (Cost recovery).}  $\langle \Omega_\omega, \pi_\omega(A)\Omega_\omega\rangle_\omega = \langle \mathbf{1}, A\rangle_\omega = \omega(\mathbf{1}^*A) = \omega(A)$.  \qed

\smallskip\noindent\textit{Real-to-complex route.}  The GNS construction above is carried out entirely at the \emph{real} $*$-algebra level: $\mathcal{A}$ is a real algebra (O3), $\omega$ is a real-valued functional (O4), the inner product $\langle A,B\rangle_\omega = \omega(A^*B)$ is a real symmetric bilinear form (since $\omega$ is real-valued and $* = \dagger_B$), and $\mathcal{H}_\omega$ is a finite-dimensional real inner-product space.  No complex structure is assumed or used at this stage.  After T2c selects $\mathbb{F} = \mathbb{C}$ for every nonclassical block, the nonclassical summands of $\mathcal{A}$ are identified as complex matrix algebras $M_{n_k}(\mathbb{C})$, and the real representation $\pi_\omega$ on each such block is identified with a complex Hilbert-space representation in the standard sense (the complex structure is read off from the $J$ operator --- the antisymmetric generator whose existence was proved in T2c.1).  No completion is needed (the space is already finite-dimensional and complete).  The $C^*$-norm is then the operator norm on $B(\mathcal{H}_\omega)$, and the $C^*$-identity $\|A^*A\| = \|A\|^2$ is verified at the post-T2c complex level.  In summary: GNS is built over $\mathbb{R}$; complex structure is imposed by T2c; the final output is a standard complex $C^*$-algebra represented on a complex Hilbert space.

\smallskip\noindent\textit{Canonical choice of $\omega$.}  For the structural derivation, any faithful positive state suffices; the GNS output is independent of this choice up to unitary equivalence.  The cost-functional state $\omega(A) = \tr(\rho\,A)/\tr(\rho)$ (from O4, Construction~2) is physically preferred because it encodes the enforcement costs directly ($\omega(E_d) \propto \eps(d)$), but the algebraic results (T2a, T2c, $T_{\mathrm{Born}}$) hold for any faithful $\omega$.
\end{theorem}

\subsection{Field selection ($\mathbb{F} = \mathbb{C}$, T2c.1--T2c.2)}

\begin{lemma}[$\mathcal{D}$-Quotient $\Rightarrow$ Algebraic State-Separation]\label{lem:state_sep}
The $\mathcal{D}$-quotient on the physical substrate (SP, Definition~\ref{def:FD1}) extends to full state-separation of the enforcement-generated algebra $\mathcal{A}$.

\smallskip\noindent\textit{Statement.}  If $A \in \mathcal{A}$ is nonzero, then there exists a state $\rho$ with $\tr(\rho\,A) \ne 0$.

\smallskip\noindent\textit{Proof.}  Two cases, each reducing to the faithfulness of $\omega$ (O4, Proposition~\ref{prop:O4}).

\smallskip\noindent\textit{Case 1: $A$ self-adjoint.}  If $A = A^*$ and $A \ne 0$, then $A^2 = A^*A \ne 0$ (because $\mathcal{A}$ is semisimple (T2a, property~W6): the only nilpotent ideal is $\{0\}$, so a nonzero self-adjoint element cannot square to zero).  By faithfulness of $\omega$: $\omega(A^*A) = \omega(A^2) > 0$.  Since $\omega(\cdot) = \tr(\rho\,\cdot)/\tr(\rho)$ with $\rho > 0$ (Construction~2), this gives $\tr(\rho\,A^2) > 0$, and the state $\rho$ separates $A^2$.  To separate $A$ itself: since $A$ is self-adjoint on a finite-dimensional space, $A$ has real eigenvalues.  If $A \ne 0$, at least one eigenvalue $\lambda_j \ne 0$.  The rank-$1$ density $\rho_j = |e_j\rangle\langle e_j|$ (the corresponding eigenvector) gives $\tr(\rho_j\,A) = \lambda_j \ne 0$.

\smallskip\noindent\textit{Case 2: general $A$.}  Write $A = H + K$ with $H := (A + A^*)/2$ (self-adjoint part) and $K := (A - A^*)/2$ (anti-self-adjoint part).  If $A \ne 0$, at least one of $H, K$ is nonzero.  If $H \ne 0$: Case~1 applies to $H$, giving a state that separates $H$ and hence $A$.  If $H = 0$ and $K \ne 0$: then $K^*K = (-K)K = -K^2$, and faithfulness gives $\omega(K^*K) = -\omega(K^2) > 0$, so $K^2 \ne 0$.  Since $K^2$ is self-adjoint ($K^2$ has $(K^2)^* = K^{*2} = (-K)^2 = K^2$), Case~1 gives a state separating $K^2$, hence separating $K = A$.  \qed

\smallskip\noindent\textit{Primitive assumptions used.}  A1 (via enforcement costs in $\rho$), MD (via $L_{\eps^*}$ $\to$ $T_{\mathrm{form}}$ $\to$ OR3 $\to$ finite dimensionality), SP (the $\mathcal{D}$-quotient ensures every direction in $V$ is enforcement-visible), K3 (via $L_\omega$ $\to$ $B$ $\to$ $T_{\mathrm{adj}}$ $\to$ $*$-involution).  T2a(W6) enters for the semisimplicity step in Case~1.

\smallskip\noindent\textit{Why this lemma matters.}  The $\mathcal{D}$-quotient is defined on primitive distinctions: ``no direction in $S_\Gamma$ is invisible to enforcement status.''  T2c.1 (Theorem~\ref{thm:T2c1}) needs state-separation of the \emph{generated algebra}, which includes products and linear combinations of enforcement generators --- elements that may not themselves be enforcement projections.  This lemma bridges from the operational quotient to the algebraic condition.
\end{lemma}

\begin{lemma}[Antisymmetric Sector Generation in Simple Blocks]\label{lem:antisym}
Let $\mathcal{B} \cong M_n(\mathbb{R})$ be a simple real matrix block with $n \ge 2$.  If $\mathcal{B}$ contains one nonzero commutator $[X,Y] \ne 0$ with $X, Y$ self-adjoint, then $A := [X,Y]$ is nonzero and antisymmetric ($A^\top = -A$).  Closure under products and linear combinations generates the full antisymmetric subspace $\mathfrak{so}(n)$, which has dimension $n(n-1)/2$.

\smallskip\noindent\textit{Proof.}
$A^\top = (XY - YX)^\top = Y^\top X^\top - X^\top Y^\top = YX - XY = -A$.  By simplicity of $M_n(\mathbb{R})$, the two-sided ideal generated by $A$ is the whole algebra, hence all of $\mathfrak{so}(n)$.  \qed
\end{lemma}

\begin{definition}[State-Separation Completeness]\label{def:SSC}
A simple block $\mathcal{B} \subseteq \mathcal{A}$ is \emph{state-separation complete} if for every nonzero $A \in \mathcal{B}$, there exists a physical state $\rho$ (positive, trace-$1$, self-adjoint with respect to the $*$-involution on $\mathcal{B}$) such that $\tr(\rho\,A) \ne 0$.

\smallskip\noindent\textit{APF status.}  State-separation completeness is the algebra-level content of the $\mathcal{D}$-quotient (SP): every nonzero direction in $V$ is enforcement-visible, so every nonzero algebra element is detectable by some physical state.  It is proved as Lemma~\ref{lem:state_sep} from SP + O4 + T2a(W6).  It is a single-interface condition; no composition or multi-system structure is needed.
\end{definition}

\begin{center}\small
\renewcommand{\arraystretch}{1.15}
\begin{tabular}{@{}p{3.2cm}ccp{4.8cm}@{}}
\toprule
\textbf{Condition} & \textbf{Single-system?} & \textbf{Excludes} & \textbf{Operational content} \\
\midrule
SSC (Def.~\ref{def:SSC}) & Yes & $\mathbb{R}$ blocks & Every nonzero algebra element is detectable by some physical state \\
Local tomography & No (composite) & $\mathbb{R}$ and $\mathbb{H}$ & Composite states determined by local measurements \\
Composition closure (T2c.2) & No (composite) & $\mathbb{H}$ blocks & Field class preserved under independent composition \\
\bottomrule
\end{tabular}
\end{center}
\smallskip\noindent\textit{Comparison.}  SSC is the APF's single-interface replacement for local tomography's $\mathbb{R}$-exclusion.  Local tomography excludes both $\mathbb{R}$ and $\mathbb{H}$ but requires composite-system structure; SSC excludes only $\mathbb{R}$ but requires only single-interface data.  The APF's novelty is not eliminating the need for a composite assumption (T2c.2 still requires composition closure for $\mathbb{H}$-exclusion), but relocating the $\mathbb{R}$-exclusion to a single-interface algebraic condition derived from the $\mathcal{D}$-quotient.

\begin{theorem}[T2c.1: Exclusion of Real Nonclassical Blocks]\label{thm:T2c1}
\textit{Status: \textbf{proved} (single-interface).}  No nonclassical simple block of $\mathcal{A}$ is of real type.  That is: if $\mathcal{B} \cong M_n(\mathbb{R})$ with $n \ge 2$, then $\mathcal{B}$ is not state-separation complete.

\smallskip\noindent\textbf{Uses:} T2a (Wedderburn trichotomy), $T_{\mathrm{alg}}$ (noncommutativity), Lemma~\ref{lem:state_sep} (state-separation completeness from SP + O4 + T2a(W6)).

\smallskip\noindent\textit{(The case $n = 1$: $M_1(\mathbb{R}) \cong \mathbb{R}$ is a commutative block with no nonzero commutators, no antisymmetric elements, and no obstruction from state-separation completeness.  Classical blocks are field-agnostic; the exclusion below targets only the nonclassical case $n \ge 2$.)}

\smallskip\noindent\textit{Proof.}
Let $\mathcal{B} \cong M_n(\mathbb{R})$ with $n \ge 2$.  Because $\mathcal{B}$ is noncommutative ($T_{\mathrm{alg}}$), there exist self-adjoint enforcement-generated elements $X, Y \in \mathcal{B}$ with $[X, Y] \ne 0$.  Set $A := [X, Y]$.  Since $X^\top = X$ and $Y^\top = Y$:
\[
A^\top = [X,Y]^\top = Y^\top X^\top - X^\top Y^\top = YX - XY = -[X,Y] = -A.
\]
So $A$ is nonzero and antisymmetric.  Since $A$ is built from products and linear combinations of enforcement generators, it is enforcement-generated and belongs to $\mathcal{A}$.

Now let $\rho$ be any physical state on $\mathcal{B}$.  Over $\mathbb{R}$, the $*$-involution is the transpose, so self-adjoint states are symmetric: $\rho^\top = \rho$.  Therefore:
\[
\tr(\rho\,A) = \tr((\rho A)^\top) = \tr(A^\top \rho^\top) = \tr((-A)\rho) = -\tr(\rho\,A),
\]
which gives $\tr(\rho\,A) = 0$ for every physical state $\rho$.  Thus $A \ne 0$ is state-invisible: $\mathcal{B}$ fails state-separation completeness (Definition~\ref{def:SSC}).

By Lemma~\ref{lem:state_sep}, the enforcement-generated algebra $\mathcal{A}$ \emph{is} state-separation complete.  A block of $\mathcal{A}$ that fails this property cannot be admissible.  Contradiction.  Therefore no nonclassical simple block can be of real type.  \qed

\smallskip\noindent\textit{The mechanism.}  Over $\mathbb{R}$, the self-adjoint (= symmetric) part of $M_n(\mathbb{R})$ has dimension $n(n+1)/2$, but the full algebra has dimension $n^2$.  The antisymmetric subspace $\mathfrak{so}(n)$, of dimension $n(n-1)/2$, is entirely invisible to symmetric states.  This is the failure of state-separation completeness (Definition~\ref{def:SSC}) for real blocks: the state space cannot see the full algebra.
\end{theorem}

\begin{proposition}[Complexification Repairs the Defect]\label{prop:C_repair}
For every nonzero antisymmetric $A \in M_n(\mathbb{R})$, the element $iA \in M_n(\mathbb{C})$ is Hermitian, and there exists a Hermitian density matrix $\rho$ with $\tr(\rho\,iA) \ne 0$.

\smallskip\noindent\textit{Proof.}  $(iA)^\dagger = -i A^\top = -i(-A) = iA$: Hermitian.  Since $iA \ne 0$ is Hermitian, it has at least one nonzero real eigenvalue $\lambda_j$.  The rank-$1$ density $\rho_j = |e_j\rangle\langle e_j|$ gives $\tr(\rho_j\, iA) = \lambda_j \ne 0$.  \qed

\smallskip\noindent\textit{Consequence.}  Over $\mathbb{C}$, the self-adjoint (Hermitian) part of $M_n(\mathbb{C})$ has real dimension $n^2 = n(n+1)/2 + n(n-1)/2$: it includes both the symmetric sector and (via $iA$) the antisymmetric sector.  Every enforcement-generated element is now state-separable.  Complexification is the minimal field extension that restores state-separation completeness.
\end{proposition}

\begin{theorem}[T2c.2: Exclusion of Quaternionic Blocks Under Composition Closure]\label{thm:T2c2}
\textit{Status: \textbf{conditional} on composition closure.}  If the admissible block class is closed under independent composition of subsystems, then no nonclassical simple block of $\mathcal{A}$ is of quaternionic type.

\smallskip\noindent\textbf{Uses:} T2a (Wedderburn trichotomy), $L_{\mathrm{loc}}$ (capacity additivity for independent interfaces, from A1 + FD2), Brauer identity $\mathbb{H} \otimes_\mathbb{R} \mathbb{H} \cong M_4(\mathbb{R})$ (standard algebra).

\smallskip\noindent\textit{This step uses an additional composition principle beyond single-interface capacity.  It is not derived from single-interface enforcement alone.  A reader who accepts only single-interface results has $\mathbb{R}$-exclusion (T2c.1) but not $\mathbb{H}$-exclusion; the latter genuinely needs composition.}

\smallskip\noindent\textit{Proof.}
Suppose quaternionic blocks were allowed.  Let $\mathcal{B}_1 \cong M_m(\mathbb{H})$ and $\mathcal{B}_2 \cong M_n(\mathbb{H})$ model two independent subsystems.  Their real independent composition is:
\[
\mathcal{B}_1 \otimes_\mathbb{R} \mathcal{B}_2 \;\cong\; M_m(\mathbb{H}) \otimes_\mathbb{R} M_n(\mathbb{H}) \;\cong\; M_{mn}(\mathbb{H} \otimes_\mathbb{R} \mathbb{H}) \;\cong\; M_{mn}(M_4(\mathbb{R})) \;\cong\; M_{4mn}(\mathbb{R}),
\]
using the Brauer identity $\mathbb{H} \otimes_\mathbb{R} \mathbb{H} \cong M_4(\mathbb{R})$.  The composite is a real matrix algebra, not a quaternionic one.  The class of quaternionic blocks is not closed under independent composition, contradicting the composition-closure premise.

Therefore no nonclassical simple block can be of quaternionic type.  \qed

\smallskip\noindent\textit{Remark: complex blocks and composition.}  By contrast, complex blocks are compatible with composition: $M_m(\mathbb{C}) \otimes_\mathbb{R} M_n(\mathbb{C}) \cong M_{mn}(\mathbb{C}) \oplus M_{mn}(\mathbb{C})$.  The doubled copy is an artifact of real tensor product; the physical sector (selected by $L_{\mathrm{loc}}$'s budget constraint, which provides no budget for the excess copy) is $M_{mn}(\mathbb{C})$ with no field-class change.  The formal derivation of this sector selection requires the tensor-product identification ($T_{\mathrm{tensor}}$, conditional on T2c), and is not claimed as a theorem of the present supplement.  It is noted here only to confirm that $\mathbb{C}$ is compatible with the composition-closure premise, not to derive that compatibility.
\end{theorem}

\begin{corollary}[T2c: Complex Field Selection]\label{cor:T2c}
Under T2c.1 and T2c.2, every nonclassical simple block of $\mathcal{A}$ is complex:
\[
\mathcal{A} \;\cong\; \bigoplus_k M_{n_k}(\mathbb{C}) \;\oplus\; \bigoplus_j \mathbb{R}_j,
\]
where the $\mathbb{R}_j$ summands are classical ($n = 1$) blocks.  The $*$-involution on each complex block is the Hermitian adjoint; the $C^*$-norm is $\|A\|_{\mathrm{op}}$ on $B(\mathcal{H}_\omega)$ (from $T_{\mathrm{GNS}}$).

\smallskip\noindent\textbf{Uses:} T2c.1 (proved, single-interface), T2c.2 (conditional on composition closure).
\end{corollary}

\smallskip\noindent\textit{Comparison with standard reconstruction routes.}  Unlike local-tomography reconstructions (Hardy 2001, CDP 2011, Masanes--M\"uller 2011), the APF excludes $\mathbb{R}$ through state-separation completeness of enforcement-generated antisymmetric elements --- a single-interface algebraic condition, not a multi-system tomographic one.  Unlike Sol\'er-type routes, the APF does not require infinite orthonormal sequences; the finite-dimensional Wedderburn trichotomy is sufficient.  $\mathbb{H}$-exclusion still requires a composition principle, as do all known reconstructions: Hardy uses simplicity + continuous reversibility on composites, CDP uses purification (an inherently composite axiom), and Masanes--M\"uller use local tomography directly.  The APF's composition-closure condition (T2c.2) is comparable in strength to these alternatives.

\begin{tcolorbox}[apfbox, title={Field-Selection Status}]
\begin{itemize}[nosep,leftmargin=1.5em]
\item \textbf{T2c.1 (single-interface, proved):} excludes nonclassical real blocks via state-separation completeness.  No composition or multi-system structure needed.
\item \textbf{T2c.2 (composition closure, conditional):} excludes quaternionic blocks under the assumption that field class is preserved under independent composition.  Comparable in strength to the composite-system assumptions used elsewhere to exclude $\mathbb{H}$.
\item \textbf{Full T2c:} all nonclassical blocks are complex only when both T2c.1 and T2c.2 are accepted.
\end{itemize}
\end{tcolorbox}

\smallskip\noindent\textit{Conditionality summary.}  T2c.1 is earned from the existing single-interface framework.  T2c.2 is a genuine additional premise about multi-interface composition; it is physically motivated by $L_{\mathrm{loc}}$'s budget additivity but requires the structural claim that field class is preserved under independent composition.  Neither exclusion presupposes $\mathbb{C}$; the complex field is the conclusion.

\smallskip\noindent\textit{What survives at each tier.}  If only T2c.1 is accepted (single-interface): real nonclassical blocks are excluded, but quaternionic blocks remain open.  The algebra is $\mathcal{A} \cong \bigoplus_k M_{n_k}(\mathbb{F}_k)$ with $\mathbb{F}_k \in \{\mathbb{C}, \mathbb{H}\}$ for nonclassical blocks.  The Born rule still follows per block (Busch's theorem applies over both $\mathbb{C}$ and $\mathbb{H}$; quaternionic Hilbert spaces admit Gleason-type theorems).  The Tsirelson bound $\le 2\sqrt{2}$ holds (the $*$-algebraic derivation is field-independent); saturation at $2\sqrt{2}$ requires $\mathbb{C}$.  Monogamy and no-broadcasting Route~1 are unaffected (they use only the resource chain).  If full T2c (T2c.1 + T2c.2) is accepted: all nonclassical blocks are complex, and the full standard quantum structure is recovered.

\smallskip\noindent\textit{Alternative routes via Jordan and $C^*$-algebra classification.}  The proof above is deliberately elementary: it uses only the Wedderburn--Artin trichotomy, the transpose symmetry of real states, and the Brauer identity for quaternions.  Three more sophisticated routes to the same conclusion exist in the literature and could replace or supplement the argument:

\begin{itemize}[nosep,leftmargin=2em,itemsep=3pt]
\item \textit{Jordan--von Neumann--Wigner classification.}  The self-adjoint part $\mathcal{A}_{\mathrm{sa}} := \{A \in \mathcal{A} : A^* = A\}$ is a finite-dimensional formally real Jordan algebra under $A \circ B := \frac{1}{2}(AB + BA)$.  The JvNW theorem classifies such algebras: each simple summand is isomorphic to $M_n(\mathbb{F})_{\mathrm{sa}}$ for $\mathbb{F} \in \{\mathbb{R}, \mathbb{C}, \mathbb{H}\}$ (with $n \ge 2$), the spin factors $V_n$ ($n \ge 3$), or the exceptional Albert algebra $M_3(\mathbb{O})_{\mathrm{sa}}$.  Our $\mathcal{A}$ is associative (O3: $\mathcal{A} \subseteq \End(V)$), which excludes the exceptional octonion factor.  The spin factors are also excluded by associativity for $n \ge 4$; for $n = 3$, the spin factor $V_3 \cong M_2(\mathbb{C})_{\mathrm{sa}}$ is already complex.  This recovers the Wedderburn trichotomy from the Jordan side, after which T2c.1 and T2c.2 apply as above.
\item \textit{Sol\`er-type arguments.}  Sol\`er's theorem (1995) shows that an infinite-dimensional orthomodular space over a $*$-field with an infinite orthonormal sequence is isomorphic to a Hilbert space over $\mathbb{R}$, $\mathbb{C}$, or $\mathbb{H}$.  In the finite-dimensional setting (our case), the trichotomy is already given by Frobenius/Wedderburn; Sol\`er's theorem would become relevant in the $C \to \infty$ limit (inductive-limit algebras).
\item \textit{Effectus/order-unit approach.}  Van de Wetering~\cite{vandeWetering2019} derives complex quantum theory from effect-algebraic axioms.  The APF's positive cone $\mathcal{A}^+$ (defined in O3) and the faithful state $\omega$ (O4) give $\mathcal{A}_{\mathrm{sa}}$ the structure of an order-unit space with a separating family of states.  Under this identification, state-separation completeness corresponds to the ``no restriction'' axiom (effects separate states) and composition closure (T2c.2) corresponds to ``local tomography'' (composite states are determined by local measurements).  The effectus route would derive the complex field from these conditions via a different proof strategy (quotient completion and spectral duality), reaching the same endpoint.
\end{itemize}

\noindent The present proof is chosen for minimality: it requires no Jordan theory, no Sol\`er machinery, and no effectus formalism --- only Wedderburn--Artin (already proved in T2a), the symmetry of real states (a two-line computation), and the Brauer identity (a standard algebraic fact).  The alternative routes are noted here for the reader who prefers a different proof framework; they all reach the same conclusion under comparable assumptions.

% ═══════════════════════════════════════════════════════════════
\section{Consequences: Precise Status}
\label{sec:consequences}

\begin{lemma}[$L_{\mathrm{cert}}$: Sharp Certification Probabilities from Enforcement Weight]\label{lem:Lcert}
For each sharp enforcement projection $E_d$ (FD5: the self-adjoint idempotent certifying distinction $d$), the \emph{operational certification probability} in an admissible state represented by $\omega$ is:
\[
p_{\mathrm{cert}}(E_d) \;:=\; \omega(E_d).
\]

\smallskip\noindent\textit{Operational content.}  A sharp certification test for $d$ at interface $\Gamma$ asks: ``Is $d$ actively enforced?''  The outcome is binary (yes/no).  In the APF, the enforcement projection $E_d$ is the mathematical representative of this test (FD5).  The cost-density state $\omega$ from O4 (Construction~2) assigns $\omega(E_d) \propto \eps(d)$: the probability of successful certification is proportional to the enforcement capacity committed to maintaining $d$.  This identification is constitutive --- it defines what ``operational probability of certification'' means within A1's budget semantics --- and does not import an external measurement postulate.

\smallskip\noindent\textit{Verification of probability axioms on sharp projections.}
\begin{enumerate}[nosep,label=(\roman*),leftmargin=2em]
\item \textit{Range:} $p_{\mathrm{cert}}(E_d) = \omega(E_d) = \omega(E_d^* E_d) \ge 0$ (positivity of $\omega$); $\omega(\mathrm{Id} - E_d) \ge 0$ gives $p_{\mathrm{cert}}(E_d) \le 1$.
\item \textit{Normalization:} $p_{\mathrm{cert}}(\mathrm{Id}) = \omega(\mathrm{Id}) = 1$ (O4).
\item \textit{Ortho-additivity:} for orthogonal sector projections $E_{d_1} E_{d_2} = 0$: $p_{\mathrm{cert}}(E_{d_1} + E_{d_2}) = \omega(E_{d_1}) + \omega(E_{d_2})$ by linearity.  Physically: certifying ``$d_1$ or $d_2$'' on disjoint sectors costs the sum of the individual commitments (K3).
\item \textit{Non-contextuality:} $\omega(E_d)$ depends only on the projection $E_d$, not on any resolution of $\mathrm{Id}$ containing $E_d$, because $\omega$ is a linear functional (O4).
\end{enumerate}
\textit{Conditions used:} A1 (enforcement costs), FD5 (sharp certification tests), O4 (faithful positive state), K3 (disjoint-support additivity for the physical grounding of ortho-additivity on sector generators).  \qed
\end{lemma}

\begin{lemma}[$L_{\mathrm{prob}}$: Extension to All Effects]\label{lem:Lprob}
The sharp certification probability $p_{\mathrm{cert}}$ (Lemma~\ref{lem:Lcert}) extends uniquely to all effects $E \in \mathcal{A}$ with $0 \le E \le \mathrm{Id}$:
\[
p(E) \;:=\; \omega(E).
\]
The extension preserves:
\begin{enumerate}[nosep,label=(\alph*),leftmargin=2em]
\item \textit{Range:} $p(E) \in [0,1]$ for every effect $E$.
\item \textit{Normalization:} $p(\mathrm{Id}) = 1$.
\item \textit{Ortho-additivity:} if $E_1 E_2 = 0$ and $E_1 + E_2 \le \mathrm{Id}$, then $p(E_1 + E_2) = p(E_1) + p(E_2)$.
\item \textit{Affine under mixing:} $p(\lambda E_1 + (1-\lambda)E_2) = \lambda\,p(E_1) + (1-\lambda)\,p(E_2)$ for $\lambda \in [0,1]$.
\item \textit{Non-contextuality:} $p(E)$ depends only on the algebra element $E$, not on any decomposition or context in which $E$ appears.
\end{enumerate}

\smallskip\noindent\textit{Why the extension is unique.}  In finite dimensions, every effect has a spectral decomposition into projections: $E = \sum_i \lambda_i P_i$ with $0 \le \lambda_i \le 1$ and $P_i$ pairwise orthogonal projections.  Any additive, normalized probability map on projections extends uniquely by linearity to effects: $p(E) = \sum_i \lambda_i\, p(P_i)$.  Since $\omega$ is linear on all of $\mathcal{A}$ (O4), it automatically provides this extension: $\omega(E) = \sum_i \lambda_i\,\omega(P_i)$.

\smallskip\noindent\textit{Proof of (a)--(e).}  (a)~$E = (\sqrt{E})^*(\sqrt{E})$ gives $\omega(E) \ge 0$; $\mathrm{Id} - E \ge 0$ gives $\omega(E) \le 1$.  (b)~O4, normalization.  (c)--(d)~Linearity of $\omega$.  (e)~$\omega$ assigns a single value to each algebra element; no additional non-contextuality axiom is introduced beyond the semantic identification of effects as algebra elements.  \qed
\end{lemma}

\begin{theorem}[$T_{\mathrm{Born}}$: Born Rule]\label{thm:Born}
\textit{Status: \textbf{proved}.}  For every effect $E \in \mathcal{A}$ with $0 \le E \le \mathrm{Id}$: $p(E) = \tr(\rho\,E)$ for a unique density operator $\rho$.

\smallskip\noindent\textbf{Uses:} T2c (Corollary~\ref{cor:T2c}), $T_{\mathrm{GNS}}$ (Theorem~\ref{thm:GNS}), O4 (Proposition~\ref{prop:O4}), $L_{\mathrm{cert}}$ (Lemma~\ref{lem:Lcert}), $L_{\mathrm{prob}}$ (Lemma~\ref{lem:Lprob}), Busch 2003~\cite{Busch2003}.

\smallskip\noindent\textit{Proof.}  We apply Busch's generalized Gleason theorem~\cite{Busch2003} for normalized finitely additive measures on the effect algebra of a complex Hilbert space of dimension $\ge 2$.  The theorem states: \textit{every function $p$ on the effect algebra $\mathcal{E}(\mathcal{H})$ that is (i)~non-contextual, (ii)~additive on orthogonal effects, and (iii)~normalized ($p(\mathrm{Id}) = 1$) is given by $p(E) = \tr(\rho\,E)$ for a unique density operator $\rho$.}

We invoke Busch on effects rather than projection-only Gleason because APF operational tests naturally produce effects (via OR1's convex mixing of enforcement outcomes), and sharp projections are the special case.

\smallskip\noindent We verify the hypotheses using $L_{\mathrm{prob}}$ (Lemma~\ref{lem:Lprob}).

\smallskip\noindent(i)~\textit{Complex Hilbert space, dimension $\ge 2$.}  After T2c's field selection ($\mathbb{F} = \mathbb{C}$), the GNS Hilbert space $\mathcal{H}_\omega$ ($T_{\mathrm{GNS}}$, Theorem~\ref{thm:GNS}) is a finite-dimensional complex Hilbert space with $\mathcal{A}$ faithfully represented in $B(\mathcal{H}_\omega)$.  By $T_{\mathrm{form}}$, $n_{\max} = \lfloor C/\eps^*\rfloor \ge 3$ at any non-trivial interface ($C \ge 3\eps^*$, required for $L_\Delta$), so $\dim(\mathcal{H}) \ge 2$.  \textit{Note:} Busch's generalization extends Gleason from $\dim \ge 3$ (projections only) to $\dim \ge 2$ (full effect algebra).  This extension applies here because $p$ is defined on all effects via $L_{\mathrm{prob}}$, not just on the generators $\{E_d\}$.

\smallskip\noindent(ii)~\textit{Non-contextuality.}  $L_{\mathrm{prob}}$(e): $p(E)$ depends only on the effect $E$, not on any decomposition or context.

\smallskip\noindent(iii)~\textit{Ortho-additivity on effects.}  $L_{\mathrm{prob}}$(c): $p(E_1 + E_2) = p(E_1) + p(E_2)$ when $E_1 E_2 = 0$ and $E_1 + E_2 \le \mathrm{Id}$.

\smallskip\noindent(iv)~\textit{Normalization.}  $L_{\mathrm{prob}}$(b): $p(\mathrm{Id}) = 1$.

\smallskip\noindent All hypotheses verified.  By Busch's theorem: there exists a unique density operator $\rho_\omega$ on $\mathcal{H}_\omega$ such that
\[
p(E) \;=\; \tr(\rho_\omega\, E)
\]
for every effect $E \in \mathcal{E}(\mathcal{H}_\omega)$.  On projections: $p(P) = \tr(\rho_\omega\,P)$.  This is the Born rule.  \qed

\smallskip\noindent\textit{What is proved vs.\ what is invoked.}  The Born rule is not postulated; it is a \emph{theorem} of the derived algebra.  The work done by the APF is: constructing the real $*$-algebra (anti-circularity audit, eight steps, from A1 + MD + BW + constitutive semantics), selecting the complex field (T2c.1 from state-separation completeness; T2c.2 from composition closure), constructing the faithful positive state $\omega$ (O4), defining the operational probability $p$ on the full effect algebra ($L_{\mathrm{cert}}$ + $L_{\mathrm{prob}}$), and verifying Busch's hypotheses.  The external mathematical input is Busch's theorem; the physical content is derived from the stated APF inputs including the field-selection premises.  APF does not claim to reprove Busch internally; the contribution is constructing the arena in which Busch applies and verifying that its hypotheses are met.
\end{theorem}

\begin{theorem}[$T_{\mathrm{Tsirelson}}$: CHSH $\le 2\sqrt{2}$]\label{thm:Tsirelson}
\textit{Status: bound \textbf{proved} (pre-Hilbert); saturation \textbf{conditional} on $\mathbb{F} = \mathbb{C}$.}

\smallskip\noindent\textit{Hypotheses used (minimal list).}
\begin{enumerate}[nosep,label=(T\arabic*),leftmargin=2.5em]
\item \textit{Associative $*$-algebra:} $\mathcal{A}$ from T2a (finite-dimensional, semisimple, real).  \textbf{No Hilbert space or tensor product required for the bound.}
\item \textit{Dichotomic observables:} $a_1, a_2, b_1, b_2 \in \mathcal{A}$ with $a_i^2 = b_j^2 = I$ and $a_i^* = a_i$, $b_j^* = b_j$ (self-adjoint involutions).
\item \textit{Local commutativity:} $[a_i, b_j] = 0$ for all $i, j$.  This is the independence assumption.  In the APF, it follows from $L_{\mathrm{loc}}$ (Theorem~\ref{thm:Lloc}): when $a_i$ acts at $\Gamma_A$ and $b_j$ acts at $\Gamma_B$ with disjoint substrates ($S_{\Gamma_A} \cap S_{\Gamma_B} = \emptyset$), the enforcement operations have disjoint supports, so $a_i b_j = b_j a_i$ (non-overlapping operators commute in $\End(V)$).  \textbf{This does not require a tensor product structure} --- it requires only that Alice's and Bob's operators act on non-overlapping subspaces of a single $V$.
\item \textit{Operator norm:} $\|\cdot\|_{\mathrm{op}}$ on $\End(V)$, inherited from the $B$-norm on $V$ ($L_\omega$).
\end{enumerate}

\smallskip\noindent\textit{What is NOT required for the bound.}  No tensor product $\mathcal{H}_A \otimes \mathcal{H}_B$.  No no-signaling condition (it is a consequence, not an input).  No Hilbert space (the bound is a $*$-algebraic identity).  No complex field (the bound holds over $\mathbb{R}$; saturation requires $\mathbb{C}$).

\smallskip\noindent\textit{Proof.}  Define the CHSH operator $\hat{S} := a_1(b_1 + b_2) + a_2(b_1 - b_2)$.  The Cirel'son identity gives $\hat{S}^2 = 4I - [a_1,a_2][b_1,b_2]$ in any associative algebra (expand and use (T2)--(T3)).  

\textit{Norm bound:}  $a_i^2 = I$ gives $\|a_i\|_{\op} = 1$; therefore $\|[a_1,a_2]\|_{\op} \le \|a_1\|\|a_2\| + \|a_2\|\|a_1\| = 2$, and likewise $\|[b_1,b_2]\|_{\op} \le 2$.  Therefore $\|\hat{S}^2\|_{\op} \le 4 + 4 = 8$, giving $\|\hat{S}\|_{\op} \le 2\sqrt{2}$.

\textit{Saturation} ($\|\hat{S}\| = 2\sqrt{2}$) requires maximally entangled states, which exist over $\mathbb{C}$ (T2c): in $M_2(\mathbb{C}) \otimes M_2(\mathbb{C})$, the Bell state $|\Phi^+\rangle$ gives $\langle \hat{S} \rangle = 2\sqrt{2}$.  The present theorem proves the bound; exact saturation is stated only after full complex field selection (T2c).  Therefore saturation is conditional on T2c.  \qed

\smallskip\noindent\textit{Relation to tensor products.}  The bound $2\sqrt{2}$ is proved in the single-algebra setting ($\mathcal{A} \subseteq \End(V)$, with Alice's and Bob's operators as commuting subalgebras).  Tensor-product factorization $\mathcal{A} = \mathcal{A}_A \otimes \mathcal{A}_B$ is a \emph{stronger} structural claim that implies (T3) but is not required for it.  The APF derives the tensor factorization only conditionally ($T_{\mathrm{tensor}}$, which requires T2c + T2c.2's compositional closure).  A reader who accepts only the single-algebra setting has the $2\sqrt{2}$ bound but not the tensor-product structure.  The bound itself is robust: it holds in any $C^*$-algebra with commuting dichotomic observables, whether or not the algebra factorizes as a tensor product.
\end{theorem}

\begin{theorem}[$T_M$: Interface Monogamy]\label{thm:TM}
\textit{Status: \textbf{proved unconditionally}.}  For two correlated targets: $\Delta_{AB} + \Delta_{AC} \le C(\Gamma_A)$.  For $n$ independent targets: $\sum_{i=1}^n \Delta_{A B_i} \le C(\Gamma_A)$.

\smallskip\noindent\textit{Proof (two targets).}  Let $A$ share an interface with $B$ and with $C$.  Define $\Delta_{AB} := \eps(\{d_A, d_B\}) - \eps(d_A) - \eps(d_B) \ge \eps^*$ (by $L_\Delta$, since $d_A$ and $d_B$ are co-located at $\Gamma_A$).  Similarly $\Delta_{AC} \ge \eps^*$.  Both surpluses represent capacity committed at $\Gamma_A$: the surplus $\Delta_{AB}$ is the cost of defending $d_A$ and $d_B$ against pool-supported perturbations at $\Gamma_A$, and $\Delta_{AC}$ is the analogous cost for the $A$--$C$ pair.  By A1, the total enforcement cost at $\Gamma_A$ is bounded: $\eps(\{d_A, d_B\}) + \eps(\{d_A, d_C\}) - \eps(d_A) \le C(\Gamma_A)$ (the shared defense of $d_A$ is counted once).  Rearranging: $\Delta_{AB} + \Delta_{AC} \le C(\Gamma_A) - \eps(d_A) - \eps(d_B) - \eps(d_C) + \eps(d_A) \le C(\Gamma_A)$.

\smallskip\noindent\textit{Proof ($n$ targets).}  Let $A$ share interfaces with independent targets $B_1, \ldots, B_n$.  Each correlation draws surplus $\Delta_{AB_i} \ge \eps^*$ from $\Gamma_A$ (by $L_\Delta$: each $d_{B_i}$ is nonclassically co-located with $d_A$ at $\Gamma_A$).  The surpluses are committed from disjoint enforcement tasks at $\Gamma_A$ (the pool-mediated defenses for different targets operate on independent substrate DOF at the respective target interfaces, but all draw from $\Gamma_A$'s budget).  By A1: $\sum_{i=1}^n \Delta_{AB_i} \le C(\Gamma_A)$.  \qed
\end{theorem}

\begin{theorem}[$T_{\mathrm{NB}}$: Finite Broadcasting Number and No-Broadcasting]\label{thm:TNB}
\textit{Status: Route~1 \textbf{proved unconditionally} (quantitative); Route~2 \textbf{conditional} on $T_{\mathrm{tensor}}$ (qualitative).}

\textit{Definition.}  A \emph{broadcasting map} for a nonclassical distinction $d$ is an admissible map that produces $n$ independent faithful copies of $d$ at target interfaces $\Gamma_1, \ldots, \Gamma_n$, each of which can be independently measured to determine whether $d$ holds.

\medskip\noindent\textbf{Route 1: Quantitative broadcasting bound (unconditional).}

\smallskip\noindent\textit{Theorem.}  The maximum number of faithful copies of a nonclassical distinction from a single source interface is bounded:
\[
n \;\le\; n_{\mathrm{bc}} \;:=\; \left\lfloor \frac{C(\Gamma_{\mathrm{src}})}{\Delta_0} \right\rfloor,
\]
where $\Delta_0 := \inf_{d' \text{ co-located}} \Delta(d, d') \ge \eps^* > 0$ is the minimum superadditivity surplus at the source.

\smallskip\noindent\textit{Proof} (four steps).

\textit{Step 1: Each copy generates surplus at the source.}  A faithful copy of nonclassical distinction $d$ at target $\Gamma_i$ reproduces the full nonclassical enforcement structure: the copy is co-located with distinctions at $\Gamma_i$ and generates superadditivity surplus $\Delta_i > 0$ at $\Gamma_i$ (by $L_\Delta$: the copy is nonclassical, so pool-mediated coupling is present at the target).  The broadcasting map connects source to target; by admissibility, the enforcement surplus at the target is mediated by capacity committed at the source.

\textit{Step 2: Capacity cannot be created.}  By $L_{\mathrm{loc}}$ (Theorem~\ref{thm:Lloc}), capacity is additive across independent interfaces and cannot be transferred between them.  The broadcasting map is admissible (A1), so the total capacity supporting the $n$ copies is bounded by $C(\Gamma_{\mathrm{src}})$.

\textit{Step 3: Surplus accounting.}  Each faithful copy draws surplus $\Delta_i \ge \Delta_0 > 0$ from the source's budget.  By $T_M$ (generalized to $n$ targets): $\sum_{i=1}^n \Delta_i \le C(\Gamma_{\mathrm{src}})$.  Since $\Delta_i \ge \Delta_0$ for each copy: $n \cdot \Delta_0 \le C(\Gamma_{\mathrm{src}})$.

\textit{Step 4: Finite bound.}  Dividing: $n \le C(\Gamma_{\mathrm{src}})/\Delta_0$.  Since $C$ and $\Delta_0$ are both finite and positive (A1 + $L_\Delta$), $n_{\mathrm{bc}} := \lfloor C/\Delta_0 \rfloor$ is a finite positive integer.  \qed

\smallskip\noindent\textit{Corollary: universal broadcasting is impossible.}  Universal broadcasting (unlimited faithful copies, $n \to \infty$) requires unbounded capacity at the source.  At any finite-capacity interface, the number of faithful copies is bounded by $n_{\mathrm{bc}}$.

\smallskip\noindent\textit{Connection to $n_{\max}$.}  Since $\Delta_0 \ge \eps^*$, we have $n_{\mathrm{bc}} \le \lfloor C/\eps^* \rfloor = n_{\max}$: the broadcasting bound is no larger than the maximum number of independently enforceable distinctions ($T_{\mathrm{form}}$).  The same capacity constraint that makes the substrate finite-dimensional also limits the number of faithful copies.

\smallskip\noindent\textbf{Uses:} A1, $L_\Delta$ (Theorem~\ref{thm:LDelta}), $T_M$ (generalized), $L_{\mathrm{loc}}$ (Theorem~\ref{thm:Lloc}).  No algebraic structure, no field selection, no Hilbert space.

\medskip\noindent\textbf{Route 2: Qualitative no-broadcasting (conditional on $T_{\mathrm{tensor}}$).}

By the Barnum--Caves--Fuchs--Jozsa--Schumacher theorem, broadcastable states in a $C^*$-algebra lie in the center $Z(\mathcal{A}) := \{A \in \mathcal{A} : [A, B] = 0\;\forall B\}$.  By $T_{\mathrm{alg}}$, $[E_{d_1}, F_\Pi] \ne 0$, so $E_{d_1} \notin Z(\mathcal{A})$: the center is a proper subalgebra.  States associated with non-central elements are not broadcastable.  This gives the standard qualitative result: not all states of a noncommutative algebra are broadcastable.

\smallskip\noindent\textbf{Uses:} $T_{\mathrm{alg}}$ (noncommutativity), $T_{\mathrm{tensor}}$ (tensor-product identification, conditional on T2c), Barnum--Caves--Fuchs--Jozsa--Schumacher (external theorem).

\medskip\noindent\textit{Complementarity.}  Route~1 and Route~2 give complementary information.  Route~1 is quantitative (gives a specific maximum copy number $n_{\mathrm{bc}}$) but does not address state-universality for a single copy.  Route~2 is qualitative (broadcastable states $\subseteq$ center) but does not give a copy-number bound.  Route~1 is purely resource-theoretic and uses no algebraic structure; Route~2 operates within the derived $C^*$-algebra.  The quantitative bound $n_{\mathrm{bc}} = \lfloor C/\Delta_0\rfloor$ is an APF-specific prediction not available from the standard algebraic framework.
\end{theorem}

\begin{center}\small
\renewcommand{\arraystretch}{1.15}
\begin{tabular}{@{}lcccc@{}}
\toprule
\textbf{Result} & \textbf{Single algebra} & \textbf{$L_{\mathrm{loc}}$} & \textbf{$T_{\mathrm{tensor}}$} & \textbf{Full T2c} \\
\midrule
Tsirelson bound $\le 2\sqrt{2}$ & \checkmark & \checkmark & --- & --- \\
Tsirelson saturation $= 2\sqrt{2}$ & --- & \checkmark & effectively & \checkmark \\
Monogamy ($T_M$) & \checkmark & \checkmark & --- & --- \\
No-broadcast Rt.~1 (finite \#) & \checkmark & \checkmark & --- & --- \\
No-broadcast Rt.~2 (qualitative) & --- & \checkmark & \checkmark & usually \\
Born rule ($T_{\mathrm{Born}}$) & \checkmark & --- & --- & \checkmark \\
\bottomrule
\end{tabular}
\end{center}
\smallskip\noindent\textit{Table key.}  \textbf{Single algebra}: result proved within a single $\mathcal{A} \subseteq \End(V)$, no composite structure needed.  \textbf{$L_{\mathrm{loc}}$}: disjoint-support locality (commuting subalgebras for spatially separated subsystems).  \textbf{$T_{\mathrm{tensor}}$}: full tensor-product identification $\mathcal{A} = \mathcal{A}_A \otimes \mathcal{A}_B$ (conditional on T2c + compositional closure).  \textbf{Full T2c}: complex field selection (T2c.1 + T2c.2).  A checkmark (\checkmark) means the assumption is required; ``---'' means it is not; ``effectively'' and ``usually'' mean the result can be stated without the assumption but saturation/full generality requires it.  The Tsirelson bound $\le 2\sqrt{2}$ and no-broadcast Route~1 are the strongest unconditional results: they need only the single-algebra + locality tier and no field selection or tensor identification.

\begin{remark}[CBH inversion: information-theoretic constraints as theorems]\label{rem:CBH_inversion}
The Clifton--Bub--Halvorson theorem~\cite{CBH2003} derives quantum structure from three information-theoretic constraints imposed on a $C^*$-algebraic framework that is taken as given.  The logical order in this supplement is the reverse.  At the point where $T_{\mathrm{NB}}$ is stated, the full $C^*$-algebraic setting is already in place: the real $*$-algebra $\mathcal{A}$ (O3, Proposition~\ref{prop:O3}), the Wedderburn--Artin decomposition (T2a, Theorem~\ref{thm:T2a}), the GNS Hilbert space $\mathcal{H}_\omega$ ($T_{\mathrm{GNS}}$, Theorem~\ref{thm:GNS}), and the complex field selection (T2c, Corollary~\ref{cor:T2c}) have all been constructed from A1~+~MD~+~BW before any information-theoretic corollary is invoked.  The anti-circularity audit (Section~\ref{sec:skeleton}, Steps~1--8) makes this construction order explicit and verifiable.

Concretely: $T_{\mathrm{NB}}$ Route~1 uses only the resource-budget chain ($T_M$, $L_\Delta$, A1) and does not reference the $C^*$-algebra at all; Route~2 uses the algebraic result $T_{\mathrm{alg}}$ (noncommutativity) together with the Barnum--Caves--Fuchs--Jozsa--Schumacher characterization of broadcastable states, which operates \emph{within} the algebra already derived.  Neither route presupposes an information-theoretic constraint as input.  The same holds for $T_{\mathrm{NS}}$ (no-signaling), which is a corollary of $L_{\mathrm{loc}}$ (locality, Theorem~\ref{thm:Lloc}) and the sector decomposition, both of which precede the skeleton in the derivation DAG.

The structural claim is therefore: APF derives the $C^*$-algebra from enforcement costs (the content of Sections~2--5 of this supplement) and then recovers CBH's three constraints as theorems within that algebra, rather than assuming the algebra and imposing the constraints.  This is a genuine inversion of logical order, not merely a rephrasing: the algebra construction uses no information-theoretic vocabulary, and the information-theoretic results use no inputs beyond the algebra already constructed.
\end{remark}

% ═══════════════════════════════════════════════════════════════
\section{Complete Dependency Map}
\label{sec:dependency}

\begin{center}\small
\begin{tabular}{@{}lclcl@{}}
\toprule
\textbf{Result} & \textbf{Type} & \textbf{Depends on} & \textbf{Reg.\ cond.} & \textbf{Status} \\
\midrule
OR1 (convex mixing) & Thm & A1 & --- & Proved \\
$L_{\eps^*}$ (cost floor) & Thm & A1, MD & \fbox{MD} & Proved \\
$T_{\mathrm{form}}$ (finite dim) & Thm & A1, MD, SP, $L_{\mathrm{pred}}$ & \fbox{MD} & Proved \\
$L_{\mathrm{pred}}$ (direction predicates) & Lemma & SP, FD1, K2, FD3, $L_{\eps^*}$ & \fbox{MD} & Proved (constitutive FD1) \\
$T_{\mathrm{embed}}$ (vector space) & Thm & OR1, $L_{\eps^*}$ & \fbox{MD} & Proved \\
$T_{\mathrm{sep}}$ (sector decomp.) & Thm & A1, FD3/FD4 & --- & Proved \\
$L_{\mathrm{cost}}$ (cost uniqueness) & Thm & A1, $L_{\mathrm{iso}}$, $T_{\mathrm{sep}}$ & \fbox{MD${}^+$} & Cond.\ on $L_{\mathrm{iso}}$ \\
$L_\Delta$ (superadditivity) & Thm & A1, K2, K3, $T_{\mathrm{sep}}$, $\mathcal{D}$-quot. & --- & Proved \\
T1 (order-dependence) & Thm & $L_\Delta$, NT, BW & \fbox{BW} & Proved \\
$T_{\mathrm{adj}}$ (self-adjointness) & Thm & $T_{\mathrm{sector}}$, K3 & --- & Proved \\
$L_{\mathrm{blk}}$ (direct-sum failure) & Thm & $T_{\mathrm{sep}}$, $T_{\mathrm{adj}}$, $L_\Delta$ & --- & Proved \\
$L_\Pi$ (codespace rotation) & Lemma & $L_{\mathrm{blk}}$, $T_{\mathrm{adj}}$, OR3 & \fbox{MD} & Proved \\
$T_{\mathrm{alg}}$ (noncommutativity) & Thm & $L_{\mathrm{blk}}$, $T_{\mathrm{adj}}$ & --- & Proved \\
O3 (enforcement algebra) & Prop & $T_{\mathrm{adj}}$, $T_{\mathrm{sep}}$ & --- & Constructed \\
O4 (faithful state $\omega$) & Prop & $\tr$ on $\End(V)$, $\eps(d)$, $E_d$ & --- & Constructed (not assumed) \\
T2a (Wedderburn--Artin) & Thm & $T_{\mathrm{alg}}$, O3, O4, OR3 & \fbox{MD} & Proved \\
$T_{\mathrm{GNS}}$ (Hilbert space) & Thm & T2a, O4 & --- & Proved \\
T2c.1 ($\mathbb{R}$-excl.) & Thm & T2a, $T_{\mathrm{alg}}$, Lem.~\ref{lem:state_sep} & --- & Proved (single-interface) \\
T2c.2 ($\mathbb{H}$-excl.) & Thm & T2a, $L_{\mathrm{loc}}$, Brauer & --- & Conditional (composition) \\
T2c ($\mathbb{F} = \mathbb{C}$) & Cor & T2c.1, T2c.2 & --- & Proved under T2c.1+T2c.2 \\
$L_{\mathrm{loc}}$ (locality) & Thm & A1, FD3 & --- & Proved \\
$L_{\mathrm{cert}}$ (sharp cert.\ prob.) & Lemma & O4, FD5, K3 & --- & Proved \\
$L_{\mathrm{prob}}$ (effect prob.) & Lemma & $L_{\mathrm{cert}}$, O4 & --- & Proved \\
$T_{\mathrm{Born}}$ (Born rule) & Thm & T2c, $T_{\mathrm{GNS}}$, $L_{\mathrm{cert}}$, $L_{\mathrm{prob}}$ & --- & Proved \\
$T_{\mathrm{Tsirelson}}$ ($\le 2\sqrt{2}$) & Thm & T2a, $L_{\mathrm{loc}}$ & --- & Proved (bound) \\
$T_M$ (monogamy) & Thm & A1, $L_{\mathrm{loc}}$, $L_\Delta$ & --- & Proved \\
$T_{\mathrm{NB}}$ Rt.~1 (finite broadcast \#) & Thm & A1, $L_\Delta$, $T_M$, $L_{\mathrm{loc}}$ & --- & Proved \\
$T_{\mathrm{NB}}$ Rt.~2 (qualitative) & Thm & $T_{\mathrm{alg}}$, $T_{\mathrm{tensor}}$ & --- & Conditional \\
\bottomrule
\end{tabular}
\end{center}

\noindent\textit{Entry points.}  MD enters at $L_{\eps^*}$ and propagates through $T_{\mathrm{form}}$, $T_{\mathrm{embed}}$, OR3.  BW enters only at T1.  Constitutive principles (SP, K3) and formal definitions (FD1--FD6) enter at the definitions and are used throughout.  Every result not boxed follows from A1 together with these constitutive principles and definitions.  \fbox{MD${}^+$} denotes dependence on the strong (isotropy) reading of MD via $L_{\mathrm{iso}}$; under the weak (floor-only) reading, $L_{\mathrm{cost}}$'s integrality is not established.  The core chain ($L_\Delta$ through $T_{\mathrm{Born}}$) does not depend on $L_{\mathrm{cost}}$ and is robust to either reading.

\medskip\noindent\textit{Tier assignment.}  The following table classifies each result by the assumptions it uses beyond A1 and constitutive semantics.  A bullet (\textbullet) indicates the assumption is load-bearing for that result.

\smallskip
\begin{center}\small
\renewcommand{\arraystretch}{1.1}
\begin{tabular}{@{}lccccccc@{}}
\toprule
\textbf{Result} & \textbf{MD} & \textbf{BW} & \textbf{$L_{\mathrm{iso}}$} & \textbf{K3\,ex.} & \textbf{SSC} & \textbf{Comp.} & \textbf{Ext.} \\
\midrule
OR1 (convex mixing) & & & & & & & \\
$L_{\eps^*}$ (cost floor) & \textbullet & & & & & & \\
$T_{\mathrm{form}}$ (finite dim) & \textbullet & & & & & & \\
$T_{\mathrm{embed}}$ (vector space) & \textbullet & & & & & & \\
$T_{\mathrm{sep}}$ (sectors) & & & & \textbullet & & & \\
$L_{\mathrm{cost}}$ (integrality) & \textbullet & & \textbullet & \textbullet & & & \\
$L_\omega$ (bilinear form) & & & & \textbullet & & & \\
$T_{\mathrm{adj}}$ (self-adjointness) & & & & \textbullet & & & \\
$L_\Delta$ (superadditivity) & & & & \textbullet & & & \\
T1 (order-dependence) & & \textbullet & & \textbullet & & & \\
$L_{\mathrm{blk}}$, $L_\Pi$ (rotation) & \textbullet & & & \textbullet & & & \\
$T_{\mathrm{alg}}$ (noncommutativity) & & & & \textbullet & & & \\
O4 (faithful state) & & & & & & & \\
T2a (Wedderburn--Artin) & \textbullet & & & \textbullet & & & \\
$T_{\mathrm{GNS}}$ (Hilbert space) & & & & & & & \\
T2c.1 ($\mathbb{R}$-exclusion) & & & & & \textbullet & & \\
T2c.2 ($\mathbb{H}$-exclusion) & & & & & & \textbullet & \\
$L_{\mathrm{cert}}$ (sharp cert.\ prob.) & & & & \textbullet & & & \\
$L_{\mathrm{prob}}$ (effect prob.) & & & & & & & \\
$T_{\mathrm{Born}}$ (Born rule) & & & & & \textbullet & \textbullet & Busch \\
$T_{\mathrm{Tsirelson}}$ (bound) & & & & & & & Cirel'son \\
$T_M$ (monogamy) & & & & \textbullet & & & \\
$T_{\mathrm{NB}}$ Rt.~1 (finite broadcast) & & & & \textbullet & & & \\
$T_{\mathrm{NB}}$ Rt.~2 (qualitative) & & & & \textbullet & & \textbullet & Barnum+ \\
\bottomrule
\end{tabular}
\end{center}
\smallskip\noindent\textit{Column key.}  \textbf{MD}: uniform per-test cost floor (regularity).  \textbf{BW}: budget-window richness.  \textbf{$L_{\mathrm{iso}}$}: cost isotropy (strong MD).  \textbf{K3\,ex.}: exact disjoint-support additivity.  \textbf{SSC}: state-separation completeness (Def.~\ref{def:SSC}).  \textbf{Comp.}: composition closure (T2c.2).  \textbf{Ext.}: imported external theorem.  Empty cells mean the result uses only A1 + constitutive semantics (SP, FD1--FD6).  Results in the first column block (through T2a/$T_{\mathrm{GNS}}$) constitute the \emph{core chain}; T2c.1--T2c.2 are the \emph{field-selection tier}; $T_{\mathrm{Born}}$ onward is the \emph{corollary tier}.

% ╔═══════════════════════════════════════════════════════════════╗
% ║               PART II: DISCUSSION AND CONTEXT                 ║
% ║  Remarks, commentary, comparisons, scope.  No new theorems.   ║
% ╚═══════════════════════════════════════════════════════════════╝

\part{Discussion and Context}
\label{part:discussion}

The proofs in Part~I are complete and self-contained.  Part~II collects interpretive remarks, comparisons with other reconstruction programs, and the scope analysis.  No theorem in Part~I depends on anything in Part~II.

% ═══════════════════════════════════════════════════════════════
\section{Comparison with Reconstruction Programs}
\label{sec:related}

The APF derivation chain is conceptually orthogonal to existing quantum reconstruction programs but arrives at the same destination.  This section maps the APF's inputs to those of four major programs and identifies where assumptions agree, differ, or are subsumed.

\paragraph{Hardy (2001)~\cite{Hardy2001}.}
Hardy derives quantum theory from five axioms: (1)~probabilities, (2)~simplicity, (3)~subspace structure, (4)~composite systems, and (5)~continuous reversibility.  In APF terms: Hardy's Axiom~1 (state determined by probabilities) corresponds to the $\mathcal{D}$-quotient (states are identified by their enforcement-cost profiles).  Hardy's Axiom~5 (continuous reversibility) plays the role of MD: both provide a regularity condition that excludes pathological infinite-dimensional or degenerate cases.  The key structural difference is that Hardy takes local tomography as a consequence of his axiom set (via simplicity + subspace structure), while APF derives the exclusion of $\mathbb{R}$ from state-separation completeness (T2c.1, Theorem~\ref{thm:T2c1}) without invoking local tomography as a premise.

\paragraph{Chiribella--D'Ariano--Perinotti (2011)~\cite{CDP2011}.}
CDP derive quantum theory from six principles: causality, perfect distinguishability, ideal compression, local distinguishability, pure conditioning, and purification.  APF's A1 (finite enforcement capacity) plays a role analogous to CDP's causality + perfect distinguishability: both constrain the operational structure of state spaces.  CDP's purification axiom (every mixed state has a purification in a larger system) is the strongest assumption with no direct APF analog; the closest APF counterpart is MD, which is strictly weaker (a uniform cost floor rather than a structural principle about extensions).  APF does not assume purification; consequences that CDP derive from purification (e.g., no-broadcasting, monogamy) are derived in APF from budget competition ($L_\Delta$, $T_M$).

\paragraph{Masanes--M\"uller (2011)~\cite{MM2011}.}
Masanes--M\"uller derive quantum theory from five axioms, including local tomography as an explicit axiom and the existence of a continuous reversible transformation between any two pure states.  In APF, local tomography is \emph{not} an axiom: the $\mathbb{R}$-exclusion (which is the content of local tomography) is derived from state-separation completeness (T2c.1, Theorem~\ref{thm:T2c1}).  APF's field-selection mechanism is closest to Masanes--M\"uller's but arrives via the $\mathcal{D}$-quotient and the enforcement algebra's internal structure rather than by external postulation.  The continuous-reversibility axiom of Masanes--M\"uller corresponds roughly to MD.

\paragraph{Van de Wetering (2019) and effectus/sequential-product frameworks~\cite{vandeWetering2019,GudderGreechie2002}.}
Van de Wetering's effect-theoretic reconstruction derives quantum theory via sequential effect algebras and the Koecher--Vinberg theorem, arriving at EJA structure.  The field selection in that framework also excludes $\mathbb{R}$ and $\mathbb{H}$ via local tomography and compositional closure.  APF's $P_{\mathrm{cls}}$ ($\mathbb{H}$-exclusion from $L_{\mathrm{loc}}$ + Brauer group) closely parallels van de Wetering's compositional closure argument.  The $\mathbb{R}$-exclusion differs: APF uses state-separation completeness (T2c.1) rather than an axiom of local tomography.  Paper~1's Appendix~E shows that the APF's enforcement operations satisfy the filter axioms of the Gudder--Greechie sequential effect algebra framework, providing a bridge between the two approaches.

The APF's claim to reach \emph{associative} $C^*$-structure --- not merely the Jordan (EJA) structure that effectus and sequential-product approaches produce directly --- rests on two features: (i)~the enforcement algebra $\mathcal{A}$ is defined as a subalgebra of $\End(V)$ (Proposition~O3, Section~\ref{sec:skeleton}), so associativity is \emph{inherited} from function composition, not axiomatized; and (ii)~the Wedderburn--Artin decomposition (T2a) operates on this associative algebra directly.  The EJA route (Appendix~E of Paper~1) serves as a cross-check: starting from the same enforcement primitives, the Koecher--Vinberg theorem recovers $\bigoplus_k M_{n_k}(\mathbb{C})_{\mathrm{sa}}$ --- whose complexification is the same $C^*$-algebra T2a delivers.  The two routes converge, confirming that the associative structure is not an artifact of the construction but a consequence of the enforcement data.

Since the effectus and sequential-product frameworks require specific categorical/sequential properties, we provide a succinct verification map from APF's enforcement primitives to those hypotheses.  Each entry identifies the APF source theorem and the categorical property it establishes:

\smallskip
\begin{center}\small
\renewcommand{\arraystretch}{1.2}
\begin{tabular}{@{}p{3.4cm}p{3cm}p{5.5cm}@{}}
\toprule
\textbf{SEA / effectus hypothesis} & \textbf{APF source} & \textbf{How established} \\
\midrule
Effect algebra structure ($E, 0, 1, \oplus, \perp$) & A1, FD2 & Budget affordability defines $\perp$; joint enforcement defines $\oplus$; $\oplus$ commutative and associative from cost arithmetic \\
Sequential product $a \mid b$ (S1--S5) & FD5, $T_{\mathrm{adj}}$, $L_\omega$ & Conditional enforcement at marginal cost $\eps(\{d_1,d_2\})-\eps(d_2)$; S1 (additivity) from K3; S4 (associativity) inherited from $\End(V)$ \\
Sharp effects ($E_d^2 = E_d$) & FD5a & Idempotency by definition: re-enforcing an active distinction costs nothing \\
Compressions ($C_{E_d}(f) = E_d f E_d$) & $T_{\mathrm{adj}}$ & L\"uders update from orthogonal projection; retains only $M_d$-supported component \\
Filters ($F_{E_d}: \mathcal{A} \to \mathcal{A}_{\le E_d}$) & $T_{\mathrm{sep}}$, $T_{\mathrm{adj}}$ & Sector decomposition + self-adjointness give $\mathcal{A}_{\le E_d} = E_d \mathcal{A} E_d$ \\
Dagger structure ($E_d = E_d^\dagger$) & $T_{\mathrm{adj}}$, $L_\omega$ & Cost bilinear form $B$ gives $B$-adjoint; enforcement projections are $B$-self-adjoint \\
Order-unit space (GG1) & A1 & $\mathbf{1} = \sum_d E_d + E_\Pi$; $\eps(d) \le C$ for all $d$ \\
Strong positivity (GG2) & $T_{\mathrm{adj}}$, O4 & $\omega(A^*A) = 0 \Rightarrow A = 0$ (faithfulness of $\omega$) \\
Sharpness (GG3) & FD5a & $E_d \mid E_d = E_d^2 = E_d$ for all enforcement effects \\
vdW SP1 (positivity) & $T_{\mathrm{adj}}$ & $a \mid a = a^2 \ge 0$ for self-adjoint $a$ \\
vdW SP2 (unit) & O3 & $\mathbf{1} \mid a = a$ (identity enforcement) \\
vdW SP3 (continuity) & OR3 & Automatic: bilinear maps on finite-dimensional spaces \\
vdW SP4 (symmetry $\Rightarrow$ Jordan) & $T_{\mathrm{alg}}$ & $[a,b] = 0 \Rightarrow a \mid b = ab = a \circ b$ in any $*$-algebra \\
\bottomrule
\end{tabular}
\end{center}
\smallskip

\noindent Every entry in the left column is verified from the indicated APF source in Paper~1's Appendix~E (Propositions E.1, E.6, and Corollary~E.7).  The Gudder--Greechie representation theorem (GG1--GG3 $\Rightarrow$ faithful Hilbert-space representation) and van de Wetering's sequential-product theorem (SP1--SP4 $\Rightarrow$ EJA) are invoked as proved mathematical results, not as analogies.  The APF does not assume any categorical or sequential-product axiom; it derives each from enforcement cost structure and then invokes the representation theorem.

\paragraph{Clifton--Bub--Halvorson (2003)~\cite{CBH2003}.}
CBH derive the structure of quantum theory from three information-theoretic constraints---no superluminal signaling, no broadcasting, and no unconditional bit commitment---imposed on a $C^*$-algebraic framework that is taken as given.  The CBH theorem is proved \emph{within} a well-defined algebraic setting: the starting point is a kinematic $C^*$-algebra of observables, and the three constraints select which states and operations on that algebra are physically admissible.  APF inverts this logical order.  The $C^*$-algebra is not assumed; it is constructed in eight explicit steps (anti-circularity audit, Section~\ref{sec:skeleton}):

\smallskip
\begin{center}\small
\begin{tabular}{@{}ll@{}}
\textbf{APF step} & \textbf{What it produces} \\
\midrule
$T_{\mathrm{embed}}$ (Thm.~\ref{thm:Tembed}) & Finite-dimensional real vector space $V$ \\
$T_{\mathrm{sep}}$ (Thm.~\ref{thm:Tsep}), $L_\omega$ (Lem.~\ref{lem:Lomega}) & Sector decomposition + cost bilinear form $B$ \\
$T_{\mathrm{adj}}$ (Thm.~\ref{thm:Tadj_sector}) & Self-adjoint enforcement projections $E_d = E_d^*$ \\
$T_{\mathrm{alg}}$ (Thm.~\ref{thm:Talg}), O3 (Prop.~\ref{prop:O3}) & Noncommutative real $*$-algebra $\mathcal{A}$ \\
T2a (Thm.~\ref{thm:T2a}) & Wedderburn--Artin decomposition $\bigoplus_k M_{n_k}(\mathbb{F}_k)$ \\
$T_{\mathrm{GNS}}$ (Thm.~\ref{thm:GNS}) & Faithful Hilbert-space representation $\mathcal{H}_\omega$ \\
T2c (Cor.~\ref{cor:T2c}) & Field selection $\mathbb{F} = \mathbb{C}$; $C^*$-norm via $B(\mathcal{H}_\omega)$ \\
\end{tabular}
\end{center}

\smallskip\noindent Only \emph{after} this construction is complete are the information-theoretic corollaries stated.  No-signaling ($T_{\mathrm{NS}}$, a corollary of $L_{\mathrm{loc}}$) and no-broadcasting ($T_{\mathrm{NB}}$, Theorem~\ref{thm:TNB}) are proved as theorems \emph{within} the derived algebraic setting, not as axioms imposed on it (Remark~\ref{rem:CBH_inversion}).  No-unconditional bit commitment requires the dynamics of admissible maps and is the intended subject of a companion paper on dynamics.  Thus two of CBH's three constraints are theorems here, and the third is deferred to a later paper.  The structural relationship is: CBH assumes the algebra and constrains it; APF derives the algebra and recovers the constraints.

\paragraph{Connes noncommutative geometry (NCG).}
The Chamseddine--Connes spectral action program~\cite{ConnesMarcolli2008} derives the Standard Model Lagrangian from a real spectral triple $(\mathcal{A}_F, \mathcal{H}_F, D_F)$ on a finite noncommutative geometry, tensored with a continuous Riemannian manifold.  The relationship to APF is \emph{derivation vs.\ classification}: NCG classifies which spectral triples reproduce SM physics and identifies the algebra $\mathcal{A}_F = \mathbb{C} \oplus \mathbb{H} \oplus M_3(\mathbb{C})$ as the unique finite geometry satisfying the axioms of a real spectral triple with KO-dimension 6; APF derives the same algebraic data from enforcement costs.

\smallskip\noindent\textit{Compatibility with almost-commutative geometry.}  The NCG framework models spacetime as a product $M \times F$, where $M$ is a smooth 4-manifold and $F$ is a finite noncommutative space with algebra $\mathcal{A}_F$.  The full algebra $C^\infty(M) \otimes \mathcal{A}_F$ is ``almost commutative'': the continuous factor commutes, and the finite factor is noncommutative.  The APF's output is structurally compatible with this picture.  At each interface $\Gamma$, the enforcement algebra is $\mathcal{A}(\Gamma) \cong \bigoplus_k M_{n_k}(\mathbb{C})$ (T2a, Theorem~\ref{thm:T2a}); the global algebra is the inductive limit $\varinjlim \mathcal{A}(\Gamma)$ indexed over interfaces (Section~\ref{sec:scope}, paragraph~(b)).  In the continuum limit where interfaces tile a manifold $M$, this inductive limit is naturally a bundle of finite-dimensional algebras over $M$ --- which is precisely the structure of an almost-commutative geometry.  The key point is that APF's per-interface finiteness ($\dim \mathcal{A}(\Gamma) < \infty$, from OR3) is \emph{consistent with} the NCG framework's finite fiber $F$, not in tension with it: the finite-dimensional algebra at each point is the enforcement algebra, and the manifold structure emerges from the network of interfaces.

\smallskip\noindent\textit{Inner automorphisms as gauge symmetries.}  In the NCG framework, gauge symmetries arise as inner automorphisms of $\mathcal{A}_F$: the group $\mathrm{Inn}(\mathcal{A}_F) = \{u(\cdot)u^* : u \in \mathcal{U}(\mathcal{A}_F)\}$ acts on the Hilbert space $\mathcal{H}_F$ and generates the gauge fields via the ``fluctuations of the metric'' (the inner perturbations of the Dirac operator $D_F \mapsto D_F + A + JAJ^{-1}$, where $A = \sum a_i[D_F, b_i]$).  In APF, the same structure appears via the Skolem--Noether theorem applied to the derived algebra: $\mathrm{Aut}(M_n(\mathbb{C})) \cong PU(n)$ (Paper~1, T3), and the gauge group is the product of these automorphism groups over the Wedderburn blocks.  The correspondence is: APF's T3 (gauge group from algebra automorphisms) $\leftrightarrow$ NCG's inner automorphism principle.  Both identify gauge symmetry with the internal symmetry of a noncommutative algebra; the difference is that NCG takes the algebra as input and reads off the symmetry, while APF derives the algebra and then applies Skolem--Noether.

\smallskip\noindent\textit{What remains.}  The intended program of the series is to derive the specific spectral triple $(\mathcal{A}_F, \mathcal{H}_F, D_F)$ for the Standard Model from APF inputs (the algebra from the gauge group, the Hilbert space from the field content, the Dirac operator from the Yukawa structure) and to show that the spectral action principle $S = \mathrm{Tr}[f(D^2/\Lambda^2)]$ follows from A1 (spectral invariance + additivity + finiteness).  These claims are developed in companion papers and should be assessed on their own merits once the core reconstruction in this supplement has been validated.  What remains as a long-term mathematical program --- independent of any companion paper claim --- is to derive the \emph{abstract formalism} of spectral triples itself from A1's canonical object operator algebra.  This is analogous to deriving Riemannian geometry from the equivalence principle, rather than using it as a pre-existing tool; it is a mathematical research direction, not a physics gap.

\paragraph{Longo's modular-theoretic program~\cite{BratteliRobinson1987}.}
Longo's operator-algebraic approach elevates faithful normal states and GNS representations as central tools, using Tomita--Takesaki modular theory to connect entropy, energy, and Bekenstein-type bounds within an assumed $C^*$- or von Neumann algebra.  APF inverts this logical order: it starts from a resource/capacity postulate and reconstructs the algebra.  Once the algebra is in place (T2a, $T_{\mathrm{GNS}}$), the APF's faithful state $\omega$ (O4) is normal (automatic in finite dimensions) and faithful (proved from the cost structure).  The modular automorphism group $\sigma_t(A) = \rho^{it}A\rho^{-it}$ is constructible from $\omega$, and KMS features are recoverable --- this is done in the companion codebase (\texttt{L\_Tomita\_finite}).  The structural difference is: Longo assumes the algebraic setting and derives thermodynamic consequences; APF derives the algebraic setting from thermodynamic-like resource constraints and then recovers the modular structure as a corollary.  Whether the APF's enforcement cost $\eps(\cdot)$ can be identified with a modular energy in Longo's sense is a question for the Bekenstein identification in a companion paper (see Section~\ref{sec:scope}).

\paragraph{Channel-capacity superadditivity~\cite{Cirelson1980}.}
The superadditivity studied in quantum information theory (e.g., Hastings' counterexample to the minimum output entropy conjecture) concerns gains in \emph{information throughput}: joint detection of channel outputs can exceed the sum of independent detections.  APF's $L_\Delta$ is a structurally different phenomenon --- superadditivity of \emph{enforcement cost}: the joint defense of co-located distinctions exceeds the sum of individual defenses because the shared pool $\Pi$ mediates cross-sector coupling.  The formal analogy is that both arise from joint-structure constraints absent in the independent-additivity regime: channel-capacity gains come from entangled measurements across output systems, while APF's surplus comes from $\Pi$-supported perturbations that threaten both distinctions simultaneously.  The mechanisms differ, but the mathematical pattern --- strict superadditivity as a signature of non-independence --- is shared.  A systematic classification of APF's $\kappa$-driven superadditivity within the resource-theoretic taxonomy of superadditive phenomena (activation, catalysis, etc.) would be a natural direction for further work; the ``pool induces strictness'' mechanism is closest to the \emph{activation} subclass, where a resource that is individually inert becomes valuable when shared.

\paragraph{Constructive noncommutativity and effect algebras.}
Programs that construct noncommutative algebras from prescribed commutation relations (e.g., deformation quantization, Weyl algebras) take noncommutativity as input and build the algebra around it.  APF derives noncommutativity as output: the cost structure forces $[E_{d_1}, F_\Pi] \ne 0$ without prescribing any commutator.  The connection to effect algebras and POVMs is made precise in Paper~1's Appendix~E: each enforcement projection $E_d$ is a sharp effect, compressions $C_{E_d}(f) = E_d f E_d$ are L\"uders updates, filters are sub-SEA morphisms, and the dagger structure is the $B$-adjoint from $L_\omega$.  The full verification of Gudder--Greechie (GG1--GG3) and van de Wetering (SP1--SP4) hypotheses from enforcement primitives is given in this supplement's Part~II comparison with effectus frameworks (Section~\ref{sec:related}).  The bridge from enforcement operations to the standard categorical primitives of effect theory is therefore explicit, not analogical.

\paragraph{Giannelli--Chiribella (2024): Dynamical Implementability (DIN).}
Giannelli and Chiribella derive complex quantum theory from a dynamical implementability principle: every physical transformation that can be realized in principle must be realizable by a collision model (repeated interactions with ancillary systems).  DIN excludes real-Hilbert theories dynamically (certain real transformations cannot be implemented by collisions), arriving at complex Hilbert space via a Jordan-algebraic route.  The APF excludes real nonclassical blocks by a complementary, non-dynamical mechanism: state-separation completeness (T2c.1) shows that the antisymmetric subspace of a real matrix algebra is state-invisible, an algebraic defect detectable at a single interface without dynamics.  The two routes are independent and mutually reinforcing: DIN provides a dynamical exclusion, APF provides a static/algebraic one.  Quaternionic exclusion in both frameworks requires composition-level assumptions.

\paragraph{Zilly (2025): Capacity and finite distinguishability.}
Zilly derives complex amplitudes from structural postulates including a Leibniz rule (structural compatibility of distinctions), basis isotropy, and metric compatibility with a finite distinguishability bound.  The APF's MD (uniform per-test cost floor) plays a role analogous to Zilly's finite distinguishability bound (both enforce finite dimensionality from capacity constraints), and K3 (disjoint-support additivity) parallels Zilly's structural Leibniz rule (both govern how costs combine under composition).  The two approaches share the conceptual thesis that finite capacity forces complex structure, but arrive via different technical routes: Zilly through metric geometry and algebraic identities, APF through cost semantics and operator-algebraic reconstruction.  A detailed comparison of which field-selection arguments can cross-validate between the two frameworks is a natural direction for future work.

\paragraph{Algebraic probability perspective (Musat--R\o rdam 2024).}
The algebraic probability viewpoint emphasizes that the classical/quantum divide is precisely the commutativity/noncommutativity divide in finite-dimensional $C^*$-algebras: classical probability theory is the theory of commutative algebras, quantum theory is the theory of noncommutative ones.  The APF's bridge chain ($L_\Delta \to L_{\mathrm{blk}} \to T_{\mathrm{alg}}$) provides a concrete operational mechanism for crossing this divide: superadditivity of enforcement costs, driven by the pool sector $\Pi$, forces noncommuting generators in the enforcement algebra.  In the algebraic-probabilistic language, the APF constructs the ``observable algebra'' by exhibiting specific enforcement operations that generate non-abelian matrix blocks.  The pool $\Pi$ is the operational ingredient that prevents the algebra from collapsing to a commutative (classical) diagonal subalgebra.

\paragraph{Information-theoretic action (Carcassi--Aidala 2024).}
The information-theoretic action approach selects complex amplitudes by requiring coherence under finite resolution and normalization constraints.  The APF arrives at complex structure from a different direction (cost/capacity constraints rather than action/coherence), but the two approaches share a structural feature: both extract the complex field from the requirement that a finite-resource constraint be self-consistently satisfiable.  In the action approach, finite resolution forces phase structure; in the APF, finite capacity forces noncommutativity and then state-separation completeness selects $\mathbb{C}$.  Whether MD/BW play roles analogous to finite-action resolution and normalization constraints is an open question that could illuminate common structural necessities across reconstruction programs.

\paragraph{Summary of assumption mapping.}
The APF's three inputs map to the reconstruction literature as follows.  A1 (finite enforcement capacity) replaces the combination of causality, perfect distinguishability, and finite state-space dimension.  MD (uniform per-test cost floor) replaces continuous reversibility (Hardy) or purification (CDP); it is strictly weaker than both.  BW (budget-window richness) has no direct analog in the other programs; it is a mild genericity condition on the cost spectrum.  The structural novelty of APF is that K3 (additivity on disjoint supports, derived from A1's additive budget semantics) and the $\mathcal{D}$-quotient (constitutive of the framework's definition of physical state) together do the work that local tomography, purification, and spectrality do in the other programs.  The price paid for fewer axioms is a longer derivation chain; the benefit is that assumptions like local tomography and compositional closure become theorems rather than inputs.

% ═══════════════════════════════════════════════════════════════
\section{Scope and Limitations}
\label{sec:scope}

\paragraph{Finite dimensionality: explicit bounds, scope, and reconciliation with infinite-dimensional QM.}
\textit{The present framework is a finite-dimensional reconstruction: every result in this supplement holds at a single finite-capacity interface.  Infinite-dimensional quantum systems (QFT, continuous variables) are recovered as inductive limits of finite-capacity interfaces, with the continuum theory emerging in the $C \to \infty$ limit.  The framework is therefore best understood as a finite-dimensional effective theory, with infinite dimensions as a limiting idealization.}

\smallskip\noindent $T_{\mathrm{form}}$ (Theorem~\ref{thm:Tform}) provides three explicit upper bounds from A1 + MD:
\[
\dim(S_\Gamma) \le \frac{C(\Gamma)}{\mu^*}, \qquad
n_{\max} := \left\lfloor \frac{C(\Gamma)}{\eps^*} \right\rfloor, \qquad
\dim(\mathcal{A}) \le (n_{\max} + 1)^2, \qquad
\dim(\mathcal{H}_\omega) \le n_{\max}^2.
\]
All four bounds are \emph{per-interface} (local), not cosmological.  Each interface $\Gamma$ has its own capacity $C(\Gamma)$, its own substrate $S_\Gamma$, and its own algebra $\mathcal{A}(\Gamma)$.  No result in this supplement asserts a global bound on the dimension of ``the'' Hilbert space of the universe.

\smallskip\noindent\textit{Reconciliation with infinite-dimensional quantum systems.}  Standard quantum mechanics routinely uses infinite-dimensional Hilbert spaces (harmonic oscillator, free fields, continuous variables).  The APF framework is compatible with this in two ways.

\noindent(a)~\textit{The $C \to \infty$ limit.}  As $C(\Gamma) \to \infty$ with $\mu^*$ fixed, $n_{\max} \to \infty$ and $\dim(\mathcal{H}_\omega) \to \infty$.  The finite-dimensional algebra $\bigoplus_k M_{n_k}(\mathbb{C})$ converges (in the appropriate inductive-limit topology) to a $C^*$-algebra of compact operators or a Type~I von Neumann factor.  The framework does not forbid infinite dimensions; it says finite capacity enforces finite dimensions, and the infinite-dimensional case arises as a limiting regime.  This limit is not developed in the present supplement; it is the intended subject of a companion paper on dynamics.

\noindent(b)~\textit{Per-interface finiteness does not require global finiteness.}  A quantum field theory can be modeled as a net of local algebras $\{\mathcal{A}(\Gamma)\}_\Gamma$, each finite-dimensional (with dimension bounded by $C(\Gamma)/\mu^*$), indexed over a directed set of interfaces.  The global algebra is the inductive limit $\varinjlim \mathcal{A}(\Gamma)$, which is infinite-dimensional even though each local algebra is finite-dimensional.  This is the standard algebraic QFT construction (Haag--Kastler nets).  APF's finite-dimensionality claim is about each local algebra, not about the inductive limit.

\noindent(c)~\textit{Bekenstein--Hawking as physical motivation.}  The Bekenstein--Hawking bound suggests $n_{\max} \le A / 4\ell_P^2$ for an interface of area $A$, which would provide a physical value of $\mu^*$.  If this identification holds (it is developed in a companion paper and is not required by any result here), then for a Planck-area interface ($A \sim \ell_P^2$), $n_{\max} \sim 1$: a single qubit; for a macroscopic interface ($A \sim 1\,\mathrm{m}^2$), $n_{\max} \sim 10^{69}$: the finite-dimensional approximation is excellent for all practical purposes.  The ``infinite-dimensional'' regime would be the idealization $A \to \infty$.

\paragraph{Proved vs.\ conditional.}  Structural results (noncommutativity, $C^*$-algebra, field selection, Born rule, Tsirelson bound, monogamy) are proved from A1 + MD + BW together with the constitutive principles (SP, K3) and formal definitions (FD1--FD6).  Real-block exclusion (T2c.1) is proved from state-separation completeness (single-interface); quaternionic-block exclusion (T2c.2) additionally requires composition closure.  Results requiring tensor-product identification ($T_{\mathrm{tensor}}$, algebraic $T_{\mathrm{NB}}$) inherit the composition condition.  $T_{\mathrm{CPTP}}$ (complete positivity of dynamics) is not proved here; it requires a dynamical framework beyond the scope of this structural derivation.

\paragraph{Where MD is load-bearing.}  Without MD: $\eps^* = 0$, $\dim(S_\Gamma) = \infty$, Wedderburn--Artin fails.  The bridge ($L_\Delta \to L_{\mathrm{blk}} \to T_{\mathrm{alg}}$) survives locally but the global skeleton is lost.

\paragraph{Robustness under approximate K3 and $\kappa$-dependence.}  Exact K3 (disjoint-support additivity) is a constitutive premise of the budget semantics.  Proposition~\ref{prop:robustness} quantifies robustness: if cross-support couplings of size $\delta$ are present, continuous quantities (bilinear form, projections, probabilities) degrade as $O(\delta)$, while discrete invariants (Wedderburn block structure, field selection $\mathbb{C}$) are exactly preserved for $\delta < \delta_{\mathrm{crit}}$.  The derivation chain is structurally invariant across the admissible $\kappa$-class in the following precise sense: any cost function satisfying K1--K3 produces the same inter-sector orthogonality, the same projection structure, and the same downstream algebra (Proposition~\ref{prop:kappa_class}).  Different $\kappa$ choices affect only the within-sector normalization of $B$, which is a free gauge.

\paragraph{Where BW is load-bearing.}  Without BW: T1's witness triple may not exist.  BW holds for any cost spectrum with $\ge 3$ distinct values.

\paragraph{Experimental signatures: what this derivation predicts and what it does not.}
The central results of this supplement --- Hilbert space, Born rule, $C^*$-algebra, Tsirelson bound, no-broadcasting --- \emph{reproduce} the structure of standard quantum mechanics.  They are not departures from it.  The framework's claim is that these structures are \emph{forced} by finite enforcement capacity together with regularity (MD, BW) and constitutive semantics, with complex field selection requiring additional composition closure.  Consequently, no experiment testing standard quantum mechanics at ordinary scales can discriminate APF from any other quantum reconstruction program (Hardy, CDP, Masanes--M\"uller, etc.): all arrive at the same destination by different routes, and the destination \emph{is} standard quantum theory.

The discriminating experimental signatures lie elsewhere --- specifically, in regimes where finite capacity becomes operationally visible.  The main paper identifies a near-term, platform-specific test: the \emph{decoherence threshold} (``kink'') in qubit dephasing as a function of environmental spectral weight $S(\omega_0)$.  The APF predicts a crossover at $S(\omega_0) = S_{\mathrm{th}} := \eps^*/C$ from a flat ($T_2$-independent) regime below threshold to the standard Bloch--Redfield ($T_2^{-1} \propto S$) regime above threshold.  The crossover form $T_2^{-1} \propto \max(S/S_{\mathrm{th}} - 1,\, 0)^{\mu(s)}$ (with $\mu = s$ the spectral index) is quantitatively distinct from the null hypothesis (no kink, slope $-1$ throughout) and is testable on existing platforms: superconducting transmon qubits with a single near-resonant TLS (Schl\"or et al.\ 2019), or trapped ions with engineered phonon reservoirs (Myatt et al.\ 2000).  The main paper gives explicit protocols (P1--P5 for transmons, Q1--Q4 for trapped ions), discriminating observables ($d\ln T_2 / d\ln S$ vs.\ $S/S_{\mathrm{th}}$), and mode-count scaling predictions ($S_{\mathrm{th}} \to S_{\mathrm{th}}/N_E$ as the number of environmental modes increases).

This supplement does not develop the experimental analysis (which requires dynamics beyond the scope of this structural derivation); it flags the logical status.  The decoherence threshold depends on the capacity-discreteness structure ($\eps^* > 0$, $n_{\max}$ finite) proved in this supplement, but the crossover form additionally requires a dynamical framework ($T_{\mathrm{CPTP}}$: complete positivity of admissible maps) and the dimensional identification $\eps^* = \alpha\ell_P^2$ (Bekenstein--Hawking connection), both of which are the intended subjects of companion papers and are not established here.  The supplement's structural results (Hilbert space, Born rule, Tsirelson bound) are necessary prerequisites --- they establish the arena within which the threshold prediction lives --- but they are not themselves the discriminating observables.

\paragraph{Bekenstein--Hawking identification: scope disclaimer.}
Several results in this supplement reference the Bekenstein--Hawking bound $S \le A/4\ell_P^2$ as physical motivation for MD and as an independent confirmation of OR3.  The precise identification $\eps^* = \alpha\ell_P^2$ (with $\alpha$ a numerical coefficient fixed by the area--entropy relation) is the intended subject of a companion paper; no result in this supplement depends on it.  That companion paper must engage with well-known subtleties: Page's counterexamples to the original Bekenstein bound (which do not apply to the covariant form used here), Casini's modular-energy reformulation (which clarifies the correct energy concept), Hayden--Wang's task-dependence analysis (which constrains the operational interpretation of $\eps(\cdot)$), and the reconciliation of finite-dimensional enforcement accounting with continuum renormalization.  These are the explicit opening tasks of Paper~6.  No result proved in this supplement is affected by their resolution.

% ═══════════════════════════════════════════════════════════════
\begin{thebibliography}{99}
\bibitem{Aczel1966} J.~Acz\'el, \textit{Lectures on Functional Equations and their Applications}, Academic Press, 1966.
\bibitem{Artin1927} E.~Artin, ``Zur Theorie der hyperkomplexen Zahlen,'' \textit{Abh.\ Math.\ Sem.\ Hamburg} \textbf{5} (1927) 251--260.
\bibitem{BarnumWilce2012} H.~Barnum and A.~Wilce, ``Post-classical probability theory,'' in \textit{Quantum Theory: Informational Foundations and Foils}, Springer, 2016.
\bibitem{Bekenstein1981} J.~D.~Bekenstein, ``Universal upper bound on the entropy-to-energy ratio for bounded systems,'' \textit{Phys.\ Rev.\ D} \textbf{23} (1981) 287.
\bibitem{Bousso2002} R.~Bousso, ``The holographic principle,'' \textit{Rev.\ Mod.\ Phys.} \textbf{74} (2002) 825--874.
\bibitem{BratteliRobinson1987} O.~Bratteli and D.~W.~Robinson, \textit{Operator Algebras and Quantum Statistical Mechanics}, 2nd ed., Springer, 1987.
\bibitem{Busch2003} P.~Busch, ``Quantum states and generalized observables: a simple proof of Gleason's theorem,'' \textit{Phys.\ Rev.\ Lett.} \textbf{91} (2003) 120403.
\bibitem{Casini2008} H.~Casini, ``Relative entropy and the Bekenstein bound,'' \textit{Class.\ Quantum Grav.} \textbf{25} (2008) 205021.
\bibitem{CBH2003} R.~Clifton, J.~Bub, and H.~Halvorson, ``Characterizing quantum theory in terms of information-theoretic constraints,'' \textit{Found.\ Phys.} \textbf{33} (2003) 1561.
\bibitem{CDP2011} G.~Chiribella, G.~M.~D'Ariano, and P.~Perinotti, ``Informational derivation of quantum theory,'' \textit{Phys.\ Rev.\ A} \textbf{84} (2011) 012311.
\bibitem{Cirelson1980} B.~S.~Cirel'son, ``Quantum generalizations of Bell's inequality,'' \textit{Lett.\ Math.\ Phys.} \textbf{4} (1980) 93--100.
\bibitem{ConnesMarcolli2008} A.~Connes and M.~Marcolli, \textit{Noncommutative Geometry, Quantum Fields and Motives}, AMS, 2008.
\bibitem{Frobenius1878} F.~G.~Frobenius, ``\"Uber lineare Substitutionen und bilineare Formen,'' \textit{J.\ reine angew.\ Math.} \textbf{84} (1878) 1--63.
\bibitem{Gleason1957} A.~M.~Gleason, ``Measures on the closed subspaces of a Hilbert space,'' \textit{J.\ Math.\ Mech.} \textbf{6} (1957) 885--893.
\bibitem{GNS1943} I.~M.~Gel'fand and M.~A.~Naimark, ``On the imbedding of normed rings into the ring of operators in Hilbert space,'' \textit{Mat.\ Sb.} \textbf{12} (1943) 197--213.
\bibitem{GudderGreechie2002} S.~Gudder and R.~Greechie, ``Sequential products on effect algebras,'' \textit{Rep.\ Math.\ Phys.} \textbf{49} (2002) 87--111.
\bibitem{Hardy2001} L.~Hardy, ``Quantum theory from five reasonable axioms,'' arXiv:quant-ph/0101012, 2001.
\bibitem{HaydenWang2024} P.~Hayden and G.~Wang, ``Operational interpretations of entropy bounds,'' arXiv:2401.XXXXX, 2024.
\bibitem{Hensen2015} B.~Hensen et al., ``Loophole-free Bell inequality violation using electron spins separated by 1.3 kilometres,'' \textit{Nature} \textbf{526} (2015) 682.
\bibitem{MM2011} L.~Masanes and M.~P.~M\"uller, ``A derivation of quantum theory from physical requirements,'' \textit{New J.\ Phys.} \textbf{13} (2011) 063001.
\bibitem{Page2005} D.~N.~Page, ``Comment on `Entropy bounds in unitary quantum evolution','' \textit{J.\ Phys.\ A} \textbf{38} (2005) 5759--5764.
\bibitem{vandeWetering2019} J.~van de Wetering, ``An effect-theoretic reconstruction of quantum theory,'' \textit{Compositionality} \textbf{1} (2019) 1.
\bibitem{Wedderburn1908} J.~H.~M.~Wedderburn, ``On hypercomplex numbers,'' \textit{Proc.\ London Math.\ Soc.} \textbf{6} (1908) 77--118.
\bibitem{Parzygnat2018} A.~J.~Parzygnat, ``From observables and states to Hilbert space and back: a 2-categorical adjunction,'' \textit{Appl.\ Categ.\ Struct.} \textbf{26} (2018) 1123--1157.  arXiv:1609.08975.
\bibitem{JonesPenneys2017} C.~Jones and D.~Penneys, ``Operator algebras in rigid $C^*$-tensor categories,'' \textit{Comm.\ Math.\ Phys.} \textbf{355} (2017) 1121--1188.  arXiv:1611.04620.
\end{thebibliography}

\end{document}
